% sec:spec-interpolation — Polynomial Interpolation
\section{Polynomial Interpolation}
\label{sec:spec-interpolation}

The discontinuous Galerkin solver at the heart of the numerical
relativity engine represents the solution within each element as a
polynomial.  Every operation the solver performs --- differentiating,
integrating, exchanging data across element faces --- reduces to
manipulating polynomials defined by their values at a set of nodes.
The quality of this representation depends entirely on where those
nodes are placed.  This section develops the interpolation theory that
governs the choice: Lagrange interpolation on general node sets, the
error bound that ties accuracy to node placement, and the spectacular
failure that results from the most natural choice of all ---
equidistant points.

\paragraph{Lagrange interpolation.}

Suppose we are given $N + 1$ distinct
nodes on an interval $[a, b]$ and the corresponding function values.
There exists a unique polynomial $p_N$ of degree at most $N$ that
passes through every data point $(x_j, f(x_j))$.  The idea is to construct $N + 1$ ``selector'' polynomials,
each of which equals one at exactly one node and vanishes at all the
others; the interpolant is then a weighted combination that picks out
the correct function value at every node.  The explicit construction
uses the \emph{Lagrange basis polynomials}
%
\begin{equation}
  \ell_j(x)
    = \prod_{\substack{k=0 \\ k \neq j}}^{N}
      \frac{x - x_k}{x_j - x_k} \,,
  \qquad j = 0, 1, \ldots, N \,.
  \label{eq:lagrange-basis}
\end{equation}
%
Each basis polynomial $\ell_j$ equals one at $x_j$ and vanishes at
every other node: $\ell_j(x_k) = \delta_{jk}$.  The interpolant is
then
%
\begin{equation}
  p_N(x)
    = \sum_{j=0}^{N} f(x_j)\,\ell_j(x) \,.
  \label{eq:lagrange-interp}
\end{equation}
%
This is not an approximation --- $p_N$ agrees with $f$ exactly at the
nodes.  The question is how well it approximates $f$ between them.

\paragraph{The interpolation error.}

For a function $f \in C^{N+1}[a,b]$, the pointwise error satisfies
%
\begin{equation}
  f(x) - p_N(x)
    = \frac{f^{(N+1)}(\xi)}{(N+1)!}\;\omega_{N+1}(x) \,,
  \label{eq:interp-error}
\end{equation}
%
where $\xi \in [a,b]$ depends on $x$ and $\omega_{N+1}$ is the
\emph{node polynomial}
%
\begin{equation}
  \omega_{N+1}(x)
    = \prod_{j=0}^{N} (x - x_j) \,.
  \label{eq:node-polynomial}
\end{equation}
%
\Cref{eq:interp-error} separates the error into two factors: the
smoothness of $f$ (through its $(N+1)$-th derivative) and the
geometry of the node set (through $\omega_{N+1}$).  The first factor
is fixed by the function being interpolated; the second is under our
control.  Choosing nodes to minimise $\max_{x \in [a,b]}
\abs{\omega_{N+1}(x)}$ is the central problem of interpolation
theory, and the subject of the next section.
\physnote{The error formula \cref{eq:interp-error} has the same
structure as the remainder term in Taylor's theorem --- both express
the approximation error through a higher derivative evaluated at an
unknown intermediate point.  The difference is that Taylor's theorem
uses a single expansion point, while Lagrange interpolation
distributes the information across $N + 1$ nodes.}

\paragraph{The Lebesgue constant.}

The error formula \cref{eq:interp-error} requires $(N+1)$-th
derivative information that is rarely available in practice.  A more
operational measure of interpolation quality is the \emph{Lebesgue
constant} $\Lambda_N$, defined as
%
\begin{equation}
  \Lambda_N
    = \max_{x \in [a,b]}
      \sum_{j=0}^{N} \abs{\ell_j(x)} \,.
  \label{eq:lebesgue-constant}
\end{equation}
%
Informally, $\Lambda_N$ measures the worst-case amplification: if the
function values $f(x_j)$ are perturbed by a small amount $\epsilon$,
the interpolant can change by as much as $\Lambda_N\,\epsilon$.  A
Lebesgue constant near one means the interpolant is stable; a
Lebesgue constant of $10^4$ means tiny perturbations at the nodes
produce enormous swings between them.  The Lebesgue constant also
controls how close the interpolant is to the best possible polynomial
approximation: if $p_N^*$ denotes the best approximation of degree $N$
in the supremum norm, then
%
\begin{equation}
  \norm{f - p_N}_\infty
    \;\leq\; (1 + \Lambda_N)\,\norm{f - p_N^*}_\infty \,.
  \label{eq:lebesgue-bound}
\end{equation}
%
A small Lebesgue constant means the interpolant is near-optimal;
a large one means the node set amplifies errors unnecessarily.  For
equidistant nodes, $\Lambda_N$ grows exponentially with $N$ --- a fact
that underlies the Runge phenomenon described below.
\warnnote{The Lebesgue constant is a property of the node set alone,
not of any particular function.  A large $\Lambda_N$ means the node
set is a poor choice for \emph{every} smooth function, not just
pathological ones.}

\paragraph{The Runge phenomenon.}

The most natural choice of nodes --- equally spaced on $[a, b]$ ---
turns out to be catastrophically bad at high polynomial degree.  Carl
Runge demonstrated this in 1901 using the innocent-looking function
%
\begin{equation}
  f(x) = \frac{1}{1 + 25\,x^2}
  \label{eq:runge-function}
\end{equation}
%
on $[-1, 1]$.  This function is smooth, bounded, and gently peaked at
the origin.  Yet polynomial interpolation on equidistant nodes
diverges as $N$ increases: the interpolant develops wild oscillations
near the endpoints that grow without bound.  At degree 15 the
endpoint error already exceeds the function's maximum value, and by
degree 25 the oscillations dwarf the function itself --- adding more
data makes the approximation worse, not better.
\histnote{Runge published this counterexample in 1901, nearly half a
century after Chebyshev had identified the optimal node distribution
in the 1850s.  The phenomenon surprised many who assumed that more
data points would always improve polynomial approximation.}

The node polynomial of \cref{eq:node-polynomial} reveals the
mechanism.  For equidistant nodes, $\omega_{N+1}$
is much larger near the endpoints than near the centre of the
interval.  The ratio between its endpoint and midpoint values grows
exponentially with $N$, so the error formula \cref{eq:interp-error}
produces large errors near $x = \pm 1$ even when the derivative
factor $f^{(N+1)}$ remains moderate.  The Lebesgue constant for
equidistant nodes grows as
$\Lambda_N \sim 2^N / (e\,N\ln N)$ \cite{Trefethen2000}, confirming
that the instability is not specific to Runge's function but is an
intrinsic property of the node distribution.

\begin{figure}[t]
\centering
%% Creator: Matplotlib, PGF backend
%%
%% To include the figure in your LaTeX document, write
%%   \input{<filename>.pgf}
%%
%% Make sure the required packages are loaded in your preamble
%%   \usepackage{pgf}
%%
%% Also ensure that all the required font packages are loaded; for instance,
%% the lmodern package is sometimes necessary when using math font.
%%   \usepackage{lmodern}
%%
%% Figures using additional raster images can only be included by \input if
%% they are in the same directory as the main LaTeX file. For loading figures
%% from other directories you can use the `import` package
%%   \usepackage{import}
%%
%% and then include the figures with
%%   \import{<path to file>}{<filename>.pgf}
%%
%% Matplotlib used the following preamble
%%   \def\mathdefault#1{#1}
%%   \everymath=\expandafter{\the\everymath\displaystyle}
%%   \IfFileExists{scrextend.sty}{
%%     \usepackage[fontsize=10.000000pt]{scrextend}
%%   }{
%%     \renewcommand{\normalsize}{\fontsize{10.000000}{12.000000}\selectfont}
%%     \normalsize
%%   }
%%   \usepackage{fontspec}
%%   \usepackage{unicode-math}
%%   \setmainfont{STIX Two Text}[Numbers=OldStyle, Ligatures=TeX]
%%   \setmathfont{STIX Two Math}
%%   \setmonofont{Source Code Pro}[Scale=MatchLowercase]
%%   \ifdefined\pdftexversion\else  % non-pdftex case.
%%     \usepackage{fontspec}
%%     \setmainfont{DejaVuSerif.ttf}[Path=\detokenize{/Users/gvn/.pyenv/versions/3.12.11/lib/python3.12/site-packages/matplotlib/mpl-data/fonts/ttf/}]
%%     \setsansfont{DejaVuSans.ttf}[Path=\detokenize{/Users/gvn/.pyenv/versions/3.12.11/lib/python3.12/site-packages/matplotlib/mpl-data/fonts/ttf/}]
%%     \setmonofont{DejaVuSansMono.ttf}[Path=\detokenize{/Users/gvn/.pyenv/versions/3.12.11/lib/python3.12/site-packages/matplotlib/mpl-data/fonts/ttf/}]
%%   \fi
%%   \makeatletter\@ifpackageloaded{underscore}{}{\usepackage[strings]{underscore}}\makeatother
%%
\begingroup%
\makeatletter%
\begin{pgfpicture}%
\pgfpathrectangle{\pgfpointorigin}{\pgfqpoint{3.788821in}{2.311421in}}%
\pgfusepath{use as bounding box, clip}%
\begin{pgfscope}%
\pgfsetbuttcap%
\pgfsetmiterjoin%
\definecolor{currentfill}{rgb}{1.000000,1.000000,1.000000}%
\pgfsetfillcolor{currentfill}%
\pgfsetlinewidth{0.000000pt}%
\definecolor{currentstroke}{rgb}{1.000000,1.000000,1.000000}%
\pgfsetstrokecolor{currentstroke}%
\pgfsetdash{}{0pt}%
\pgfpathmoveto{\pgfqpoint{0.000000in}{0.000000in}}%
\pgfpathlineto{\pgfqpoint{3.788821in}{0.000000in}}%
\pgfpathlineto{\pgfqpoint{3.788821in}{2.311421in}}%
\pgfpathlineto{\pgfqpoint{0.000000in}{2.311421in}}%
\pgfpathlineto{\pgfqpoint{0.000000in}{0.000000in}}%
\pgfpathclose%
\pgfusepath{fill}%
\end{pgfscope}%
\begin{pgfscope}%
\pgfsetbuttcap%
\pgfsetmiterjoin%
\definecolor{currentfill}{rgb}{1.000000,1.000000,1.000000}%
\pgfsetfillcolor{currentfill}%
\pgfsetlinewidth{0.000000pt}%
\definecolor{currentstroke}{rgb}{0.000000,0.000000,0.000000}%
\pgfsetstrokecolor{currentstroke}%
\pgfsetstrokeopacity{0.000000}%
\pgfsetdash{}{0pt}%
\pgfpathmoveto{\pgfqpoint{0.507620in}{0.352669in}}%
\pgfpathlineto{\pgfqpoint{1.857769in}{0.352669in}}%
\pgfpathlineto{\pgfqpoint{1.857769in}{2.123669in}}%
\pgfpathlineto{\pgfqpoint{0.507620in}{2.123669in}}%
\pgfpathlineto{\pgfqpoint{0.507620in}{0.352669in}}%
\pgfpathclose%
\pgfusepath{fill}%
\end{pgfscope}%
\begin{pgfscope}%
\pgfsetbuttcap%
\pgfsetroundjoin%
\definecolor{currentfill}{rgb}{0.000000,0.000000,0.000000}%
\pgfsetfillcolor{currentfill}%
\pgfsetlinewidth{0.501875pt}%
\definecolor{currentstroke}{rgb}{0.000000,0.000000,0.000000}%
\pgfsetstrokecolor{currentstroke}%
\pgfsetdash{}{0pt}%
\pgfsys@defobject{currentmarker}{\pgfqpoint{0.000000in}{0.000000in}}{\pgfqpoint{0.000000in}{0.041667in}}{%
\pgfpathmoveto{\pgfqpoint{0.000000in}{0.000000in}}%
\pgfpathlineto{\pgfqpoint{0.000000in}{0.041667in}}%
\pgfusepath{stroke,fill}%
}%
\begin{pgfscope}%
\pgfsys@transformshift{0.507620in}{0.352669in}%
\pgfsys@useobject{currentmarker}{}%
\end{pgfscope}%
\end{pgfscope}%
\begin{pgfscope}%
\definecolor{textcolor}{rgb}{0.000000,0.000000,0.000000}%
\pgfsetstrokecolor{textcolor}%
\pgfsetfillcolor{textcolor}%
\pgftext[x=0.507620in,y=0.304058in,,top]{\color{textcolor}{\sffamily\fontsize{8.000000}{9.600000}\selectfont\catcode`\^=\active\def^{\ifmmode\sp\else\^{}\fi}\catcode`\%=\active\def%{\%}\ensuremath{-}1.0}}%
\end{pgfscope}%
\begin{pgfscope}%
\pgfsetbuttcap%
\pgfsetroundjoin%
\definecolor{currentfill}{rgb}{0.000000,0.000000,0.000000}%
\pgfsetfillcolor{currentfill}%
\pgfsetlinewidth{0.501875pt}%
\definecolor{currentstroke}{rgb}{0.000000,0.000000,0.000000}%
\pgfsetstrokecolor{currentstroke}%
\pgfsetdash{}{0pt}%
\pgfsys@defobject{currentmarker}{\pgfqpoint{0.000000in}{0.000000in}}{\pgfqpoint{0.000000in}{0.041667in}}{%
\pgfpathmoveto{\pgfqpoint{0.000000in}{0.000000in}}%
\pgfpathlineto{\pgfqpoint{0.000000in}{0.041667in}}%
\pgfusepath{stroke,fill}%
}%
\begin{pgfscope}%
\pgfsys@transformshift{0.845157in}{0.352669in}%
\pgfsys@useobject{currentmarker}{}%
\end{pgfscope}%
\end{pgfscope}%
\begin{pgfscope}%
\definecolor{textcolor}{rgb}{0.000000,0.000000,0.000000}%
\pgfsetstrokecolor{textcolor}%
\pgfsetfillcolor{textcolor}%
\pgftext[x=0.845157in,y=0.304058in,,top]{\color{textcolor}{\sffamily\fontsize{8.000000}{9.600000}\selectfont\catcode`\^=\active\def^{\ifmmode\sp\else\^{}\fi}\catcode`\%=\active\def%{\%}\ensuremath{-}0.5}}%
\end{pgfscope}%
\begin{pgfscope}%
\pgfsetbuttcap%
\pgfsetroundjoin%
\definecolor{currentfill}{rgb}{0.000000,0.000000,0.000000}%
\pgfsetfillcolor{currentfill}%
\pgfsetlinewidth{0.501875pt}%
\definecolor{currentstroke}{rgb}{0.000000,0.000000,0.000000}%
\pgfsetstrokecolor{currentstroke}%
\pgfsetdash{}{0pt}%
\pgfsys@defobject{currentmarker}{\pgfqpoint{0.000000in}{0.000000in}}{\pgfqpoint{0.000000in}{0.041667in}}{%
\pgfpathmoveto{\pgfqpoint{0.000000in}{0.000000in}}%
\pgfpathlineto{\pgfqpoint{0.000000in}{0.041667in}}%
\pgfusepath{stroke,fill}%
}%
\begin{pgfscope}%
\pgfsys@transformshift{1.182694in}{0.352669in}%
\pgfsys@useobject{currentmarker}{}%
\end{pgfscope}%
\end{pgfscope}%
\begin{pgfscope}%
\definecolor{textcolor}{rgb}{0.000000,0.000000,0.000000}%
\pgfsetstrokecolor{textcolor}%
\pgfsetfillcolor{textcolor}%
\pgftext[x=1.182694in,y=0.304058in,,top]{\color{textcolor}{\sffamily\fontsize{8.000000}{9.600000}\selectfont\catcode`\^=\active\def^{\ifmmode\sp\else\^{}\fi}\catcode`\%=\active\def%{\%}0.0}}%
\end{pgfscope}%
\begin{pgfscope}%
\pgfsetbuttcap%
\pgfsetroundjoin%
\definecolor{currentfill}{rgb}{0.000000,0.000000,0.000000}%
\pgfsetfillcolor{currentfill}%
\pgfsetlinewidth{0.501875pt}%
\definecolor{currentstroke}{rgb}{0.000000,0.000000,0.000000}%
\pgfsetstrokecolor{currentstroke}%
\pgfsetdash{}{0pt}%
\pgfsys@defobject{currentmarker}{\pgfqpoint{0.000000in}{0.000000in}}{\pgfqpoint{0.000000in}{0.041667in}}{%
\pgfpathmoveto{\pgfqpoint{0.000000in}{0.000000in}}%
\pgfpathlineto{\pgfqpoint{0.000000in}{0.041667in}}%
\pgfusepath{stroke,fill}%
}%
\begin{pgfscope}%
\pgfsys@transformshift{1.520231in}{0.352669in}%
\pgfsys@useobject{currentmarker}{}%
\end{pgfscope}%
\end{pgfscope}%
\begin{pgfscope}%
\definecolor{textcolor}{rgb}{0.000000,0.000000,0.000000}%
\pgfsetstrokecolor{textcolor}%
\pgfsetfillcolor{textcolor}%
\pgftext[x=1.520231in,y=0.304058in,,top]{\color{textcolor}{\sffamily\fontsize{8.000000}{9.600000}\selectfont\catcode`\^=\active\def^{\ifmmode\sp\else\^{}\fi}\catcode`\%=\active\def%{\%}0.5}}%
\end{pgfscope}%
\begin{pgfscope}%
\pgfsetbuttcap%
\pgfsetroundjoin%
\definecolor{currentfill}{rgb}{0.000000,0.000000,0.000000}%
\pgfsetfillcolor{currentfill}%
\pgfsetlinewidth{0.501875pt}%
\definecolor{currentstroke}{rgb}{0.000000,0.000000,0.000000}%
\pgfsetstrokecolor{currentstroke}%
\pgfsetdash{}{0pt}%
\pgfsys@defobject{currentmarker}{\pgfqpoint{0.000000in}{0.000000in}}{\pgfqpoint{0.000000in}{0.041667in}}{%
\pgfpathmoveto{\pgfqpoint{0.000000in}{0.000000in}}%
\pgfpathlineto{\pgfqpoint{0.000000in}{0.041667in}}%
\pgfusepath{stroke,fill}%
}%
\begin{pgfscope}%
\pgfsys@transformshift{1.857769in}{0.352669in}%
\pgfsys@useobject{currentmarker}{}%
\end{pgfscope}%
\end{pgfscope}%
\begin{pgfscope}%
\definecolor{textcolor}{rgb}{0.000000,0.000000,0.000000}%
\pgfsetstrokecolor{textcolor}%
\pgfsetfillcolor{textcolor}%
\pgftext[x=1.857769in,y=0.304058in,,top]{\color{textcolor}{\sffamily\fontsize{8.000000}{9.600000}\selectfont\catcode`\^=\active\def^{\ifmmode\sp\else\^{}\fi}\catcode`\%=\active\def%{\%}1.0}}%
\end{pgfscope}%
\begin{pgfscope}%
\pgfsetbuttcap%
\pgfsetroundjoin%
\definecolor{currentfill}{rgb}{0.000000,0.000000,0.000000}%
\pgfsetfillcolor{currentfill}%
\pgfsetlinewidth{0.301125pt}%
\definecolor{currentstroke}{rgb}{0.000000,0.000000,0.000000}%
\pgfsetstrokecolor{currentstroke}%
\pgfsetdash{}{0pt}%
\pgfsys@defobject{currentmarker}{\pgfqpoint{0.000000in}{0.000000in}}{\pgfqpoint{0.000000in}{0.027778in}}{%
\pgfpathmoveto{\pgfqpoint{0.000000in}{0.000000in}}%
\pgfpathlineto{\pgfqpoint{0.000000in}{0.027778in}}%
\pgfusepath{stroke,fill}%
}%
\begin{pgfscope}%
\pgfsys@transformshift{0.575127in}{0.352669in}%
\pgfsys@useobject{currentmarker}{}%
\end{pgfscope}%
\end{pgfscope}%
\begin{pgfscope}%
\pgfsetbuttcap%
\pgfsetroundjoin%
\definecolor{currentfill}{rgb}{0.000000,0.000000,0.000000}%
\pgfsetfillcolor{currentfill}%
\pgfsetlinewidth{0.301125pt}%
\definecolor{currentstroke}{rgb}{0.000000,0.000000,0.000000}%
\pgfsetstrokecolor{currentstroke}%
\pgfsetdash{}{0pt}%
\pgfsys@defobject{currentmarker}{\pgfqpoint{0.000000in}{0.000000in}}{\pgfqpoint{0.000000in}{0.027778in}}{%
\pgfpathmoveto{\pgfqpoint{0.000000in}{0.000000in}}%
\pgfpathlineto{\pgfqpoint{0.000000in}{0.027778in}}%
\pgfusepath{stroke,fill}%
}%
\begin{pgfscope}%
\pgfsys@transformshift{0.642634in}{0.352669in}%
\pgfsys@useobject{currentmarker}{}%
\end{pgfscope}%
\end{pgfscope}%
\begin{pgfscope}%
\pgfsetbuttcap%
\pgfsetroundjoin%
\definecolor{currentfill}{rgb}{0.000000,0.000000,0.000000}%
\pgfsetfillcolor{currentfill}%
\pgfsetlinewidth{0.301125pt}%
\definecolor{currentstroke}{rgb}{0.000000,0.000000,0.000000}%
\pgfsetstrokecolor{currentstroke}%
\pgfsetdash{}{0pt}%
\pgfsys@defobject{currentmarker}{\pgfqpoint{0.000000in}{0.000000in}}{\pgfqpoint{0.000000in}{0.027778in}}{%
\pgfpathmoveto{\pgfqpoint{0.000000in}{0.000000in}}%
\pgfpathlineto{\pgfqpoint{0.000000in}{0.027778in}}%
\pgfusepath{stroke,fill}%
}%
\begin{pgfscope}%
\pgfsys@transformshift{0.710142in}{0.352669in}%
\pgfsys@useobject{currentmarker}{}%
\end{pgfscope}%
\end{pgfscope}%
\begin{pgfscope}%
\pgfsetbuttcap%
\pgfsetroundjoin%
\definecolor{currentfill}{rgb}{0.000000,0.000000,0.000000}%
\pgfsetfillcolor{currentfill}%
\pgfsetlinewidth{0.301125pt}%
\definecolor{currentstroke}{rgb}{0.000000,0.000000,0.000000}%
\pgfsetstrokecolor{currentstroke}%
\pgfsetdash{}{0pt}%
\pgfsys@defobject{currentmarker}{\pgfqpoint{0.000000in}{0.000000in}}{\pgfqpoint{0.000000in}{0.027778in}}{%
\pgfpathmoveto{\pgfqpoint{0.000000in}{0.000000in}}%
\pgfpathlineto{\pgfqpoint{0.000000in}{0.027778in}}%
\pgfusepath{stroke,fill}%
}%
\begin{pgfscope}%
\pgfsys@transformshift{0.777649in}{0.352669in}%
\pgfsys@useobject{currentmarker}{}%
\end{pgfscope}%
\end{pgfscope}%
\begin{pgfscope}%
\pgfsetbuttcap%
\pgfsetroundjoin%
\definecolor{currentfill}{rgb}{0.000000,0.000000,0.000000}%
\pgfsetfillcolor{currentfill}%
\pgfsetlinewidth{0.301125pt}%
\definecolor{currentstroke}{rgb}{0.000000,0.000000,0.000000}%
\pgfsetstrokecolor{currentstroke}%
\pgfsetdash{}{0pt}%
\pgfsys@defobject{currentmarker}{\pgfqpoint{0.000000in}{0.000000in}}{\pgfqpoint{0.000000in}{0.027778in}}{%
\pgfpathmoveto{\pgfqpoint{0.000000in}{0.000000in}}%
\pgfpathlineto{\pgfqpoint{0.000000in}{0.027778in}}%
\pgfusepath{stroke,fill}%
}%
\begin{pgfscope}%
\pgfsys@transformshift{0.912664in}{0.352669in}%
\pgfsys@useobject{currentmarker}{}%
\end{pgfscope}%
\end{pgfscope}%
\begin{pgfscope}%
\pgfsetbuttcap%
\pgfsetroundjoin%
\definecolor{currentfill}{rgb}{0.000000,0.000000,0.000000}%
\pgfsetfillcolor{currentfill}%
\pgfsetlinewidth{0.301125pt}%
\definecolor{currentstroke}{rgb}{0.000000,0.000000,0.000000}%
\pgfsetstrokecolor{currentstroke}%
\pgfsetdash{}{0pt}%
\pgfsys@defobject{currentmarker}{\pgfqpoint{0.000000in}{0.000000in}}{\pgfqpoint{0.000000in}{0.027778in}}{%
\pgfpathmoveto{\pgfqpoint{0.000000in}{0.000000in}}%
\pgfpathlineto{\pgfqpoint{0.000000in}{0.027778in}}%
\pgfusepath{stroke,fill}%
}%
\begin{pgfscope}%
\pgfsys@transformshift{0.980172in}{0.352669in}%
\pgfsys@useobject{currentmarker}{}%
\end{pgfscope}%
\end{pgfscope}%
\begin{pgfscope}%
\pgfsetbuttcap%
\pgfsetroundjoin%
\definecolor{currentfill}{rgb}{0.000000,0.000000,0.000000}%
\pgfsetfillcolor{currentfill}%
\pgfsetlinewidth{0.301125pt}%
\definecolor{currentstroke}{rgb}{0.000000,0.000000,0.000000}%
\pgfsetstrokecolor{currentstroke}%
\pgfsetdash{}{0pt}%
\pgfsys@defobject{currentmarker}{\pgfqpoint{0.000000in}{0.000000in}}{\pgfqpoint{0.000000in}{0.027778in}}{%
\pgfpathmoveto{\pgfqpoint{0.000000in}{0.000000in}}%
\pgfpathlineto{\pgfqpoint{0.000000in}{0.027778in}}%
\pgfusepath{stroke,fill}%
}%
\begin{pgfscope}%
\pgfsys@transformshift{1.047679in}{0.352669in}%
\pgfsys@useobject{currentmarker}{}%
\end{pgfscope}%
\end{pgfscope}%
\begin{pgfscope}%
\pgfsetbuttcap%
\pgfsetroundjoin%
\definecolor{currentfill}{rgb}{0.000000,0.000000,0.000000}%
\pgfsetfillcolor{currentfill}%
\pgfsetlinewidth{0.301125pt}%
\definecolor{currentstroke}{rgb}{0.000000,0.000000,0.000000}%
\pgfsetstrokecolor{currentstroke}%
\pgfsetdash{}{0pt}%
\pgfsys@defobject{currentmarker}{\pgfqpoint{0.000000in}{0.000000in}}{\pgfqpoint{0.000000in}{0.027778in}}{%
\pgfpathmoveto{\pgfqpoint{0.000000in}{0.000000in}}%
\pgfpathlineto{\pgfqpoint{0.000000in}{0.027778in}}%
\pgfusepath{stroke,fill}%
}%
\begin{pgfscope}%
\pgfsys@transformshift{1.115187in}{0.352669in}%
\pgfsys@useobject{currentmarker}{}%
\end{pgfscope}%
\end{pgfscope}%
\begin{pgfscope}%
\pgfsetbuttcap%
\pgfsetroundjoin%
\definecolor{currentfill}{rgb}{0.000000,0.000000,0.000000}%
\pgfsetfillcolor{currentfill}%
\pgfsetlinewidth{0.301125pt}%
\definecolor{currentstroke}{rgb}{0.000000,0.000000,0.000000}%
\pgfsetstrokecolor{currentstroke}%
\pgfsetdash{}{0pt}%
\pgfsys@defobject{currentmarker}{\pgfqpoint{0.000000in}{0.000000in}}{\pgfqpoint{0.000000in}{0.027778in}}{%
\pgfpathmoveto{\pgfqpoint{0.000000in}{0.000000in}}%
\pgfpathlineto{\pgfqpoint{0.000000in}{0.027778in}}%
\pgfusepath{stroke,fill}%
}%
\begin{pgfscope}%
\pgfsys@transformshift{1.250201in}{0.352669in}%
\pgfsys@useobject{currentmarker}{}%
\end{pgfscope}%
\end{pgfscope}%
\begin{pgfscope}%
\pgfsetbuttcap%
\pgfsetroundjoin%
\definecolor{currentfill}{rgb}{0.000000,0.000000,0.000000}%
\pgfsetfillcolor{currentfill}%
\pgfsetlinewidth{0.301125pt}%
\definecolor{currentstroke}{rgb}{0.000000,0.000000,0.000000}%
\pgfsetstrokecolor{currentstroke}%
\pgfsetdash{}{0pt}%
\pgfsys@defobject{currentmarker}{\pgfqpoint{0.000000in}{0.000000in}}{\pgfqpoint{0.000000in}{0.027778in}}{%
\pgfpathmoveto{\pgfqpoint{0.000000in}{0.000000in}}%
\pgfpathlineto{\pgfqpoint{0.000000in}{0.027778in}}%
\pgfusepath{stroke,fill}%
}%
\begin{pgfscope}%
\pgfsys@transformshift{1.317709in}{0.352669in}%
\pgfsys@useobject{currentmarker}{}%
\end{pgfscope}%
\end{pgfscope}%
\begin{pgfscope}%
\pgfsetbuttcap%
\pgfsetroundjoin%
\definecolor{currentfill}{rgb}{0.000000,0.000000,0.000000}%
\pgfsetfillcolor{currentfill}%
\pgfsetlinewidth{0.301125pt}%
\definecolor{currentstroke}{rgb}{0.000000,0.000000,0.000000}%
\pgfsetstrokecolor{currentstroke}%
\pgfsetdash{}{0pt}%
\pgfsys@defobject{currentmarker}{\pgfqpoint{0.000000in}{0.000000in}}{\pgfqpoint{0.000000in}{0.027778in}}{%
\pgfpathmoveto{\pgfqpoint{0.000000in}{0.000000in}}%
\pgfpathlineto{\pgfqpoint{0.000000in}{0.027778in}}%
\pgfusepath{stroke,fill}%
}%
\begin{pgfscope}%
\pgfsys@transformshift{1.385216in}{0.352669in}%
\pgfsys@useobject{currentmarker}{}%
\end{pgfscope}%
\end{pgfscope}%
\begin{pgfscope}%
\pgfsetbuttcap%
\pgfsetroundjoin%
\definecolor{currentfill}{rgb}{0.000000,0.000000,0.000000}%
\pgfsetfillcolor{currentfill}%
\pgfsetlinewidth{0.301125pt}%
\definecolor{currentstroke}{rgb}{0.000000,0.000000,0.000000}%
\pgfsetstrokecolor{currentstroke}%
\pgfsetdash{}{0pt}%
\pgfsys@defobject{currentmarker}{\pgfqpoint{0.000000in}{0.000000in}}{\pgfqpoint{0.000000in}{0.027778in}}{%
\pgfpathmoveto{\pgfqpoint{0.000000in}{0.000000in}}%
\pgfpathlineto{\pgfqpoint{0.000000in}{0.027778in}}%
\pgfusepath{stroke,fill}%
}%
\begin{pgfscope}%
\pgfsys@transformshift{1.452724in}{0.352669in}%
\pgfsys@useobject{currentmarker}{}%
\end{pgfscope}%
\end{pgfscope}%
\begin{pgfscope}%
\pgfsetbuttcap%
\pgfsetroundjoin%
\definecolor{currentfill}{rgb}{0.000000,0.000000,0.000000}%
\pgfsetfillcolor{currentfill}%
\pgfsetlinewidth{0.301125pt}%
\definecolor{currentstroke}{rgb}{0.000000,0.000000,0.000000}%
\pgfsetstrokecolor{currentstroke}%
\pgfsetdash{}{0pt}%
\pgfsys@defobject{currentmarker}{\pgfqpoint{0.000000in}{0.000000in}}{\pgfqpoint{0.000000in}{0.027778in}}{%
\pgfpathmoveto{\pgfqpoint{0.000000in}{0.000000in}}%
\pgfpathlineto{\pgfqpoint{0.000000in}{0.027778in}}%
\pgfusepath{stroke,fill}%
}%
\begin{pgfscope}%
\pgfsys@transformshift{1.587739in}{0.352669in}%
\pgfsys@useobject{currentmarker}{}%
\end{pgfscope}%
\end{pgfscope}%
\begin{pgfscope}%
\pgfsetbuttcap%
\pgfsetroundjoin%
\definecolor{currentfill}{rgb}{0.000000,0.000000,0.000000}%
\pgfsetfillcolor{currentfill}%
\pgfsetlinewidth{0.301125pt}%
\definecolor{currentstroke}{rgb}{0.000000,0.000000,0.000000}%
\pgfsetstrokecolor{currentstroke}%
\pgfsetdash{}{0pt}%
\pgfsys@defobject{currentmarker}{\pgfqpoint{0.000000in}{0.000000in}}{\pgfqpoint{0.000000in}{0.027778in}}{%
\pgfpathmoveto{\pgfqpoint{0.000000in}{0.000000in}}%
\pgfpathlineto{\pgfqpoint{0.000000in}{0.027778in}}%
\pgfusepath{stroke,fill}%
}%
\begin{pgfscope}%
\pgfsys@transformshift{1.655246in}{0.352669in}%
\pgfsys@useobject{currentmarker}{}%
\end{pgfscope}%
\end{pgfscope}%
\begin{pgfscope}%
\pgfsetbuttcap%
\pgfsetroundjoin%
\definecolor{currentfill}{rgb}{0.000000,0.000000,0.000000}%
\pgfsetfillcolor{currentfill}%
\pgfsetlinewidth{0.301125pt}%
\definecolor{currentstroke}{rgb}{0.000000,0.000000,0.000000}%
\pgfsetstrokecolor{currentstroke}%
\pgfsetdash{}{0pt}%
\pgfsys@defobject{currentmarker}{\pgfqpoint{0.000000in}{0.000000in}}{\pgfqpoint{0.000000in}{0.027778in}}{%
\pgfpathmoveto{\pgfqpoint{0.000000in}{0.000000in}}%
\pgfpathlineto{\pgfqpoint{0.000000in}{0.027778in}}%
\pgfusepath{stroke,fill}%
}%
\begin{pgfscope}%
\pgfsys@transformshift{1.722754in}{0.352669in}%
\pgfsys@useobject{currentmarker}{}%
\end{pgfscope}%
\end{pgfscope}%
\begin{pgfscope}%
\pgfsetbuttcap%
\pgfsetroundjoin%
\definecolor{currentfill}{rgb}{0.000000,0.000000,0.000000}%
\pgfsetfillcolor{currentfill}%
\pgfsetlinewidth{0.301125pt}%
\definecolor{currentstroke}{rgb}{0.000000,0.000000,0.000000}%
\pgfsetstrokecolor{currentstroke}%
\pgfsetdash{}{0pt}%
\pgfsys@defobject{currentmarker}{\pgfqpoint{0.000000in}{0.000000in}}{\pgfqpoint{0.000000in}{0.027778in}}{%
\pgfpathmoveto{\pgfqpoint{0.000000in}{0.000000in}}%
\pgfpathlineto{\pgfqpoint{0.000000in}{0.027778in}}%
\pgfusepath{stroke,fill}%
}%
\begin{pgfscope}%
\pgfsys@transformshift{1.790261in}{0.352669in}%
\pgfsys@useobject{currentmarker}{}%
\end{pgfscope}%
\end{pgfscope}%
\begin{pgfscope}%
\definecolor{textcolor}{rgb}{0.000000,0.000000,0.000000}%
\pgfsetstrokecolor{textcolor}%
\pgfsetfillcolor{textcolor}%
\pgftext[x=1.182694in,y=0.140972in,,top]{\color{textcolor}{\sffamily\fontsize{9.000000}{10.800000}\selectfont\catcode`\^=\active\def^{\ifmmode\sp\else\^{}\fi}\catcode`\%=\active\def%{\%}$x$}}%
\end{pgfscope}%
\begin{pgfscope}%
\pgfsetbuttcap%
\pgfsetroundjoin%
\definecolor{currentfill}{rgb}{0.000000,0.000000,0.000000}%
\pgfsetfillcolor{currentfill}%
\pgfsetlinewidth{0.501875pt}%
\definecolor{currentstroke}{rgb}{0.000000,0.000000,0.000000}%
\pgfsetstrokecolor{currentstroke}%
\pgfsetdash{}{0pt}%
\pgfsys@defobject{currentmarker}{\pgfqpoint{0.000000in}{0.000000in}}{\pgfqpoint{0.041667in}{0.000000in}}{%
\pgfpathmoveto{\pgfqpoint{0.000000in}{0.000000in}}%
\pgfpathlineto{\pgfqpoint{0.041667in}{0.000000in}}%
\pgfusepath{stroke,fill}%
}%
\begin{pgfscope}%
\pgfsys@transformshift{0.507620in}{0.352669in}%
\pgfsys@useobject{currentmarker}{}%
\end{pgfscope}%
\end{pgfscope}%
\begin{pgfscope}%
\definecolor{textcolor}{rgb}{0.000000,0.000000,0.000000}%
\pgfsetstrokecolor{textcolor}%
\pgfsetfillcolor{textcolor}%
\pgftext[x=0.196527in, y=0.310460in, left, base]{\color{textcolor}{\sffamily\fontsize{8.000000}{9.600000}\selectfont\catcode`\^=\active\def^{\ifmmode\sp\else\^{}\fi}\catcode`\%=\active\def%{\%}\ensuremath{-}1.5}}%
\end{pgfscope}%
\begin{pgfscope}%
\pgfsetbuttcap%
\pgfsetroundjoin%
\definecolor{currentfill}{rgb}{0.000000,0.000000,0.000000}%
\pgfsetfillcolor{currentfill}%
\pgfsetlinewidth{0.501875pt}%
\definecolor{currentstroke}{rgb}{0.000000,0.000000,0.000000}%
\pgfsetstrokecolor{currentstroke}%
\pgfsetdash{}{0pt}%
\pgfsys@defobject{currentmarker}{\pgfqpoint{0.000000in}{0.000000in}}{\pgfqpoint{0.041667in}{0.000000in}}{%
\pgfpathmoveto{\pgfqpoint{0.000000in}{0.000000in}}%
\pgfpathlineto{\pgfqpoint{0.041667in}{0.000000in}}%
\pgfusepath{stroke,fill}%
}%
\begin{pgfscope}%
\pgfsys@transformshift{0.507620in}{0.605669in}%
\pgfsys@useobject{currentmarker}{}%
\end{pgfscope}%
\end{pgfscope}%
\begin{pgfscope}%
\definecolor{textcolor}{rgb}{0.000000,0.000000,0.000000}%
\pgfsetstrokecolor{textcolor}%
\pgfsetfillcolor{textcolor}%
\pgftext[x=0.196527in, y=0.563460in, left, base]{\color{textcolor}{\sffamily\fontsize{8.000000}{9.600000}\selectfont\catcode`\^=\active\def^{\ifmmode\sp\else\^{}\fi}\catcode`\%=\active\def%{\%}\ensuremath{-}1.0}}%
\end{pgfscope}%
\begin{pgfscope}%
\pgfsetbuttcap%
\pgfsetroundjoin%
\definecolor{currentfill}{rgb}{0.000000,0.000000,0.000000}%
\pgfsetfillcolor{currentfill}%
\pgfsetlinewidth{0.501875pt}%
\definecolor{currentstroke}{rgb}{0.000000,0.000000,0.000000}%
\pgfsetstrokecolor{currentstroke}%
\pgfsetdash{}{0pt}%
\pgfsys@defobject{currentmarker}{\pgfqpoint{0.000000in}{0.000000in}}{\pgfqpoint{0.041667in}{0.000000in}}{%
\pgfpathmoveto{\pgfqpoint{0.000000in}{0.000000in}}%
\pgfpathlineto{\pgfqpoint{0.041667in}{0.000000in}}%
\pgfusepath{stroke,fill}%
}%
\begin{pgfscope}%
\pgfsys@transformshift{0.507620in}{0.858669in}%
\pgfsys@useobject{currentmarker}{}%
\end{pgfscope}%
\end{pgfscope}%
\begin{pgfscope}%
\definecolor{textcolor}{rgb}{0.000000,0.000000,0.000000}%
\pgfsetstrokecolor{textcolor}%
\pgfsetfillcolor{textcolor}%
\pgftext[x=0.196527in, y=0.816460in, left, base]{\color{textcolor}{\sffamily\fontsize{8.000000}{9.600000}\selectfont\catcode`\^=\active\def^{\ifmmode\sp\else\^{}\fi}\catcode`\%=\active\def%{\%}\ensuremath{-}0.5}}%
\end{pgfscope}%
\begin{pgfscope}%
\pgfsetbuttcap%
\pgfsetroundjoin%
\definecolor{currentfill}{rgb}{0.000000,0.000000,0.000000}%
\pgfsetfillcolor{currentfill}%
\pgfsetlinewidth{0.501875pt}%
\definecolor{currentstroke}{rgb}{0.000000,0.000000,0.000000}%
\pgfsetstrokecolor{currentstroke}%
\pgfsetdash{}{0pt}%
\pgfsys@defobject{currentmarker}{\pgfqpoint{0.000000in}{0.000000in}}{\pgfqpoint{0.041667in}{0.000000in}}{%
\pgfpathmoveto{\pgfqpoint{0.000000in}{0.000000in}}%
\pgfpathlineto{\pgfqpoint{0.041667in}{0.000000in}}%
\pgfusepath{stroke,fill}%
}%
\begin{pgfscope}%
\pgfsys@transformshift{0.507620in}{1.111669in}%
\pgfsys@useobject{currentmarker}{}%
\end{pgfscope}%
\end{pgfscope}%
\begin{pgfscope}%
\definecolor{textcolor}{rgb}{0.000000,0.000000,0.000000}%
\pgfsetstrokecolor{textcolor}%
\pgfsetfillcolor{textcolor}%
\pgftext[x=0.282305in, y=1.069460in, left, base]{\color{textcolor}{\sffamily\fontsize{8.000000}{9.600000}\selectfont\catcode`\^=\active\def^{\ifmmode\sp\else\^{}\fi}\catcode`\%=\active\def%{\%}0.0}}%
\end{pgfscope}%
\begin{pgfscope}%
\pgfsetbuttcap%
\pgfsetroundjoin%
\definecolor{currentfill}{rgb}{0.000000,0.000000,0.000000}%
\pgfsetfillcolor{currentfill}%
\pgfsetlinewidth{0.501875pt}%
\definecolor{currentstroke}{rgb}{0.000000,0.000000,0.000000}%
\pgfsetstrokecolor{currentstroke}%
\pgfsetdash{}{0pt}%
\pgfsys@defobject{currentmarker}{\pgfqpoint{0.000000in}{0.000000in}}{\pgfqpoint{0.041667in}{0.000000in}}{%
\pgfpathmoveto{\pgfqpoint{0.000000in}{0.000000in}}%
\pgfpathlineto{\pgfqpoint{0.041667in}{0.000000in}}%
\pgfusepath{stroke,fill}%
}%
\begin{pgfscope}%
\pgfsys@transformshift{0.507620in}{1.364669in}%
\pgfsys@useobject{currentmarker}{}%
\end{pgfscope}%
\end{pgfscope}%
\begin{pgfscope}%
\definecolor{textcolor}{rgb}{0.000000,0.000000,0.000000}%
\pgfsetstrokecolor{textcolor}%
\pgfsetfillcolor{textcolor}%
\pgftext[x=0.282305in, y=1.322460in, left, base]{\color{textcolor}{\sffamily\fontsize{8.000000}{9.600000}\selectfont\catcode`\^=\active\def^{\ifmmode\sp\else\^{}\fi}\catcode`\%=\active\def%{\%}0.5}}%
\end{pgfscope}%
\begin{pgfscope}%
\pgfsetbuttcap%
\pgfsetroundjoin%
\definecolor{currentfill}{rgb}{0.000000,0.000000,0.000000}%
\pgfsetfillcolor{currentfill}%
\pgfsetlinewidth{0.501875pt}%
\definecolor{currentstroke}{rgb}{0.000000,0.000000,0.000000}%
\pgfsetstrokecolor{currentstroke}%
\pgfsetdash{}{0pt}%
\pgfsys@defobject{currentmarker}{\pgfqpoint{0.000000in}{0.000000in}}{\pgfqpoint{0.041667in}{0.000000in}}{%
\pgfpathmoveto{\pgfqpoint{0.000000in}{0.000000in}}%
\pgfpathlineto{\pgfqpoint{0.041667in}{0.000000in}}%
\pgfusepath{stroke,fill}%
}%
\begin{pgfscope}%
\pgfsys@transformshift{0.507620in}{1.617669in}%
\pgfsys@useobject{currentmarker}{}%
\end{pgfscope}%
\end{pgfscope}%
\begin{pgfscope}%
\definecolor{textcolor}{rgb}{0.000000,0.000000,0.000000}%
\pgfsetstrokecolor{textcolor}%
\pgfsetfillcolor{textcolor}%
\pgftext[x=0.282305in, y=1.575460in, left, base]{\color{textcolor}{\sffamily\fontsize{8.000000}{9.600000}\selectfont\catcode`\^=\active\def^{\ifmmode\sp\else\^{}\fi}\catcode`\%=\active\def%{\%}1.0}}%
\end{pgfscope}%
\begin{pgfscope}%
\pgfsetbuttcap%
\pgfsetroundjoin%
\definecolor{currentfill}{rgb}{0.000000,0.000000,0.000000}%
\pgfsetfillcolor{currentfill}%
\pgfsetlinewidth{0.501875pt}%
\definecolor{currentstroke}{rgb}{0.000000,0.000000,0.000000}%
\pgfsetstrokecolor{currentstroke}%
\pgfsetdash{}{0pt}%
\pgfsys@defobject{currentmarker}{\pgfqpoint{0.000000in}{0.000000in}}{\pgfqpoint{0.041667in}{0.000000in}}{%
\pgfpathmoveto{\pgfqpoint{0.000000in}{0.000000in}}%
\pgfpathlineto{\pgfqpoint{0.041667in}{0.000000in}}%
\pgfusepath{stroke,fill}%
}%
\begin{pgfscope}%
\pgfsys@transformshift{0.507620in}{1.870669in}%
\pgfsys@useobject{currentmarker}{}%
\end{pgfscope}%
\end{pgfscope}%
\begin{pgfscope}%
\definecolor{textcolor}{rgb}{0.000000,0.000000,0.000000}%
\pgfsetstrokecolor{textcolor}%
\pgfsetfillcolor{textcolor}%
\pgftext[x=0.282305in, y=1.828460in, left, base]{\color{textcolor}{\sffamily\fontsize{8.000000}{9.600000}\selectfont\catcode`\^=\active\def^{\ifmmode\sp\else\^{}\fi}\catcode`\%=\active\def%{\%}1.5}}%
\end{pgfscope}%
\begin{pgfscope}%
\pgfsetbuttcap%
\pgfsetroundjoin%
\definecolor{currentfill}{rgb}{0.000000,0.000000,0.000000}%
\pgfsetfillcolor{currentfill}%
\pgfsetlinewidth{0.501875pt}%
\definecolor{currentstroke}{rgb}{0.000000,0.000000,0.000000}%
\pgfsetstrokecolor{currentstroke}%
\pgfsetdash{}{0pt}%
\pgfsys@defobject{currentmarker}{\pgfqpoint{0.000000in}{0.000000in}}{\pgfqpoint{0.041667in}{0.000000in}}{%
\pgfpathmoveto{\pgfqpoint{0.000000in}{0.000000in}}%
\pgfpathlineto{\pgfqpoint{0.041667in}{0.000000in}}%
\pgfusepath{stroke,fill}%
}%
\begin{pgfscope}%
\pgfsys@transformshift{0.507620in}{2.123669in}%
\pgfsys@useobject{currentmarker}{}%
\end{pgfscope}%
\end{pgfscope}%
\begin{pgfscope}%
\definecolor{textcolor}{rgb}{0.000000,0.000000,0.000000}%
\pgfsetstrokecolor{textcolor}%
\pgfsetfillcolor{textcolor}%
\pgftext[x=0.282305in, y=2.081460in, left, base]{\color{textcolor}{\sffamily\fontsize{8.000000}{9.600000}\selectfont\catcode`\^=\active\def^{\ifmmode\sp\else\^{}\fi}\catcode`\%=\active\def%{\%}2.0}}%
\end{pgfscope}%
\begin{pgfscope}%
\pgfsetbuttcap%
\pgfsetroundjoin%
\definecolor{currentfill}{rgb}{0.000000,0.000000,0.000000}%
\pgfsetfillcolor{currentfill}%
\pgfsetlinewidth{0.301125pt}%
\definecolor{currentstroke}{rgb}{0.000000,0.000000,0.000000}%
\pgfsetstrokecolor{currentstroke}%
\pgfsetdash{}{0pt}%
\pgfsys@defobject{currentmarker}{\pgfqpoint{0.000000in}{0.000000in}}{\pgfqpoint{0.027778in}{0.000000in}}{%
\pgfpathmoveto{\pgfqpoint{0.000000in}{0.000000in}}%
\pgfpathlineto{\pgfqpoint{0.027778in}{0.000000in}}%
\pgfusepath{stroke,fill}%
}%
\begin{pgfscope}%
\pgfsys@transformshift{0.507620in}{0.403269in}%
\pgfsys@useobject{currentmarker}{}%
\end{pgfscope}%
\end{pgfscope}%
\begin{pgfscope}%
\pgfsetbuttcap%
\pgfsetroundjoin%
\definecolor{currentfill}{rgb}{0.000000,0.000000,0.000000}%
\pgfsetfillcolor{currentfill}%
\pgfsetlinewidth{0.301125pt}%
\definecolor{currentstroke}{rgb}{0.000000,0.000000,0.000000}%
\pgfsetstrokecolor{currentstroke}%
\pgfsetdash{}{0pt}%
\pgfsys@defobject{currentmarker}{\pgfqpoint{0.000000in}{0.000000in}}{\pgfqpoint{0.027778in}{0.000000in}}{%
\pgfpathmoveto{\pgfqpoint{0.000000in}{0.000000in}}%
\pgfpathlineto{\pgfqpoint{0.027778in}{0.000000in}}%
\pgfusepath{stroke,fill}%
}%
\begin{pgfscope}%
\pgfsys@transformshift{0.507620in}{0.453869in}%
\pgfsys@useobject{currentmarker}{}%
\end{pgfscope}%
\end{pgfscope}%
\begin{pgfscope}%
\pgfsetbuttcap%
\pgfsetroundjoin%
\definecolor{currentfill}{rgb}{0.000000,0.000000,0.000000}%
\pgfsetfillcolor{currentfill}%
\pgfsetlinewidth{0.301125pt}%
\definecolor{currentstroke}{rgb}{0.000000,0.000000,0.000000}%
\pgfsetstrokecolor{currentstroke}%
\pgfsetdash{}{0pt}%
\pgfsys@defobject{currentmarker}{\pgfqpoint{0.000000in}{0.000000in}}{\pgfqpoint{0.027778in}{0.000000in}}{%
\pgfpathmoveto{\pgfqpoint{0.000000in}{0.000000in}}%
\pgfpathlineto{\pgfqpoint{0.027778in}{0.000000in}}%
\pgfusepath{stroke,fill}%
}%
\begin{pgfscope}%
\pgfsys@transformshift{0.507620in}{0.504469in}%
\pgfsys@useobject{currentmarker}{}%
\end{pgfscope}%
\end{pgfscope}%
\begin{pgfscope}%
\pgfsetbuttcap%
\pgfsetroundjoin%
\definecolor{currentfill}{rgb}{0.000000,0.000000,0.000000}%
\pgfsetfillcolor{currentfill}%
\pgfsetlinewidth{0.301125pt}%
\definecolor{currentstroke}{rgb}{0.000000,0.000000,0.000000}%
\pgfsetstrokecolor{currentstroke}%
\pgfsetdash{}{0pt}%
\pgfsys@defobject{currentmarker}{\pgfqpoint{0.000000in}{0.000000in}}{\pgfqpoint{0.027778in}{0.000000in}}{%
\pgfpathmoveto{\pgfqpoint{0.000000in}{0.000000in}}%
\pgfpathlineto{\pgfqpoint{0.027778in}{0.000000in}}%
\pgfusepath{stroke,fill}%
}%
\begin{pgfscope}%
\pgfsys@transformshift{0.507620in}{0.555069in}%
\pgfsys@useobject{currentmarker}{}%
\end{pgfscope}%
\end{pgfscope}%
\begin{pgfscope}%
\pgfsetbuttcap%
\pgfsetroundjoin%
\definecolor{currentfill}{rgb}{0.000000,0.000000,0.000000}%
\pgfsetfillcolor{currentfill}%
\pgfsetlinewidth{0.301125pt}%
\definecolor{currentstroke}{rgb}{0.000000,0.000000,0.000000}%
\pgfsetstrokecolor{currentstroke}%
\pgfsetdash{}{0pt}%
\pgfsys@defobject{currentmarker}{\pgfqpoint{0.000000in}{0.000000in}}{\pgfqpoint{0.027778in}{0.000000in}}{%
\pgfpathmoveto{\pgfqpoint{0.000000in}{0.000000in}}%
\pgfpathlineto{\pgfqpoint{0.027778in}{0.000000in}}%
\pgfusepath{stroke,fill}%
}%
\begin{pgfscope}%
\pgfsys@transformshift{0.507620in}{0.656269in}%
\pgfsys@useobject{currentmarker}{}%
\end{pgfscope}%
\end{pgfscope}%
\begin{pgfscope}%
\pgfsetbuttcap%
\pgfsetroundjoin%
\definecolor{currentfill}{rgb}{0.000000,0.000000,0.000000}%
\pgfsetfillcolor{currentfill}%
\pgfsetlinewidth{0.301125pt}%
\definecolor{currentstroke}{rgb}{0.000000,0.000000,0.000000}%
\pgfsetstrokecolor{currentstroke}%
\pgfsetdash{}{0pt}%
\pgfsys@defobject{currentmarker}{\pgfqpoint{0.000000in}{0.000000in}}{\pgfqpoint{0.027778in}{0.000000in}}{%
\pgfpathmoveto{\pgfqpoint{0.000000in}{0.000000in}}%
\pgfpathlineto{\pgfqpoint{0.027778in}{0.000000in}}%
\pgfusepath{stroke,fill}%
}%
\begin{pgfscope}%
\pgfsys@transformshift{0.507620in}{0.706869in}%
\pgfsys@useobject{currentmarker}{}%
\end{pgfscope}%
\end{pgfscope}%
\begin{pgfscope}%
\pgfsetbuttcap%
\pgfsetroundjoin%
\definecolor{currentfill}{rgb}{0.000000,0.000000,0.000000}%
\pgfsetfillcolor{currentfill}%
\pgfsetlinewidth{0.301125pt}%
\definecolor{currentstroke}{rgb}{0.000000,0.000000,0.000000}%
\pgfsetstrokecolor{currentstroke}%
\pgfsetdash{}{0pt}%
\pgfsys@defobject{currentmarker}{\pgfqpoint{0.000000in}{0.000000in}}{\pgfqpoint{0.027778in}{0.000000in}}{%
\pgfpathmoveto{\pgfqpoint{0.000000in}{0.000000in}}%
\pgfpathlineto{\pgfqpoint{0.027778in}{0.000000in}}%
\pgfusepath{stroke,fill}%
}%
\begin{pgfscope}%
\pgfsys@transformshift{0.507620in}{0.757469in}%
\pgfsys@useobject{currentmarker}{}%
\end{pgfscope}%
\end{pgfscope}%
\begin{pgfscope}%
\pgfsetbuttcap%
\pgfsetroundjoin%
\definecolor{currentfill}{rgb}{0.000000,0.000000,0.000000}%
\pgfsetfillcolor{currentfill}%
\pgfsetlinewidth{0.301125pt}%
\definecolor{currentstroke}{rgb}{0.000000,0.000000,0.000000}%
\pgfsetstrokecolor{currentstroke}%
\pgfsetdash{}{0pt}%
\pgfsys@defobject{currentmarker}{\pgfqpoint{0.000000in}{0.000000in}}{\pgfqpoint{0.027778in}{0.000000in}}{%
\pgfpathmoveto{\pgfqpoint{0.000000in}{0.000000in}}%
\pgfpathlineto{\pgfqpoint{0.027778in}{0.000000in}}%
\pgfusepath{stroke,fill}%
}%
\begin{pgfscope}%
\pgfsys@transformshift{0.507620in}{0.808069in}%
\pgfsys@useobject{currentmarker}{}%
\end{pgfscope}%
\end{pgfscope}%
\begin{pgfscope}%
\pgfsetbuttcap%
\pgfsetroundjoin%
\definecolor{currentfill}{rgb}{0.000000,0.000000,0.000000}%
\pgfsetfillcolor{currentfill}%
\pgfsetlinewidth{0.301125pt}%
\definecolor{currentstroke}{rgb}{0.000000,0.000000,0.000000}%
\pgfsetstrokecolor{currentstroke}%
\pgfsetdash{}{0pt}%
\pgfsys@defobject{currentmarker}{\pgfqpoint{0.000000in}{0.000000in}}{\pgfqpoint{0.027778in}{0.000000in}}{%
\pgfpathmoveto{\pgfqpoint{0.000000in}{0.000000in}}%
\pgfpathlineto{\pgfqpoint{0.027778in}{0.000000in}}%
\pgfusepath{stroke,fill}%
}%
\begin{pgfscope}%
\pgfsys@transformshift{0.507620in}{0.909269in}%
\pgfsys@useobject{currentmarker}{}%
\end{pgfscope}%
\end{pgfscope}%
\begin{pgfscope}%
\pgfsetbuttcap%
\pgfsetroundjoin%
\definecolor{currentfill}{rgb}{0.000000,0.000000,0.000000}%
\pgfsetfillcolor{currentfill}%
\pgfsetlinewidth{0.301125pt}%
\definecolor{currentstroke}{rgb}{0.000000,0.000000,0.000000}%
\pgfsetstrokecolor{currentstroke}%
\pgfsetdash{}{0pt}%
\pgfsys@defobject{currentmarker}{\pgfqpoint{0.000000in}{0.000000in}}{\pgfqpoint{0.027778in}{0.000000in}}{%
\pgfpathmoveto{\pgfqpoint{0.000000in}{0.000000in}}%
\pgfpathlineto{\pgfqpoint{0.027778in}{0.000000in}}%
\pgfusepath{stroke,fill}%
}%
\begin{pgfscope}%
\pgfsys@transformshift{0.507620in}{0.959869in}%
\pgfsys@useobject{currentmarker}{}%
\end{pgfscope}%
\end{pgfscope}%
\begin{pgfscope}%
\pgfsetbuttcap%
\pgfsetroundjoin%
\definecolor{currentfill}{rgb}{0.000000,0.000000,0.000000}%
\pgfsetfillcolor{currentfill}%
\pgfsetlinewidth{0.301125pt}%
\definecolor{currentstroke}{rgb}{0.000000,0.000000,0.000000}%
\pgfsetstrokecolor{currentstroke}%
\pgfsetdash{}{0pt}%
\pgfsys@defobject{currentmarker}{\pgfqpoint{0.000000in}{0.000000in}}{\pgfqpoint{0.027778in}{0.000000in}}{%
\pgfpathmoveto{\pgfqpoint{0.000000in}{0.000000in}}%
\pgfpathlineto{\pgfqpoint{0.027778in}{0.000000in}}%
\pgfusepath{stroke,fill}%
}%
\begin{pgfscope}%
\pgfsys@transformshift{0.507620in}{1.010469in}%
\pgfsys@useobject{currentmarker}{}%
\end{pgfscope}%
\end{pgfscope}%
\begin{pgfscope}%
\pgfsetbuttcap%
\pgfsetroundjoin%
\definecolor{currentfill}{rgb}{0.000000,0.000000,0.000000}%
\pgfsetfillcolor{currentfill}%
\pgfsetlinewidth{0.301125pt}%
\definecolor{currentstroke}{rgb}{0.000000,0.000000,0.000000}%
\pgfsetstrokecolor{currentstroke}%
\pgfsetdash{}{0pt}%
\pgfsys@defobject{currentmarker}{\pgfqpoint{0.000000in}{0.000000in}}{\pgfqpoint{0.027778in}{0.000000in}}{%
\pgfpathmoveto{\pgfqpoint{0.000000in}{0.000000in}}%
\pgfpathlineto{\pgfqpoint{0.027778in}{0.000000in}}%
\pgfusepath{stroke,fill}%
}%
\begin{pgfscope}%
\pgfsys@transformshift{0.507620in}{1.061069in}%
\pgfsys@useobject{currentmarker}{}%
\end{pgfscope}%
\end{pgfscope}%
\begin{pgfscope}%
\pgfsetbuttcap%
\pgfsetroundjoin%
\definecolor{currentfill}{rgb}{0.000000,0.000000,0.000000}%
\pgfsetfillcolor{currentfill}%
\pgfsetlinewidth{0.301125pt}%
\definecolor{currentstroke}{rgb}{0.000000,0.000000,0.000000}%
\pgfsetstrokecolor{currentstroke}%
\pgfsetdash{}{0pt}%
\pgfsys@defobject{currentmarker}{\pgfqpoint{0.000000in}{0.000000in}}{\pgfqpoint{0.027778in}{0.000000in}}{%
\pgfpathmoveto{\pgfqpoint{0.000000in}{0.000000in}}%
\pgfpathlineto{\pgfqpoint{0.027778in}{0.000000in}}%
\pgfusepath{stroke,fill}%
}%
\begin{pgfscope}%
\pgfsys@transformshift{0.507620in}{1.162269in}%
\pgfsys@useobject{currentmarker}{}%
\end{pgfscope}%
\end{pgfscope}%
\begin{pgfscope}%
\pgfsetbuttcap%
\pgfsetroundjoin%
\definecolor{currentfill}{rgb}{0.000000,0.000000,0.000000}%
\pgfsetfillcolor{currentfill}%
\pgfsetlinewidth{0.301125pt}%
\definecolor{currentstroke}{rgb}{0.000000,0.000000,0.000000}%
\pgfsetstrokecolor{currentstroke}%
\pgfsetdash{}{0pt}%
\pgfsys@defobject{currentmarker}{\pgfqpoint{0.000000in}{0.000000in}}{\pgfqpoint{0.027778in}{0.000000in}}{%
\pgfpathmoveto{\pgfqpoint{0.000000in}{0.000000in}}%
\pgfpathlineto{\pgfqpoint{0.027778in}{0.000000in}}%
\pgfusepath{stroke,fill}%
}%
\begin{pgfscope}%
\pgfsys@transformshift{0.507620in}{1.212869in}%
\pgfsys@useobject{currentmarker}{}%
\end{pgfscope}%
\end{pgfscope}%
\begin{pgfscope}%
\pgfsetbuttcap%
\pgfsetroundjoin%
\definecolor{currentfill}{rgb}{0.000000,0.000000,0.000000}%
\pgfsetfillcolor{currentfill}%
\pgfsetlinewidth{0.301125pt}%
\definecolor{currentstroke}{rgb}{0.000000,0.000000,0.000000}%
\pgfsetstrokecolor{currentstroke}%
\pgfsetdash{}{0pt}%
\pgfsys@defobject{currentmarker}{\pgfqpoint{0.000000in}{0.000000in}}{\pgfqpoint{0.027778in}{0.000000in}}{%
\pgfpathmoveto{\pgfqpoint{0.000000in}{0.000000in}}%
\pgfpathlineto{\pgfqpoint{0.027778in}{0.000000in}}%
\pgfusepath{stroke,fill}%
}%
\begin{pgfscope}%
\pgfsys@transformshift{0.507620in}{1.263469in}%
\pgfsys@useobject{currentmarker}{}%
\end{pgfscope}%
\end{pgfscope}%
\begin{pgfscope}%
\pgfsetbuttcap%
\pgfsetroundjoin%
\definecolor{currentfill}{rgb}{0.000000,0.000000,0.000000}%
\pgfsetfillcolor{currentfill}%
\pgfsetlinewidth{0.301125pt}%
\definecolor{currentstroke}{rgb}{0.000000,0.000000,0.000000}%
\pgfsetstrokecolor{currentstroke}%
\pgfsetdash{}{0pt}%
\pgfsys@defobject{currentmarker}{\pgfqpoint{0.000000in}{0.000000in}}{\pgfqpoint{0.027778in}{0.000000in}}{%
\pgfpathmoveto{\pgfqpoint{0.000000in}{0.000000in}}%
\pgfpathlineto{\pgfqpoint{0.027778in}{0.000000in}}%
\pgfusepath{stroke,fill}%
}%
\begin{pgfscope}%
\pgfsys@transformshift{0.507620in}{1.314069in}%
\pgfsys@useobject{currentmarker}{}%
\end{pgfscope}%
\end{pgfscope}%
\begin{pgfscope}%
\pgfsetbuttcap%
\pgfsetroundjoin%
\definecolor{currentfill}{rgb}{0.000000,0.000000,0.000000}%
\pgfsetfillcolor{currentfill}%
\pgfsetlinewidth{0.301125pt}%
\definecolor{currentstroke}{rgb}{0.000000,0.000000,0.000000}%
\pgfsetstrokecolor{currentstroke}%
\pgfsetdash{}{0pt}%
\pgfsys@defobject{currentmarker}{\pgfqpoint{0.000000in}{0.000000in}}{\pgfqpoint{0.027778in}{0.000000in}}{%
\pgfpathmoveto{\pgfqpoint{0.000000in}{0.000000in}}%
\pgfpathlineto{\pgfqpoint{0.027778in}{0.000000in}}%
\pgfusepath{stroke,fill}%
}%
\begin{pgfscope}%
\pgfsys@transformshift{0.507620in}{1.415269in}%
\pgfsys@useobject{currentmarker}{}%
\end{pgfscope}%
\end{pgfscope}%
\begin{pgfscope}%
\pgfsetbuttcap%
\pgfsetroundjoin%
\definecolor{currentfill}{rgb}{0.000000,0.000000,0.000000}%
\pgfsetfillcolor{currentfill}%
\pgfsetlinewidth{0.301125pt}%
\definecolor{currentstroke}{rgb}{0.000000,0.000000,0.000000}%
\pgfsetstrokecolor{currentstroke}%
\pgfsetdash{}{0pt}%
\pgfsys@defobject{currentmarker}{\pgfqpoint{0.000000in}{0.000000in}}{\pgfqpoint{0.027778in}{0.000000in}}{%
\pgfpathmoveto{\pgfqpoint{0.000000in}{0.000000in}}%
\pgfpathlineto{\pgfqpoint{0.027778in}{0.000000in}}%
\pgfusepath{stroke,fill}%
}%
\begin{pgfscope}%
\pgfsys@transformshift{0.507620in}{1.465869in}%
\pgfsys@useobject{currentmarker}{}%
\end{pgfscope}%
\end{pgfscope}%
\begin{pgfscope}%
\pgfsetbuttcap%
\pgfsetroundjoin%
\definecolor{currentfill}{rgb}{0.000000,0.000000,0.000000}%
\pgfsetfillcolor{currentfill}%
\pgfsetlinewidth{0.301125pt}%
\definecolor{currentstroke}{rgb}{0.000000,0.000000,0.000000}%
\pgfsetstrokecolor{currentstroke}%
\pgfsetdash{}{0pt}%
\pgfsys@defobject{currentmarker}{\pgfqpoint{0.000000in}{0.000000in}}{\pgfqpoint{0.027778in}{0.000000in}}{%
\pgfpathmoveto{\pgfqpoint{0.000000in}{0.000000in}}%
\pgfpathlineto{\pgfqpoint{0.027778in}{0.000000in}}%
\pgfusepath{stroke,fill}%
}%
\begin{pgfscope}%
\pgfsys@transformshift{0.507620in}{1.516469in}%
\pgfsys@useobject{currentmarker}{}%
\end{pgfscope}%
\end{pgfscope}%
\begin{pgfscope}%
\pgfsetbuttcap%
\pgfsetroundjoin%
\definecolor{currentfill}{rgb}{0.000000,0.000000,0.000000}%
\pgfsetfillcolor{currentfill}%
\pgfsetlinewidth{0.301125pt}%
\definecolor{currentstroke}{rgb}{0.000000,0.000000,0.000000}%
\pgfsetstrokecolor{currentstroke}%
\pgfsetdash{}{0pt}%
\pgfsys@defobject{currentmarker}{\pgfqpoint{0.000000in}{0.000000in}}{\pgfqpoint{0.027778in}{0.000000in}}{%
\pgfpathmoveto{\pgfqpoint{0.000000in}{0.000000in}}%
\pgfpathlineto{\pgfqpoint{0.027778in}{0.000000in}}%
\pgfusepath{stroke,fill}%
}%
\begin{pgfscope}%
\pgfsys@transformshift{0.507620in}{1.567069in}%
\pgfsys@useobject{currentmarker}{}%
\end{pgfscope}%
\end{pgfscope}%
\begin{pgfscope}%
\pgfsetbuttcap%
\pgfsetroundjoin%
\definecolor{currentfill}{rgb}{0.000000,0.000000,0.000000}%
\pgfsetfillcolor{currentfill}%
\pgfsetlinewidth{0.301125pt}%
\definecolor{currentstroke}{rgb}{0.000000,0.000000,0.000000}%
\pgfsetstrokecolor{currentstroke}%
\pgfsetdash{}{0pt}%
\pgfsys@defobject{currentmarker}{\pgfqpoint{0.000000in}{0.000000in}}{\pgfqpoint{0.027778in}{0.000000in}}{%
\pgfpathmoveto{\pgfqpoint{0.000000in}{0.000000in}}%
\pgfpathlineto{\pgfqpoint{0.027778in}{0.000000in}}%
\pgfusepath{stroke,fill}%
}%
\begin{pgfscope}%
\pgfsys@transformshift{0.507620in}{1.668269in}%
\pgfsys@useobject{currentmarker}{}%
\end{pgfscope}%
\end{pgfscope}%
\begin{pgfscope}%
\pgfsetbuttcap%
\pgfsetroundjoin%
\definecolor{currentfill}{rgb}{0.000000,0.000000,0.000000}%
\pgfsetfillcolor{currentfill}%
\pgfsetlinewidth{0.301125pt}%
\definecolor{currentstroke}{rgb}{0.000000,0.000000,0.000000}%
\pgfsetstrokecolor{currentstroke}%
\pgfsetdash{}{0pt}%
\pgfsys@defobject{currentmarker}{\pgfqpoint{0.000000in}{0.000000in}}{\pgfqpoint{0.027778in}{0.000000in}}{%
\pgfpathmoveto{\pgfqpoint{0.000000in}{0.000000in}}%
\pgfpathlineto{\pgfqpoint{0.027778in}{0.000000in}}%
\pgfusepath{stroke,fill}%
}%
\begin{pgfscope}%
\pgfsys@transformshift{0.507620in}{1.718869in}%
\pgfsys@useobject{currentmarker}{}%
\end{pgfscope}%
\end{pgfscope}%
\begin{pgfscope}%
\pgfsetbuttcap%
\pgfsetroundjoin%
\definecolor{currentfill}{rgb}{0.000000,0.000000,0.000000}%
\pgfsetfillcolor{currentfill}%
\pgfsetlinewidth{0.301125pt}%
\definecolor{currentstroke}{rgb}{0.000000,0.000000,0.000000}%
\pgfsetstrokecolor{currentstroke}%
\pgfsetdash{}{0pt}%
\pgfsys@defobject{currentmarker}{\pgfqpoint{0.000000in}{0.000000in}}{\pgfqpoint{0.027778in}{0.000000in}}{%
\pgfpathmoveto{\pgfqpoint{0.000000in}{0.000000in}}%
\pgfpathlineto{\pgfqpoint{0.027778in}{0.000000in}}%
\pgfusepath{stroke,fill}%
}%
\begin{pgfscope}%
\pgfsys@transformshift{0.507620in}{1.769469in}%
\pgfsys@useobject{currentmarker}{}%
\end{pgfscope}%
\end{pgfscope}%
\begin{pgfscope}%
\pgfsetbuttcap%
\pgfsetroundjoin%
\definecolor{currentfill}{rgb}{0.000000,0.000000,0.000000}%
\pgfsetfillcolor{currentfill}%
\pgfsetlinewidth{0.301125pt}%
\definecolor{currentstroke}{rgb}{0.000000,0.000000,0.000000}%
\pgfsetstrokecolor{currentstroke}%
\pgfsetdash{}{0pt}%
\pgfsys@defobject{currentmarker}{\pgfqpoint{0.000000in}{0.000000in}}{\pgfqpoint{0.027778in}{0.000000in}}{%
\pgfpathmoveto{\pgfqpoint{0.000000in}{0.000000in}}%
\pgfpathlineto{\pgfqpoint{0.027778in}{0.000000in}}%
\pgfusepath{stroke,fill}%
}%
\begin{pgfscope}%
\pgfsys@transformshift{0.507620in}{1.820069in}%
\pgfsys@useobject{currentmarker}{}%
\end{pgfscope}%
\end{pgfscope}%
\begin{pgfscope}%
\pgfsetbuttcap%
\pgfsetroundjoin%
\definecolor{currentfill}{rgb}{0.000000,0.000000,0.000000}%
\pgfsetfillcolor{currentfill}%
\pgfsetlinewidth{0.301125pt}%
\definecolor{currentstroke}{rgb}{0.000000,0.000000,0.000000}%
\pgfsetstrokecolor{currentstroke}%
\pgfsetdash{}{0pt}%
\pgfsys@defobject{currentmarker}{\pgfqpoint{0.000000in}{0.000000in}}{\pgfqpoint{0.027778in}{0.000000in}}{%
\pgfpathmoveto{\pgfqpoint{0.000000in}{0.000000in}}%
\pgfpathlineto{\pgfqpoint{0.027778in}{0.000000in}}%
\pgfusepath{stroke,fill}%
}%
\begin{pgfscope}%
\pgfsys@transformshift{0.507620in}{1.921269in}%
\pgfsys@useobject{currentmarker}{}%
\end{pgfscope}%
\end{pgfscope}%
\begin{pgfscope}%
\pgfsetbuttcap%
\pgfsetroundjoin%
\definecolor{currentfill}{rgb}{0.000000,0.000000,0.000000}%
\pgfsetfillcolor{currentfill}%
\pgfsetlinewidth{0.301125pt}%
\definecolor{currentstroke}{rgb}{0.000000,0.000000,0.000000}%
\pgfsetstrokecolor{currentstroke}%
\pgfsetdash{}{0pt}%
\pgfsys@defobject{currentmarker}{\pgfqpoint{0.000000in}{0.000000in}}{\pgfqpoint{0.027778in}{0.000000in}}{%
\pgfpathmoveto{\pgfqpoint{0.000000in}{0.000000in}}%
\pgfpathlineto{\pgfqpoint{0.027778in}{0.000000in}}%
\pgfusepath{stroke,fill}%
}%
\begin{pgfscope}%
\pgfsys@transformshift{0.507620in}{1.971869in}%
\pgfsys@useobject{currentmarker}{}%
\end{pgfscope}%
\end{pgfscope}%
\begin{pgfscope}%
\pgfsetbuttcap%
\pgfsetroundjoin%
\definecolor{currentfill}{rgb}{0.000000,0.000000,0.000000}%
\pgfsetfillcolor{currentfill}%
\pgfsetlinewidth{0.301125pt}%
\definecolor{currentstroke}{rgb}{0.000000,0.000000,0.000000}%
\pgfsetstrokecolor{currentstroke}%
\pgfsetdash{}{0pt}%
\pgfsys@defobject{currentmarker}{\pgfqpoint{0.000000in}{0.000000in}}{\pgfqpoint{0.027778in}{0.000000in}}{%
\pgfpathmoveto{\pgfqpoint{0.000000in}{0.000000in}}%
\pgfpathlineto{\pgfqpoint{0.027778in}{0.000000in}}%
\pgfusepath{stroke,fill}%
}%
\begin{pgfscope}%
\pgfsys@transformshift{0.507620in}{2.022469in}%
\pgfsys@useobject{currentmarker}{}%
\end{pgfscope}%
\end{pgfscope}%
\begin{pgfscope}%
\pgfsetbuttcap%
\pgfsetroundjoin%
\definecolor{currentfill}{rgb}{0.000000,0.000000,0.000000}%
\pgfsetfillcolor{currentfill}%
\pgfsetlinewidth{0.301125pt}%
\definecolor{currentstroke}{rgb}{0.000000,0.000000,0.000000}%
\pgfsetstrokecolor{currentstroke}%
\pgfsetdash{}{0pt}%
\pgfsys@defobject{currentmarker}{\pgfqpoint{0.000000in}{0.000000in}}{\pgfqpoint{0.027778in}{0.000000in}}{%
\pgfpathmoveto{\pgfqpoint{0.000000in}{0.000000in}}%
\pgfpathlineto{\pgfqpoint{0.027778in}{0.000000in}}%
\pgfusepath{stroke,fill}%
}%
\begin{pgfscope}%
\pgfsys@transformshift{0.507620in}{2.073069in}%
\pgfsys@useobject{currentmarker}{}%
\end{pgfscope}%
\end{pgfscope}%
\begin{pgfscope}%
\definecolor{textcolor}{rgb}{0.000000,0.000000,0.000000}%
\pgfsetstrokecolor{textcolor}%
\pgfsetfillcolor{textcolor}%
\pgftext[x=0.140972in,y=1.238169in,,bottom,rotate=90.000000]{\color{textcolor}{\sffamily\fontsize{9.000000}{10.800000}\selectfont\catcode`\^=\active\def^{\ifmmode\sp\else\^{}\fi}\catcode`\%=\active\def%{\%}$p_N(x)$}}%
\end{pgfscope}%
\begin{pgfscope}%
\pgfsetrectcap%
\pgfsetmiterjoin%
\pgfsetlinewidth{0.602250pt}%
\definecolor{currentstroke}{rgb}{0.000000,0.000000,0.000000}%
\pgfsetstrokecolor{currentstroke}%
\pgfsetdash{}{0pt}%
\pgfpathmoveto{\pgfqpoint{0.507620in}{0.352669in}}%
\pgfpathlineto{\pgfqpoint{0.507620in}{2.123669in}}%
\pgfusepath{stroke}%
\end{pgfscope}%
\begin{pgfscope}%
\pgfsetrectcap%
\pgfsetmiterjoin%
\pgfsetlinewidth{0.602250pt}%
\definecolor{currentstroke}{rgb}{0.000000,0.000000,0.000000}%
\pgfsetstrokecolor{currentstroke}%
\pgfsetdash{}{0pt}%
\pgfpathmoveto{\pgfqpoint{0.507620in}{0.352669in}}%
\pgfpathlineto{\pgfqpoint{1.857769in}{0.352669in}}%
\pgfusepath{stroke}%
\end{pgfscope}%
\begin{pgfscope}%
\pgfpathrectangle{\pgfqpoint{0.507620in}{0.352669in}}{\pgfqpoint{1.350149in}{1.771000in}}%
\pgfusepath{clip}%
\pgfsetrectcap%
\pgfsetroundjoin%
\pgfsetlinewidth{0.903375pt}%
\definecolor{currentstroke}{rgb}{0.250980,0.400000,0.600000}%
\pgfsetstrokecolor{currentstroke}%
\pgfsetdash{}{0pt}%
\pgfpathmoveto{\pgfqpoint{0.507620in}{1.131130in}}%
\pgfpathlineto{\pgfqpoint{0.521135in}{1.118586in}}%
\pgfpathlineto{\pgfqpoint{0.534650in}{1.108148in}}%
\pgfpathlineto{\pgfqpoint{0.548165in}{1.099705in}}%
\pgfpathlineto{\pgfqpoint{0.561680in}{1.093144in}}%
\pgfpathlineto{\pgfqpoint{0.575195in}{1.088356in}}%
\pgfpathlineto{\pgfqpoint{0.588710in}{1.085235in}}%
\pgfpathlineto{\pgfqpoint{0.602225in}{1.083677in}}%
\pgfpathlineto{\pgfqpoint{0.615740in}{1.083580in}}%
\pgfpathlineto{\pgfqpoint{0.630606in}{1.085042in}}%
\pgfpathlineto{\pgfqpoint{0.645473in}{1.088022in}}%
\pgfpathlineto{\pgfqpoint{0.661691in}{1.092856in}}%
\pgfpathlineto{\pgfqpoint{0.679260in}{1.099780in}}%
\pgfpathlineto{\pgfqpoint{0.698181in}{1.108978in}}%
\pgfpathlineto{\pgfqpoint{0.718454in}{1.120579in}}%
\pgfpathlineto{\pgfqpoint{0.741429in}{1.135567in}}%
\pgfpathlineto{\pgfqpoint{0.767108in}{1.154176in}}%
\pgfpathlineto{\pgfqpoint{0.798192in}{1.178624in}}%
\pgfpathlineto{\pgfqpoint{0.842792in}{1.215793in}}%
\pgfpathlineto{\pgfqpoint{0.917124in}{1.277760in}}%
\pgfpathlineto{\pgfqpoint{0.952263in}{1.304943in}}%
\pgfpathlineto{\pgfqpoint{0.981996in}{1.326072in}}%
\pgfpathlineto{\pgfqpoint{1.009026in}{1.343431in}}%
\pgfpathlineto{\pgfqpoint{1.033353in}{1.357324in}}%
\pgfpathlineto{\pgfqpoint{1.056329in}{1.368787in}}%
\pgfpathlineto{\pgfqpoint{1.077953in}{1.377996in}}%
\pgfpathlineto{\pgfqpoint{1.099577in}{1.385590in}}%
\pgfpathlineto{\pgfqpoint{1.119849in}{1.391182in}}%
\pgfpathlineto{\pgfqpoint{1.140122in}{1.395253in}}%
\pgfpathlineto{\pgfqpoint{1.160394in}{1.397772in}}%
\pgfpathlineto{\pgfqpoint{1.180667in}{1.398719in}}%
\pgfpathlineto{\pgfqpoint{1.200939in}{1.398087in}}%
\pgfpathlineto{\pgfqpoint{1.221212in}{1.395882in}}%
\pgfpathlineto{\pgfqpoint{1.241484in}{1.392119in}}%
\pgfpathlineto{\pgfqpoint{1.261757in}{1.386828in}}%
\pgfpathlineto{\pgfqpoint{1.282029in}{1.380049in}}%
\pgfpathlineto{\pgfqpoint{1.303653in}{1.371237in}}%
\pgfpathlineto{\pgfqpoint{1.326629in}{1.360171in}}%
\pgfpathlineto{\pgfqpoint{1.349604in}{1.347462in}}%
\pgfpathlineto{\pgfqpoint{1.373931in}{1.332363in}}%
\pgfpathlineto{\pgfqpoint{1.400961in}{1.313822in}}%
\pgfpathlineto{\pgfqpoint{1.430694in}{1.291611in}}%
\pgfpathlineto{\pgfqpoint{1.467185in}{1.262375in}}%
\pgfpathlineto{\pgfqpoint{1.523948in}{1.214649in}}%
\pgfpathlineto{\pgfqpoint{1.576656in}{1.171007in}}%
\pgfpathlineto{\pgfqpoint{1.607741in}{1.147122in}}%
\pgfpathlineto{\pgfqpoint{1.633419in}{1.129182in}}%
\pgfpathlineto{\pgfqpoint{1.656395in}{1.114958in}}%
\pgfpathlineto{\pgfqpoint{1.676667in}{1.104167in}}%
\pgfpathlineto{\pgfqpoint{1.694237in}{1.096378in}}%
\pgfpathlineto{\pgfqpoint{1.711806in}{1.090242in}}%
\pgfpathlineto{\pgfqpoint{1.728024in}{1.086216in}}%
\pgfpathlineto{\pgfqpoint{1.742891in}{1.084048in}}%
\pgfpathlineto{\pgfqpoint{1.757757in}{1.083469in}}%
\pgfpathlineto{\pgfqpoint{1.771272in}{1.084430in}}%
\pgfpathlineto{\pgfqpoint{1.784787in}{1.086914in}}%
\pgfpathlineto{\pgfqpoint{1.798303in}{1.091022in}}%
\pgfpathlineto{\pgfqpoint{1.810466in}{1.096196in}}%
\pgfpathlineto{\pgfqpoint{1.822630in}{1.102850in}}%
\pgfpathlineto{\pgfqpoint{1.834793in}{1.111065in}}%
\pgfpathlineto{\pgfqpoint{1.846957in}{1.120922in}}%
\pgfpathlineto{\pgfqpoint{1.857769in}{1.131130in}}%
\pgfpathlineto{\pgfqpoint{1.857769in}{1.131130in}}%
\pgfusepath{stroke}%
\end{pgfscope}%
\begin{pgfscope}%
\pgfpathrectangle{\pgfqpoint{0.507620in}{0.352669in}}{\pgfqpoint{1.350149in}{1.771000in}}%
\pgfusepath{clip}%
\pgfsetrectcap%
\pgfsetroundjoin%
\pgfsetlinewidth{0.903375pt}%
\definecolor{currentstroke}{rgb}{0.768627,0.556863,0.211765}%
\pgfsetstrokecolor{currentstroke}%
\pgfsetdash{}{0pt}%
\pgfpathmoveto{\pgfqpoint{0.507620in}{1.131130in}}%
\pgfpathlineto{\pgfqpoint{0.514377in}{1.479360in}}%
\pgfpathlineto{\pgfqpoint{0.519783in}{1.690266in}}%
\pgfpathlineto{\pgfqpoint{0.525189in}{1.849550in}}%
\pgfpathlineto{\pgfqpoint{0.530595in}{1.964305in}}%
\pgfpathlineto{\pgfqpoint{0.534650in}{2.025027in}}%
\pgfpathlineto{\pgfqpoint{0.538704in}{2.066783in}}%
\pgfpathlineto{\pgfqpoint{0.542759in}{2.091846in}}%
\pgfpathlineto{\pgfqpoint{0.545462in}{2.100324in}}%
\pgfpathlineto{\pgfqpoint{0.546813in}{2.102313in}}%
\pgfpathlineto{\pgfqpoint{0.548165in}{2.102896in}}%
\pgfpathlineto{\pgfqpoint{0.549516in}{2.102140in}}%
\pgfpathlineto{\pgfqpoint{0.552219in}{2.096869in}}%
\pgfpathlineto{\pgfqpoint{0.554922in}{2.087001in}}%
\pgfpathlineto{\pgfqpoint{0.558977in}{2.064605in}}%
\pgfpathlineto{\pgfqpoint{0.564383in}{2.022825in}}%
\pgfpathlineto{\pgfqpoint{0.571140in}{1.955785in}}%
\pgfpathlineto{\pgfqpoint{0.580601in}{1.843512in}}%
\pgfpathlineto{\pgfqpoint{0.622497in}{1.319672in}}%
\pgfpathlineto{\pgfqpoint{0.631958in}{1.227454in}}%
\pgfpathlineto{\pgfqpoint{0.641418in}{1.150241in}}%
\pgfpathlineto{\pgfqpoint{0.649527in}{1.096344in}}%
\pgfpathlineto{\pgfqpoint{0.656285in}{1.059984in}}%
\pgfpathlineto{\pgfqpoint{0.663042in}{1.031122in}}%
\pgfpathlineto{\pgfqpoint{0.669800in}{1.009370in}}%
\pgfpathlineto{\pgfqpoint{0.675206in}{0.996767in}}%
\pgfpathlineto{\pgfqpoint{0.680612in}{0.988129in}}%
\pgfpathlineto{\pgfqpoint{0.684666in}{0.984073in}}%
\pgfpathlineto{\pgfqpoint{0.688721in}{0.981946in}}%
\pgfpathlineto{\pgfqpoint{0.692775in}{0.981613in}}%
\pgfpathlineto{\pgfqpoint{0.696830in}{0.982936in}}%
\pgfpathlineto{\pgfqpoint{0.700884in}{0.985777in}}%
\pgfpathlineto{\pgfqpoint{0.706290in}{0.991687in}}%
\pgfpathlineto{\pgfqpoint{0.713048in}{1.002027in}}%
\pgfpathlineto{\pgfqpoint{0.721157in}{1.017950in}}%
\pgfpathlineto{\pgfqpoint{0.730617in}{1.040092in}}%
\pgfpathlineto{\pgfqpoint{0.745484in}{1.079240in}}%
\pgfpathlineto{\pgfqpoint{0.771162in}{1.146955in}}%
\pgfpathlineto{\pgfqpoint{0.783326in}{1.174729in}}%
\pgfpathlineto{\pgfqpoint{0.794138in}{1.195683in}}%
\pgfpathlineto{\pgfqpoint{0.803598in}{1.210743in}}%
\pgfpathlineto{\pgfqpoint{0.811707in}{1.221104in}}%
\pgfpathlineto{\pgfqpoint{0.819816in}{1.229109in}}%
\pgfpathlineto{\pgfqpoint{0.827925in}{1.234830in}}%
\pgfpathlineto{\pgfqpoint{0.834683in}{1.237945in}}%
\pgfpathlineto{\pgfqpoint{0.841440in}{1.239669in}}%
\pgfpathlineto{\pgfqpoint{0.849549in}{1.240080in}}%
\pgfpathlineto{\pgfqpoint{0.857658in}{1.238923in}}%
\pgfpathlineto{\pgfqpoint{0.867119in}{1.235967in}}%
\pgfpathlineto{\pgfqpoint{0.879282in}{1.230371in}}%
\pgfpathlineto{\pgfqpoint{0.918476in}{1.210610in}}%
\pgfpathlineto{\pgfqpoint{0.927936in}{1.208029in}}%
\pgfpathlineto{\pgfqpoint{0.936045in}{1.207138in}}%
\pgfpathlineto{\pgfqpoint{0.944154in}{1.207675in}}%
\pgfpathlineto{\pgfqpoint{0.952263in}{1.209803in}}%
\pgfpathlineto{\pgfqpoint{0.960372in}{1.213651in}}%
\pgfpathlineto{\pgfqpoint{0.968481in}{1.219310in}}%
\pgfpathlineto{\pgfqpoint{0.976590in}{1.226833in}}%
\pgfpathlineto{\pgfqpoint{0.984699in}{1.236236in}}%
\pgfpathlineto{\pgfqpoint{0.994160in}{1.249552in}}%
\pgfpathlineto{\pgfqpoint{1.004972in}{1.267744in}}%
\pgfpathlineto{\pgfqpoint{1.015784in}{1.288872in}}%
\pgfpathlineto{\pgfqpoint{1.029299in}{1.318888in}}%
\pgfpathlineto{\pgfqpoint{1.045517in}{1.359056in}}%
\pgfpathlineto{\pgfqpoint{1.069844in}{1.424018in}}%
\pgfpathlineto{\pgfqpoint{1.096874in}{1.495339in}}%
\pgfpathlineto{\pgfqpoint{1.111740in}{1.530812in}}%
\pgfpathlineto{\pgfqpoint{1.123904in}{1.556492in}}%
\pgfpathlineto{\pgfqpoint{1.134716in}{1.576157in}}%
\pgfpathlineto{\pgfqpoint{1.144176in}{1.590558in}}%
\pgfpathlineto{\pgfqpoint{1.152285in}{1.600619in}}%
\pgfpathlineto{\pgfqpoint{1.160394in}{1.608436in}}%
\pgfpathlineto{\pgfqpoint{1.167152in}{1.613166in}}%
\pgfpathlineto{\pgfqpoint{1.173909in}{1.616226in}}%
\pgfpathlineto{\pgfqpoint{1.180667in}{1.617592in}}%
\pgfpathlineto{\pgfqpoint{1.187424in}{1.617250in}}%
\pgfpathlineto{\pgfqpoint{1.194182in}{1.615204in}}%
\pgfpathlineto{\pgfqpoint{1.200939in}{1.611472in}}%
\pgfpathlineto{\pgfqpoint{1.207697in}{1.606086in}}%
\pgfpathlineto{\pgfqpoint{1.215806in}{1.597509in}}%
\pgfpathlineto{\pgfqpoint{1.223915in}{1.586728in}}%
\pgfpathlineto{\pgfqpoint{1.233375in}{1.571547in}}%
\pgfpathlineto{\pgfqpoint{1.244187in}{1.551089in}}%
\pgfpathlineto{\pgfqpoint{1.256351in}{1.524662in}}%
\pgfpathlineto{\pgfqpoint{1.271217in}{1.488511in}}%
\pgfpathlineto{\pgfqpoint{1.292841in}{1.431337in}}%
\pgfpathlineto{\pgfqpoint{1.326629in}{1.341866in}}%
\pgfpathlineto{\pgfqpoint{1.342847in}{1.303422in}}%
\pgfpathlineto{\pgfqpoint{1.356362in}{1.275339in}}%
\pgfpathlineto{\pgfqpoint{1.367174in}{1.256012in}}%
\pgfpathlineto{\pgfqpoint{1.377986in}{1.239785in}}%
\pgfpathlineto{\pgfqpoint{1.387446in}{1.228269in}}%
\pgfpathlineto{\pgfqpoint{1.395555in}{1.220433in}}%
\pgfpathlineto{\pgfqpoint{1.403664in}{1.214466in}}%
\pgfpathlineto{\pgfqpoint{1.411773in}{1.210322in}}%
\pgfpathlineto{\pgfqpoint{1.419882in}{1.207915in}}%
\pgfpathlineto{\pgfqpoint{1.427991in}{1.207123in}}%
\pgfpathlineto{\pgfqpoint{1.436100in}{1.207790in}}%
\pgfpathlineto{\pgfqpoint{1.445561in}{1.210152in}}%
\pgfpathlineto{\pgfqpoint{1.456373in}{1.214498in}}%
\pgfpathlineto{\pgfqpoint{1.471239in}{1.222238in}}%
\pgfpathlineto{\pgfqpoint{1.494215in}{1.234283in}}%
\pgfpathlineto{\pgfqpoint{1.505027in}{1.238237in}}%
\pgfpathlineto{\pgfqpoint{1.513136in}{1.239855in}}%
\pgfpathlineto{\pgfqpoint{1.521245in}{1.239997in}}%
\pgfpathlineto{\pgfqpoint{1.528002in}{1.238795in}}%
\pgfpathlineto{\pgfqpoint{1.534760in}{1.236250in}}%
\pgfpathlineto{\pgfqpoint{1.541517in}{1.232248in}}%
\pgfpathlineto{\pgfqpoint{1.548275in}{1.226700in}}%
\pgfpathlineto{\pgfqpoint{1.556384in}{1.217914in}}%
\pgfpathlineto{\pgfqpoint{1.564493in}{1.206764in}}%
\pgfpathlineto{\pgfqpoint{1.573953in}{1.190807in}}%
\pgfpathlineto{\pgfqpoint{1.583414in}{1.171848in}}%
\pgfpathlineto{\pgfqpoint{1.594226in}{1.146955in}}%
\pgfpathlineto{\pgfqpoint{1.609092in}{1.108585in}}%
\pgfpathlineto{\pgfqpoint{1.638825in}{1.030223in}}%
\pgfpathlineto{\pgfqpoint{1.648286in}{1.009567in}}%
\pgfpathlineto{\pgfqpoint{1.656395in}{0.995459in}}%
\pgfpathlineto{\pgfqpoint{1.661801in}{0.988446in}}%
\pgfpathlineto{\pgfqpoint{1.667207in}{0.983722in}}%
\pgfpathlineto{\pgfqpoint{1.671261in}{0.981877in}}%
\pgfpathlineto{\pgfqpoint{1.675316in}{0.981642in}}%
\pgfpathlineto{\pgfqpoint{1.679370in}{0.983156in}}%
\pgfpathlineto{\pgfqpoint{1.683425in}{0.986554in}}%
\pgfpathlineto{\pgfqpoint{1.687479in}{0.991971in}}%
\pgfpathlineto{\pgfqpoint{1.691534in}{0.999535in}}%
\pgfpathlineto{\pgfqpoint{1.696940in}{1.013174in}}%
\pgfpathlineto{\pgfqpoint{1.702346in}{1.031122in}}%
\pgfpathlineto{\pgfqpoint{1.709103in}{1.059984in}}%
\pgfpathlineto{\pgfqpoint{1.715861in}{1.096344in}}%
\pgfpathlineto{\pgfqpoint{1.722618in}{1.140471in}}%
\pgfpathlineto{\pgfqpoint{1.730727in}{1.203828in}}%
\pgfpathlineto{\pgfqpoint{1.740188in}{1.291853in}}%
\pgfpathlineto{\pgfqpoint{1.749648in}{1.394012in}}%
\pgfpathlineto{\pgfqpoint{1.761812in}{1.542427in}}%
\pgfpathlineto{\pgfqpoint{1.795599in}{1.970275in}}%
\pgfpathlineto{\pgfqpoint{1.802357in}{2.034396in}}%
\pgfpathlineto{\pgfqpoint{1.807763in}{2.073009in}}%
\pgfpathlineto{\pgfqpoint{1.811818in}{2.092480in}}%
\pgfpathlineto{\pgfqpoint{1.814521in}{2.100109in}}%
\pgfpathlineto{\pgfqpoint{1.817224in}{2.102896in}}%
\pgfpathlineto{\pgfqpoint{1.818575in}{2.102313in}}%
\pgfpathlineto{\pgfqpoint{1.819927in}{2.100324in}}%
\pgfpathlineto{\pgfqpoint{1.822630in}{2.091846in}}%
\pgfpathlineto{\pgfqpoint{1.825333in}{2.076885in}}%
\pgfpathlineto{\pgfqpoint{1.828036in}{2.054827in}}%
\pgfpathlineto{\pgfqpoint{1.832090in}{2.007012in}}%
\pgfpathlineto{\pgfqpoint{1.836145in}{1.939430in}}%
\pgfpathlineto{\pgfqpoint{1.840199in}{1.849550in}}%
\pgfpathlineto{\pgfqpoint{1.845605in}{1.690266in}}%
\pgfpathlineto{\pgfqpoint{1.851011in}{1.479360in}}%
\pgfpathlineto{\pgfqpoint{1.856417in}{1.209002in}}%
\pgfpathlineto{\pgfqpoint{1.857769in}{1.131130in}}%
\pgfpathlineto{\pgfqpoint{1.857769in}{1.131130in}}%
\pgfusepath{stroke}%
\end{pgfscope}%
\begin{pgfscope}%
\pgfpathrectangle{\pgfqpoint{0.507620in}{0.352669in}}{\pgfqpoint{1.350149in}{1.771000in}}%
\pgfusepath{clip}%
\pgfsetrectcap%
\pgfsetroundjoin%
\pgfsetlinewidth{0.903375pt}%
\definecolor{currentstroke}{rgb}{0.619608,0.188235,0.172549}%
\pgfsetstrokecolor{currentstroke}%
\pgfsetdash{}{0pt}%
\pgfpathmoveto{\pgfqpoint{0.507620in}{1.131130in}}%
\pgfpathlineto{\pgfqpoint{0.513026in}{1.611858in}}%
\pgfpathlineto{\pgfqpoint{0.517080in}{1.859589in}}%
\pgfpathlineto{\pgfqpoint{0.521135in}{2.028887in}}%
\pgfpathlineto{\pgfqpoint{0.525021in}{2.130335in}}%
\pgfpathmoveto{\pgfqpoint{0.541541in}{2.130335in}}%
\pgfpathlineto{\pgfqpoint{0.546813in}{2.044592in}}%
\pgfpathlineto{\pgfqpoint{0.554922in}{1.879641in}}%
\pgfpathlineto{\pgfqpoint{0.575195in}{1.452385in}}%
\pgfpathlineto{\pgfqpoint{0.583304in}{1.313992in}}%
\pgfpathlineto{\pgfqpoint{0.590061in}{1.219667in}}%
\pgfpathlineto{\pgfqpoint{0.596819in}{1.144933in}}%
\pgfpathlineto{\pgfqpoint{0.602225in}{1.098736in}}%
\pgfpathlineto{\pgfqpoint{0.607631in}{1.063769in}}%
\pgfpathlineto{\pgfqpoint{0.613037in}{1.039004in}}%
\pgfpathlineto{\pgfqpoint{0.617091in}{1.026430in}}%
\pgfpathlineto{\pgfqpoint{0.621146in}{1.018423in}}%
\pgfpathlineto{\pgfqpoint{0.623849in}{1.015362in}}%
\pgfpathlineto{\pgfqpoint{0.626552in}{1.013942in}}%
\pgfpathlineto{\pgfqpoint{0.629255in}{1.014008in}}%
\pgfpathlineto{\pgfqpoint{0.631958in}{1.015411in}}%
\pgfpathlineto{\pgfqpoint{0.636012in}{1.019701in}}%
\pgfpathlineto{\pgfqpoint{0.640067in}{1.026196in}}%
\pgfpathlineto{\pgfqpoint{0.645473in}{1.037519in}}%
\pgfpathlineto{\pgfqpoint{0.653582in}{1.058295in}}%
\pgfpathlineto{\pgfqpoint{0.681963in}{1.134682in}}%
\pgfpathlineto{\pgfqpoint{0.690072in}{1.151326in}}%
\pgfpathlineto{\pgfqpoint{0.696830in}{1.162445in}}%
\pgfpathlineto{\pgfqpoint{0.703587in}{1.170989in}}%
\pgfpathlineto{\pgfqpoint{0.710345in}{1.177034in}}%
\pgfpathlineto{\pgfqpoint{0.717102in}{1.180765in}}%
\pgfpathlineto{\pgfqpoint{0.723860in}{1.182445in}}%
\pgfpathlineto{\pgfqpoint{0.730617in}{1.182393in}}%
\pgfpathlineto{\pgfqpoint{0.738726in}{1.180534in}}%
\pgfpathlineto{\pgfqpoint{0.748187in}{1.176668in}}%
\pgfpathlineto{\pgfqpoint{0.784677in}{1.159901in}}%
\pgfpathlineto{\pgfqpoint{0.794138in}{1.157984in}}%
\pgfpathlineto{\pgfqpoint{0.803598in}{1.157694in}}%
\pgfpathlineto{\pgfqpoint{0.813059in}{1.159072in}}%
\pgfpathlineto{\pgfqpoint{0.822519in}{1.162025in}}%
\pgfpathlineto{\pgfqpoint{0.833331in}{1.167075in}}%
\pgfpathlineto{\pgfqpoint{0.846846in}{1.175303in}}%
\pgfpathlineto{\pgfqpoint{0.867119in}{1.189787in}}%
\pgfpathlineto{\pgfqpoint{0.895500in}{1.209895in}}%
\pgfpathlineto{\pgfqpoint{0.915773in}{1.222249in}}%
\pgfpathlineto{\pgfqpoint{0.967130in}{1.251968in}}%
\pgfpathlineto{\pgfqpoint{0.979293in}{1.261624in}}%
\pgfpathlineto{\pgfqpoint{0.990105in}{1.272013in}}%
\pgfpathlineto{\pgfqpoint{1.000917in}{1.284434in}}%
\pgfpathlineto{\pgfqpoint{1.011729in}{1.299114in}}%
\pgfpathlineto{\pgfqpoint{1.022541in}{1.316165in}}%
\pgfpathlineto{\pgfqpoint{1.034705in}{1.338164in}}%
\pgfpathlineto{\pgfqpoint{1.048220in}{1.365834in}}%
\pgfpathlineto{\pgfqpoint{1.064438in}{1.402667in}}%
\pgfpathlineto{\pgfqpoint{1.094171in}{1.475018in}}%
\pgfpathlineto{\pgfqpoint{1.113092in}{1.519175in}}%
\pgfpathlineto{\pgfqpoint{1.126607in}{1.547258in}}%
\pgfpathlineto{\pgfqpoint{1.137419in}{1.566626in}}%
\pgfpathlineto{\pgfqpoint{1.146879in}{1.580793in}}%
\pgfpathlineto{\pgfqpoint{1.154988in}{1.590596in}}%
\pgfpathlineto{\pgfqpoint{1.163097in}{1.598051in}}%
\pgfpathlineto{\pgfqpoint{1.169855in}{1.602371in}}%
\pgfpathlineto{\pgfqpoint{1.176612in}{1.604912in}}%
\pgfpathlineto{\pgfqpoint{1.183370in}{1.605641in}}%
\pgfpathlineto{\pgfqpoint{1.190127in}{1.604548in}}%
\pgfpathlineto{\pgfqpoint{1.196885in}{1.601648in}}%
\pgfpathlineto{\pgfqpoint{1.203642in}{1.596978in}}%
\pgfpathlineto{\pgfqpoint{1.210400in}{1.590596in}}%
\pgfpathlineto{\pgfqpoint{1.218509in}{1.580793in}}%
\pgfpathlineto{\pgfqpoint{1.227969in}{1.566626in}}%
\pgfpathlineto{\pgfqpoint{1.238781in}{1.547258in}}%
\pgfpathlineto{\pgfqpoint{1.250945in}{1.522144in}}%
\pgfpathlineto{\pgfqpoint{1.267163in}{1.484794in}}%
\pgfpathlineto{\pgfqpoint{1.319871in}{1.360052in}}%
\pgfpathlineto{\pgfqpoint{1.333386in}{1.333025in}}%
\pgfpathlineto{\pgfqpoint{1.345550in}{1.311678in}}%
\pgfpathlineto{\pgfqpoint{1.357713in}{1.293335in}}%
\pgfpathlineto{\pgfqpoint{1.368525in}{1.279522in}}%
\pgfpathlineto{\pgfqpoint{1.379337in}{1.267894in}}%
\pgfpathlineto{\pgfqpoint{1.391501in}{1.257095in}}%
\pgfpathlineto{\pgfqpoint{1.403664in}{1.248232in}}%
\pgfpathlineto{\pgfqpoint{1.419882in}{1.238431in}}%
\pgfpathlineto{\pgfqpoint{1.469888in}{1.209895in}}%
\pgfpathlineto{\pgfqpoint{1.491512in}{1.194761in}}%
\pgfpathlineto{\pgfqpoint{1.523948in}{1.171806in}}%
\pgfpathlineto{\pgfqpoint{1.536111in}{1.164994in}}%
\pgfpathlineto{\pgfqpoint{1.546923in}{1.160577in}}%
\pgfpathlineto{\pgfqpoint{1.556384in}{1.158280in}}%
\pgfpathlineto{\pgfqpoint{1.565844in}{1.157613in}}%
\pgfpathlineto{\pgfqpoint{1.575305in}{1.158614in}}%
\pgfpathlineto{\pgfqpoint{1.584765in}{1.161178in}}%
\pgfpathlineto{\pgfqpoint{1.595577in}{1.165669in}}%
\pgfpathlineto{\pgfqpoint{1.629365in}{1.181337in}}%
\pgfpathlineto{\pgfqpoint{1.637474in}{1.182600in}}%
\pgfpathlineto{\pgfqpoint{1.644231in}{1.181999in}}%
\pgfpathlineto{\pgfqpoint{1.649637in}{1.180192in}}%
\pgfpathlineto{\pgfqpoint{1.655043in}{1.177034in}}%
\pgfpathlineto{\pgfqpoint{1.660449in}{1.172393in}}%
\pgfpathlineto{\pgfqpoint{1.667207in}{1.164359in}}%
\pgfpathlineto{\pgfqpoint{1.673964in}{1.153755in}}%
\pgfpathlineto{\pgfqpoint{1.682073in}{1.137693in}}%
\pgfpathlineto{\pgfqpoint{1.691534in}{1.114861in}}%
\pgfpathlineto{\pgfqpoint{1.705049in}{1.077328in}}%
\pgfpathlineto{\pgfqpoint{1.718564in}{1.040724in}}%
\pgfpathlineto{\pgfqpoint{1.725321in}{1.026196in}}%
\pgfpathlineto{\pgfqpoint{1.730727in}{1.018003in}}%
\pgfpathlineto{\pgfqpoint{1.734782in}{1.014552in}}%
\pgfpathlineto{\pgfqpoint{1.737485in}{1.013799in}}%
\pgfpathlineto{\pgfqpoint{1.740188in}{1.014457in}}%
\pgfpathlineto{\pgfqpoint{1.742891in}{1.016678in}}%
\pgfpathlineto{\pgfqpoint{1.745594in}{1.020617in}}%
\pgfpathlineto{\pgfqpoint{1.749648in}{1.030088in}}%
\pgfpathlineto{\pgfqpoint{1.753703in}{1.044301in}}%
\pgfpathlineto{\pgfqpoint{1.757757in}{1.063769in}}%
\pgfpathlineto{\pgfqpoint{1.763163in}{1.098736in}}%
\pgfpathlineto{\pgfqpoint{1.768569in}{1.144933in}}%
\pgfpathlineto{\pgfqpoint{1.773975in}{1.203164in}}%
\pgfpathlineto{\pgfqpoint{1.780733in}{1.293558in}}%
\pgfpathlineto{\pgfqpoint{1.787490in}{1.403351in}}%
\pgfpathlineto{\pgfqpoint{1.795599in}{1.558169in}}%
\pgfpathlineto{\pgfqpoint{1.810466in}{1.879641in}}%
\pgfpathlineto{\pgfqpoint{1.819927in}{2.068784in}}%
\pgfpathlineto{\pgfqpoint{1.823847in}{2.130335in}}%
\pgfpathmoveto{\pgfqpoint{1.840367in}{2.130335in}}%
\pgfpathlineto{\pgfqpoint{1.842902in}{2.070392in}}%
\pgfpathlineto{\pgfqpoint{1.846957in}{1.923997in}}%
\pgfpathlineto{\pgfqpoint{1.851011in}{1.704107in}}%
\pgfpathlineto{\pgfqpoint{1.855066in}{1.395200in}}%
\pgfpathlineto{\pgfqpoint{1.857769in}{1.131130in}}%
\pgfpathlineto{\pgfqpoint{1.857769in}{1.131130in}}%
\pgfusepath{stroke}%
\end{pgfscope}%
\begin{pgfscope}%
\definecolor{textcolor}{rgb}{0.000000,0.000000,0.000000}%
\pgfsetstrokecolor{textcolor}%
\pgfsetfillcolor{textcolor}%
\pgftext[x=1.182694in,y=2.207002in,,base]{\color{textcolor}{\sffamily\fontsize{8.000000}{9.600000}\selectfont\catcode`\^=\active\def^{\ifmmode\sp\else\^{}\fi}\catcode`\%=\active\def%{\%}(a) Equidistant nodes}}%
\end{pgfscope}%
\begin{pgfscope}%
\pgfpathrectangle{\pgfqpoint{0.507620in}{0.352669in}}{\pgfqpoint{1.350149in}{1.771000in}}%
\pgfusepath{clip}%
\pgfsetbuttcap%
\pgfsetroundjoin%
\definecolor{currentfill}{rgb}{0.250980,0.400000,0.600000}%
\pgfsetfillcolor{currentfill}%
\pgfsetlinewidth{1.003750pt}%
\definecolor{currentstroke}{rgb}{0.250980,0.400000,0.600000}%
\pgfsetstrokecolor{currentstroke}%
\pgfsetdash{}{0pt}%
\pgfsys@defobject{currentmarker}{\pgfqpoint{-0.013889in}{-0.013889in}}{\pgfqpoint{0.013889in}{0.013889in}}{%
\pgfpathmoveto{\pgfqpoint{0.000000in}{-0.013889in}}%
\pgfpathcurveto{\pgfqpoint{0.003683in}{-0.013889in}}{\pgfqpoint{0.007216in}{-0.012425in}}{\pgfqpoint{0.009821in}{-0.009821in}}%
\pgfpathcurveto{\pgfqpoint{0.012425in}{-0.007216in}}{\pgfqpoint{0.013889in}{-0.003683in}}{\pgfqpoint{0.013889in}{0.000000in}}%
\pgfpathcurveto{\pgfqpoint{0.013889in}{0.003683in}}{\pgfqpoint{0.012425in}{0.007216in}}{\pgfqpoint{0.009821in}{0.009821in}}%
\pgfpathcurveto{\pgfqpoint{0.007216in}{0.012425in}}{\pgfqpoint{0.003683in}{0.013889in}}{\pgfqpoint{0.000000in}{0.013889in}}%
\pgfpathcurveto{\pgfqpoint{-0.003683in}{0.013889in}}{\pgfqpoint{-0.007216in}{0.012425in}}{\pgfqpoint{-0.009821in}{0.009821in}}%
\pgfpathcurveto{\pgfqpoint{-0.012425in}{0.007216in}}{\pgfqpoint{-0.013889in}{0.003683in}}{\pgfqpoint{-0.013889in}{0.000000in}}%
\pgfpathcurveto{\pgfqpoint{-0.013889in}{-0.003683in}}{\pgfqpoint{-0.012425in}{-0.007216in}}{\pgfqpoint{-0.009821in}{-0.009821in}}%
\pgfpathcurveto{\pgfqpoint{-0.007216in}{-0.012425in}}{\pgfqpoint{-0.003683in}{-0.013889in}}{\pgfqpoint{0.000000in}{-0.013889in}}%
\pgfpathlineto{\pgfqpoint{0.000000in}{-0.013889in}}%
\pgfpathclose%
\pgfusepath{stroke,fill}%
}%
\begin{pgfscope}%
\pgfsys@transformshift{0.507620in}{1.131130in}%
\pgfsys@useobject{currentmarker}{}%
\end{pgfscope}%
\begin{pgfscope}%
\pgfsys@transformshift{0.777649in}{1.162269in}%
\pgfsys@useobject{currentmarker}{}%
\end{pgfscope}%
\begin{pgfscope}%
\pgfsys@transformshift{1.047679in}{1.364669in}%
\pgfsys@useobject{currentmarker}{}%
\end{pgfscope}%
\begin{pgfscope}%
\pgfsys@transformshift{1.317709in}{1.364669in}%
\pgfsys@useobject{currentmarker}{}%
\end{pgfscope}%
\begin{pgfscope}%
\pgfsys@transformshift{1.587739in}{1.162269in}%
\pgfsys@useobject{currentmarker}{}%
\end{pgfscope}%
\begin{pgfscope}%
\pgfsys@transformshift{1.857769in}{1.131130in}%
\pgfsys@useobject{currentmarker}{}%
\end{pgfscope}%
\end{pgfscope}%
\begin{pgfscope}%
\pgfpathrectangle{\pgfqpoint{0.507620in}{0.352669in}}{\pgfqpoint{1.350149in}{1.771000in}}%
\pgfusepath{clip}%
\pgfsetbuttcap%
\pgfsetroundjoin%
\definecolor{currentfill}{rgb}{0.768627,0.556863,0.211765}%
\pgfsetfillcolor{currentfill}%
\pgfsetlinewidth{1.003750pt}%
\definecolor{currentstroke}{rgb}{0.768627,0.556863,0.211765}%
\pgfsetstrokecolor{currentstroke}%
\pgfsetdash{}{0pt}%
\pgfsys@defobject{currentmarker}{\pgfqpoint{-0.013889in}{-0.013889in}}{\pgfqpoint{0.013889in}{0.013889in}}{%
\pgfpathmoveto{\pgfqpoint{0.000000in}{-0.013889in}}%
\pgfpathcurveto{\pgfqpoint{0.003683in}{-0.013889in}}{\pgfqpoint{0.007216in}{-0.012425in}}{\pgfqpoint{0.009821in}{-0.009821in}}%
\pgfpathcurveto{\pgfqpoint{0.012425in}{-0.007216in}}{\pgfqpoint{0.013889in}{-0.003683in}}{\pgfqpoint{0.013889in}{0.000000in}}%
\pgfpathcurveto{\pgfqpoint{0.013889in}{0.003683in}}{\pgfqpoint{0.012425in}{0.007216in}}{\pgfqpoint{0.009821in}{0.009821in}}%
\pgfpathcurveto{\pgfqpoint{0.007216in}{0.012425in}}{\pgfqpoint{0.003683in}{0.013889in}}{\pgfqpoint{0.000000in}{0.013889in}}%
\pgfpathcurveto{\pgfqpoint{-0.003683in}{0.013889in}}{\pgfqpoint{-0.007216in}{0.012425in}}{\pgfqpoint{-0.009821in}{0.009821in}}%
\pgfpathcurveto{\pgfqpoint{-0.012425in}{0.007216in}}{\pgfqpoint{-0.013889in}{0.003683in}}{\pgfqpoint{-0.013889in}{0.000000in}}%
\pgfpathcurveto{\pgfqpoint{-0.013889in}{-0.003683in}}{\pgfqpoint{-0.012425in}{-0.007216in}}{\pgfqpoint{-0.009821in}{-0.009821in}}%
\pgfpathcurveto{\pgfqpoint{-0.007216in}{-0.012425in}}{\pgfqpoint{-0.003683in}{-0.013889in}}{\pgfqpoint{0.000000in}{-0.013889in}}%
\pgfpathlineto{\pgfqpoint{0.000000in}{-0.013889in}}%
\pgfpathclose%
\pgfusepath{stroke,fill}%
}%
\begin{pgfscope}%
\pgfsys@transformshift{0.507620in}{1.131130in}%
\pgfsys@useobject{currentmarker}{}%
\end{pgfscope}%
\begin{pgfscope}%
\pgfsys@transformshift{0.642634in}{1.141433in}%
\pgfsys@useobject{currentmarker}{}%
\end{pgfscope}%
\begin{pgfscope}%
\pgfsys@transformshift{0.777649in}{1.162269in}%
\pgfsys@useobject{currentmarker}{}%
\end{pgfscope}%
\begin{pgfscope}%
\pgfsys@transformshift{0.912664in}{1.212869in}%
\pgfsys@useobject{currentmarker}{}%
\end{pgfscope}%
\begin{pgfscope}%
\pgfsys@transformshift{1.047679in}{1.364669in}%
\pgfsys@useobject{currentmarker}{}%
\end{pgfscope}%
\begin{pgfscope}%
\pgfsys@transformshift{1.182694in}{1.617669in}%
\pgfsys@useobject{currentmarker}{}%
\end{pgfscope}%
\begin{pgfscope}%
\pgfsys@transformshift{1.317709in}{1.364669in}%
\pgfsys@useobject{currentmarker}{}%
\end{pgfscope}%
\begin{pgfscope}%
\pgfsys@transformshift{1.452724in}{1.212869in}%
\pgfsys@useobject{currentmarker}{}%
\end{pgfscope}%
\begin{pgfscope}%
\pgfsys@transformshift{1.587739in}{1.162269in}%
\pgfsys@useobject{currentmarker}{}%
\end{pgfscope}%
\begin{pgfscope}%
\pgfsys@transformshift{1.722754in}{1.141433in}%
\pgfsys@useobject{currentmarker}{}%
\end{pgfscope}%
\begin{pgfscope}%
\pgfsys@transformshift{1.857769in}{1.131130in}%
\pgfsys@useobject{currentmarker}{}%
\end{pgfscope}%
\end{pgfscope}%
\begin{pgfscope}%
\pgfpathrectangle{\pgfqpoint{0.507620in}{0.352669in}}{\pgfqpoint{1.350149in}{1.771000in}}%
\pgfusepath{clip}%
\pgfsetbuttcap%
\pgfsetroundjoin%
\definecolor{currentfill}{rgb}{0.619608,0.188235,0.172549}%
\pgfsetfillcolor{currentfill}%
\pgfsetlinewidth{1.003750pt}%
\definecolor{currentstroke}{rgb}{0.619608,0.188235,0.172549}%
\pgfsetstrokecolor{currentstroke}%
\pgfsetdash{}{0pt}%
\pgfsys@defobject{currentmarker}{\pgfqpoint{-0.013889in}{-0.013889in}}{\pgfqpoint{0.013889in}{0.013889in}}{%
\pgfpathmoveto{\pgfqpoint{0.000000in}{-0.013889in}}%
\pgfpathcurveto{\pgfqpoint{0.003683in}{-0.013889in}}{\pgfqpoint{0.007216in}{-0.012425in}}{\pgfqpoint{0.009821in}{-0.009821in}}%
\pgfpathcurveto{\pgfqpoint{0.012425in}{-0.007216in}}{\pgfqpoint{0.013889in}{-0.003683in}}{\pgfqpoint{0.013889in}{0.000000in}}%
\pgfpathcurveto{\pgfqpoint{0.013889in}{0.003683in}}{\pgfqpoint{0.012425in}{0.007216in}}{\pgfqpoint{0.009821in}{0.009821in}}%
\pgfpathcurveto{\pgfqpoint{0.007216in}{0.012425in}}{\pgfqpoint{0.003683in}{0.013889in}}{\pgfqpoint{0.000000in}{0.013889in}}%
\pgfpathcurveto{\pgfqpoint{-0.003683in}{0.013889in}}{\pgfqpoint{-0.007216in}{0.012425in}}{\pgfqpoint{-0.009821in}{0.009821in}}%
\pgfpathcurveto{\pgfqpoint{-0.012425in}{0.007216in}}{\pgfqpoint{-0.013889in}{0.003683in}}{\pgfqpoint{-0.013889in}{0.000000in}}%
\pgfpathcurveto{\pgfqpoint{-0.013889in}{-0.003683in}}{\pgfqpoint{-0.012425in}{-0.007216in}}{\pgfqpoint{-0.009821in}{-0.009821in}}%
\pgfpathcurveto{\pgfqpoint{-0.007216in}{-0.012425in}}{\pgfqpoint{-0.003683in}{-0.013889in}}{\pgfqpoint{0.000000in}{-0.013889in}}%
\pgfpathlineto{\pgfqpoint{0.000000in}{-0.013889in}}%
\pgfpathclose%
\pgfusepath{stroke,fill}%
}%
\begin{pgfscope}%
\pgfsys@transformshift{0.507620in}{1.131130in}%
\pgfsys@useobject{currentmarker}{}%
\end{pgfscope}%
\begin{pgfscope}%
\pgfsys@transformshift{0.597630in}{1.137253in}%
\pgfsys@useobject{currentmarker}{}%
\end{pgfscope}%
\begin{pgfscope}%
\pgfsys@transformshift{0.687639in}{1.146700in}%
\pgfsys@useobject{currentmarker}{}%
\end{pgfscope}%
\begin{pgfscope}%
\pgfsys@transformshift{0.777649in}{1.162269in}%
\pgfsys@useobject{currentmarker}{}%
\end{pgfscope}%
\begin{pgfscope}%
\pgfsys@transformshift{0.867659in}{1.190186in}%
\pgfsys@useobject{currentmarker}{}%
\end{pgfscope}%
\begin{pgfscope}%
\pgfsys@transformshift{0.957669in}{1.245610in}%
\pgfsys@useobject{currentmarker}{}%
\end{pgfscope}%
\begin{pgfscope}%
\pgfsys@transformshift{1.047679in}{1.364669in}%
\pgfsys@useobject{currentmarker}{}%
\end{pgfscope}%
\begin{pgfscope}%
\pgfsys@transformshift{1.137689in}{1.567069in}%
\pgfsys@useobject{currentmarker}{}%
\end{pgfscope}%
\begin{pgfscope}%
\pgfsys@transformshift{1.227699in}{1.567069in}%
\pgfsys@useobject{currentmarker}{}%
\end{pgfscope}%
\begin{pgfscope}%
\pgfsys@transformshift{1.317709in}{1.364669in}%
\pgfsys@useobject{currentmarker}{}%
\end{pgfscope}%
\begin{pgfscope}%
\pgfsys@transformshift{1.407719in}{1.245610in}%
\pgfsys@useobject{currentmarker}{}%
\end{pgfscope}%
\begin{pgfscope}%
\pgfsys@transformshift{1.497729in}{1.190186in}%
\pgfsys@useobject{currentmarker}{}%
\end{pgfscope}%
\begin{pgfscope}%
\pgfsys@transformshift{1.587739in}{1.162269in}%
\pgfsys@useobject{currentmarker}{}%
\end{pgfscope}%
\begin{pgfscope}%
\pgfsys@transformshift{1.677749in}{1.146700in}%
\pgfsys@useobject{currentmarker}{}%
\end{pgfscope}%
\begin{pgfscope}%
\pgfsys@transformshift{1.767759in}{1.137253in}%
\pgfsys@useobject{currentmarker}{}%
\end{pgfscope}%
\begin{pgfscope}%
\pgfsys@transformshift{1.857769in}{1.131130in}%
\pgfsys@useobject{currentmarker}{}%
\end{pgfscope}%
\end{pgfscope}%
\begin{pgfscope}%
\pgfpathrectangle{\pgfqpoint{0.507620in}{0.352669in}}{\pgfqpoint{1.350149in}{1.771000in}}%
\pgfusepath{clip}%
\pgfsetrectcap%
\pgfsetroundjoin%
\pgfsetlinewidth{1.405250pt}%
\definecolor{currentstroke}{rgb}{0.000000,0.000000,0.000000}%
\pgfsetstrokecolor{currentstroke}%
\pgfsetdash{}{0pt}%
\pgfpathmoveto{\pgfqpoint{0.507620in}{1.131130in}}%
\pgfpathlineto{\pgfqpoint{0.581952in}{1.135998in}}%
\pgfpathlineto{\pgfqpoint{0.641418in}{1.141308in}}%
\pgfpathlineto{\pgfqpoint{0.691424in}{1.147203in}}%
\pgfpathlineto{\pgfqpoint{0.733320in}{1.153564in}}%
\pgfpathlineto{\pgfqpoint{0.769811in}{1.160550in}}%
\pgfpathlineto{\pgfqpoint{0.802247in}{1.168268in}}%
\pgfpathlineto{\pgfqpoint{0.830628in}{1.176544in}}%
\pgfpathlineto{\pgfqpoint{0.856307in}{1.185603in}}%
\pgfpathlineto{\pgfqpoint{0.879282in}{1.195303in}}%
\pgfpathlineto{\pgfqpoint{0.899555in}{1.205410in}}%
\pgfpathlineto{\pgfqpoint{0.918476in}{1.216438in}}%
\pgfpathlineto{\pgfqpoint{0.936045in}{1.228331in}}%
\pgfpathlineto{\pgfqpoint{0.952263in}{1.240986in}}%
\pgfpathlineto{\pgfqpoint{0.968481in}{1.255531in}}%
\pgfpathlineto{\pgfqpoint{0.983348in}{1.270789in}}%
\pgfpathlineto{\pgfqpoint{0.998214in}{1.288162in}}%
\pgfpathlineto{\pgfqpoint{1.013081in}{1.307931in}}%
\pgfpathlineto{\pgfqpoint{1.027947in}{1.330371in}}%
\pgfpathlineto{\pgfqpoint{1.042814in}{1.355716in}}%
\pgfpathlineto{\pgfqpoint{1.057680in}{1.384102in}}%
\pgfpathlineto{\pgfqpoint{1.073898in}{1.418461in}}%
\pgfpathlineto{\pgfqpoint{1.092819in}{1.462300in}}%
\pgfpathlineto{\pgfqpoint{1.134716in}{1.560936in}}%
\pgfpathlineto{\pgfqpoint{1.145528in}{1.582027in}}%
\pgfpathlineto{\pgfqpoint{1.154988in}{1.597222in}}%
\pgfpathlineto{\pgfqpoint{1.163097in}{1.607229in}}%
\pgfpathlineto{\pgfqpoint{1.169855in}{1.613134in}}%
\pgfpathlineto{\pgfqpoint{1.175261in}{1.616140in}}%
\pgfpathlineto{\pgfqpoint{1.180667in}{1.617555in}}%
\pgfpathlineto{\pgfqpoint{1.186073in}{1.617352in}}%
\pgfpathlineto{\pgfqpoint{1.191479in}{1.615536in}}%
\pgfpathlineto{\pgfqpoint{1.196885in}{1.612140in}}%
\pgfpathlineto{\pgfqpoint{1.203642in}{1.605774in}}%
\pgfpathlineto{\pgfqpoint{1.210400in}{1.597222in}}%
\pgfpathlineto{\pgfqpoint{1.218509in}{1.584404in}}%
\pgfpathlineto{\pgfqpoint{1.227969in}{1.566521in}}%
\pgfpathlineto{\pgfqpoint{1.240133in}{1.540124in}}%
\pgfpathlineto{\pgfqpoint{1.260405in}{1.491752in}}%
\pgfpathlineto{\pgfqpoint{1.288787in}{1.424505in}}%
\pgfpathlineto{\pgfqpoint{1.306356in}{1.386833in}}%
\pgfpathlineto{\pgfqpoint{1.322574in}{1.355716in}}%
\pgfpathlineto{\pgfqpoint{1.337441in}{1.330371in}}%
\pgfpathlineto{\pgfqpoint{1.352307in}{1.307931in}}%
\pgfpathlineto{\pgfqpoint{1.367174in}{1.288162in}}%
\pgfpathlineto{\pgfqpoint{1.382040in}{1.270789in}}%
\pgfpathlineto{\pgfqpoint{1.396907in}{1.255531in}}%
\pgfpathlineto{\pgfqpoint{1.413125in}{1.240986in}}%
\pgfpathlineto{\pgfqpoint{1.430694in}{1.227353in}}%
\pgfpathlineto{\pgfqpoint{1.448264in}{1.215592in}}%
\pgfpathlineto{\pgfqpoint{1.467185in}{1.204685in}}%
\pgfpathlineto{\pgfqpoint{1.487457in}{1.194685in}}%
\pgfpathlineto{\pgfqpoint{1.510433in}{1.185083in}}%
\pgfpathlineto{\pgfqpoint{1.534760in}{1.176544in}}%
\pgfpathlineto{\pgfqpoint{1.561790in}{1.168626in}}%
\pgfpathlineto{\pgfqpoint{1.592874in}{1.161133in}}%
\pgfpathlineto{\pgfqpoint{1.628013in}{1.154266in}}%
\pgfpathlineto{\pgfqpoint{1.668558in}{1.147941in}}%
\pgfpathlineto{\pgfqpoint{1.714509in}{1.142307in}}%
\pgfpathlineto{\pgfqpoint{1.769921in}{1.137074in}}%
\pgfpathlineto{\pgfqpoint{1.836145in}{1.132386in}}%
\pgfpathlineto{\pgfqpoint{1.857769in}{1.131130in}}%
\pgfpathlineto{\pgfqpoint{1.857769in}{1.131130in}}%
\pgfusepath{stroke}%
\end{pgfscope}%
\begin{pgfscope}%
\pgfsetrectcap%
\pgfsetroundjoin%
\pgfsetlinewidth{1.405250pt}%
\definecolor{currentstroke}{rgb}{0.000000,0.000000,0.000000}%
\pgfsetstrokecolor{currentstroke}%
\pgfsetdash{}{0pt}%
\pgfpathmoveto{\pgfqpoint{0.574286in}{2.016419in}}%
\pgfpathlineto{\pgfqpoint{0.657620in}{2.016419in}}%
\pgfpathlineto{\pgfqpoint{0.740953in}{2.016419in}}%
\pgfusepath{stroke}%
\end{pgfscope}%
\begin{pgfscope}%
\definecolor{textcolor}{rgb}{0.000000,0.000000,0.000000}%
\pgfsetstrokecolor{textcolor}%
\pgfsetfillcolor{textcolor}%
\pgftext[x=0.807620in,y=1.987252in,left,base]{\color{textcolor}{\sffamily\fontsize{6.000000}{7.200000}\selectfont\catcode`\^=\active\def^{\ifmmode\sp\else\^{}\fi}\catcode`\%=\active\def%{\%}$f(x)$}}%
\end{pgfscope}%
\begin{pgfscope}%
\pgfsetrectcap%
\pgfsetroundjoin%
\pgfsetlinewidth{0.903375pt}%
\definecolor{currentstroke}{rgb}{0.250980,0.400000,0.600000}%
\pgfsetstrokecolor{currentstroke}%
\pgfsetdash{}{0pt}%
\pgfpathmoveto{\pgfqpoint{0.574286in}{1.892022in}}%
\pgfpathlineto{\pgfqpoint{0.657620in}{1.892022in}}%
\pgfpathlineto{\pgfqpoint{0.740953in}{1.892022in}}%
\pgfusepath{stroke}%
\end{pgfscope}%
\begin{pgfscope}%
\definecolor{textcolor}{rgb}{0.000000,0.000000,0.000000}%
\pgfsetstrokecolor{textcolor}%
\pgfsetfillcolor{textcolor}%
\pgftext[x=0.807620in,y=1.862855in,left,base]{\color{textcolor}{\sffamily\fontsize{6.000000}{7.200000}\selectfont\catcode`\^=\active\def^{\ifmmode\sp\else\^{}\fi}\catcode`\%=\active\def%{\%}$N = 5$}}%
\end{pgfscope}%
\begin{pgfscope}%
\pgfsetrectcap%
\pgfsetroundjoin%
\pgfsetlinewidth{0.903375pt}%
\definecolor{currentstroke}{rgb}{0.768627,0.556863,0.211765}%
\pgfsetstrokecolor{currentstroke}%
\pgfsetdash{}{0pt}%
\pgfpathmoveto{\pgfqpoint{1.214832in}{2.016419in}}%
\pgfpathlineto{\pgfqpoint{1.298165in}{2.016419in}}%
\pgfpathlineto{\pgfqpoint{1.381498in}{2.016419in}}%
\pgfusepath{stroke}%
\end{pgfscope}%
\begin{pgfscope}%
\definecolor{textcolor}{rgb}{0.000000,0.000000,0.000000}%
\pgfsetstrokecolor{textcolor}%
\pgfsetfillcolor{textcolor}%
\pgftext[x=1.448165in,y=1.987252in,left,base]{\color{textcolor}{\sffamily\fontsize{6.000000}{7.200000}\selectfont\catcode`\^=\active\def^{\ifmmode\sp\else\^{}\fi}\catcode`\%=\active\def%{\%}$N = 10$}}%
\end{pgfscope}%
\begin{pgfscope}%
\pgfsetrectcap%
\pgfsetroundjoin%
\pgfsetlinewidth{0.903375pt}%
\definecolor{currentstroke}{rgb}{0.619608,0.188235,0.172549}%
\pgfsetstrokecolor{currentstroke}%
\pgfsetdash{}{0pt}%
\pgfpathmoveto{\pgfqpoint{1.214832in}{1.894104in}}%
\pgfpathlineto{\pgfqpoint{1.298165in}{1.894104in}}%
\pgfpathlineto{\pgfqpoint{1.381498in}{1.894104in}}%
\pgfusepath{stroke}%
\end{pgfscope}%
\begin{pgfscope}%
\definecolor{textcolor}{rgb}{0.000000,0.000000,0.000000}%
\pgfsetstrokecolor{textcolor}%
\pgfsetfillcolor{textcolor}%
\pgftext[x=1.448165in,y=1.864938in,left,base]{\color{textcolor}{\sffamily\fontsize{6.000000}{7.200000}\selectfont\catcode`\^=\active\def^{\ifmmode\sp\else\^{}\fi}\catcode`\%=\active\def%{\%}$N = 15$}}%
\end{pgfscope}%
\begin{pgfscope}%
\pgfsetbuttcap%
\pgfsetmiterjoin%
\definecolor{currentfill}{rgb}{1.000000,1.000000,1.000000}%
\pgfsetfillcolor{currentfill}%
\pgfsetlinewidth{0.000000pt}%
\definecolor{currentstroke}{rgb}{0.000000,0.000000,0.000000}%
\pgfsetstrokecolor{currentstroke}%
\pgfsetstrokeopacity{0.000000}%
\pgfsetdash{}{0pt}%
\pgfpathmoveto{\pgfqpoint{2.330321in}{0.352669in}}%
\pgfpathlineto{\pgfqpoint{3.680470in}{0.352669in}}%
\pgfpathlineto{\pgfqpoint{3.680470in}{2.123669in}}%
\pgfpathlineto{\pgfqpoint{2.330321in}{2.123669in}}%
\pgfpathlineto{\pgfqpoint{2.330321in}{0.352669in}}%
\pgfpathclose%
\pgfusepath{fill}%
\end{pgfscope}%
\begin{pgfscope}%
\pgfsetbuttcap%
\pgfsetroundjoin%
\definecolor{currentfill}{rgb}{0.000000,0.000000,0.000000}%
\pgfsetfillcolor{currentfill}%
\pgfsetlinewidth{0.501875pt}%
\definecolor{currentstroke}{rgb}{0.000000,0.000000,0.000000}%
\pgfsetstrokecolor{currentstroke}%
\pgfsetdash{}{0pt}%
\pgfsys@defobject{currentmarker}{\pgfqpoint{0.000000in}{0.000000in}}{\pgfqpoint{0.000000in}{0.041667in}}{%
\pgfpathmoveto{\pgfqpoint{0.000000in}{0.000000in}}%
\pgfpathlineto{\pgfqpoint{0.000000in}{0.041667in}}%
\pgfusepath{stroke,fill}%
}%
\begin{pgfscope}%
\pgfsys@transformshift{2.330321in}{0.352669in}%
\pgfsys@useobject{currentmarker}{}%
\end{pgfscope}%
\end{pgfscope}%
\begin{pgfscope}%
\definecolor{textcolor}{rgb}{0.000000,0.000000,0.000000}%
\pgfsetstrokecolor{textcolor}%
\pgfsetfillcolor{textcolor}%
\pgftext[x=2.330321in,y=0.304058in,,top]{\color{textcolor}{\sffamily\fontsize{8.000000}{9.600000}\selectfont\catcode`\^=\active\def^{\ifmmode\sp\else\^{}\fi}\catcode`\%=\active\def%{\%}\ensuremath{-}1.0}}%
\end{pgfscope}%
\begin{pgfscope}%
\pgfsetbuttcap%
\pgfsetroundjoin%
\definecolor{currentfill}{rgb}{0.000000,0.000000,0.000000}%
\pgfsetfillcolor{currentfill}%
\pgfsetlinewidth{0.501875pt}%
\definecolor{currentstroke}{rgb}{0.000000,0.000000,0.000000}%
\pgfsetstrokecolor{currentstroke}%
\pgfsetdash{}{0pt}%
\pgfsys@defobject{currentmarker}{\pgfqpoint{0.000000in}{0.000000in}}{\pgfqpoint{0.000000in}{0.041667in}}{%
\pgfpathmoveto{\pgfqpoint{0.000000in}{0.000000in}}%
\pgfpathlineto{\pgfqpoint{0.000000in}{0.041667in}}%
\pgfusepath{stroke,fill}%
}%
\begin{pgfscope}%
\pgfsys@transformshift{2.667858in}{0.352669in}%
\pgfsys@useobject{currentmarker}{}%
\end{pgfscope}%
\end{pgfscope}%
\begin{pgfscope}%
\definecolor{textcolor}{rgb}{0.000000,0.000000,0.000000}%
\pgfsetstrokecolor{textcolor}%
\pgfsetfillcolor{textcolor}%
\pgftext[x=2.667858in,y=0.304058in,,top]{\color{textcolor}{\sffamily\fontsize{8.000000}{9.600000}\selectfont\catcode`\^=\active\def^{\ifmmode\sp\else\^{}\fi}\catcode`\%=\active\def%{\%}\ensuremath{-}0.5}}%
\end{pgfscope}%
\begin{pgfscope}%
\pgfsetbuttcap%
\pgfsetroundjoin%
\definecolor{currentfill}{rgb}{0.000000,0.000000,0.000000}%
\pgfsetfillcolor{currentfill}%
\pgfsetlinewidth{0.501875pt}%
\definecolor{currentstroke}{rgb}{0.000000,0.000000,0.000000}%
\pgfsetstrokecolor{currentstroke}%
\pgfsetdash{}{0pt}%
\pgfsys@defobject{currentmarker}{\pgfqpoint{0.000000in}{0.000000in}}{\pgfqpoint{0.000000in}{0.041667in}}{%
\pgfpathmoveto{\pgfqpoint{0.000000in}{0.000000in}}%
\pgfpathlineto{\pgfqpoint{0.000000in}{0.041667in}}%
\pgfusepath{stroke,fill}%
}%
\begin{pgfscope}%
\pgfsys@transformshift{3.005395in}{0.352669in}%
\pgfsys@useobject{currentmarker}{}%
\end{pgfscope}%
\end{pgfscope}%
\begin{pgfscope}%
\definecolor{textcolor}{rgb}{0.000000,0.000000,0.000000}%
\pgfsetstrokecolor{textcolor}%
\pgfsetfillcolor{textcolor}%
\pgftext[x=3.005395in,y=0.304058in,,top]{\color{textcolor}{\sffamily\fontsize{8.000000}{9.600000}\selectfont\catcode`\^=\active\def^{\ifmmode\sp\else\^{}\fi}\catcode`\%=\active\def%{\%}0.0}}%
\end{pgfscope}%
\begin{pgfscope}%
\pgfsetbuttcap%
\pgfsetroundjoin%
\definecolor{currentfill}{rgb}{0.000000,0.000000,0.000000}%
\pgfsetfillcolor{currentfill}%
\pgfsetlinewidth{0.501875pt}%
\definecolor{currentstroke}{rgb}{0.000000,0.000000,0.000000}%
\pgfsetstrokecolor{currentstroke}%
\pgfsetdash{}{0pt}%
\pgfsys@defobject{currentmarker}{\pgfqpoint{0.000000in}{0.000000in}}{\pgfqpoint{0.000000in}{0.041667in}}{%
\pgfpathmoveto{\pgfqpoint{0.000000in}{0.000000in}}%
\pgfpathlineto{\pgfqpoint{0.000000in}{0.041667in}}%
\pgfusepath{stroke,fill}%
}%
\begin{pgfscope}%
\pgfsys@transformshift{3.342932in}{0.352669in}%
\pgfsys@useobject{currentmarker}{}%
\end{pgfscope}%
\end{pgfscope}%
\begin{pgfscope}%
\definecolor{textcolor}{rgb}{0.000000,0.000000,0.000000}%
\pgfsetstrokecolor{textcolor}%
\pgfsetfillcolor{textcolor}%
\pgftext[x=3.342932in,y=0.304058in,,top]{\color{textcolor}{\sffamily\fontsize{8.000000}{9.600000}\selectfont\catcode`\^=\active\def^{\ifmmode\sp\else\^{}\fi}\catcode`\%=\active\def%{\%}0.5}}%
\end{pgfscope}%
\begin{pgfscope}%
\pgfsetbuttcap%
\pgfsetroundjoin%
\definecolor{currentfill}{rgb}{0.000000,0.000000,0.000000}%
\pgfsetfillcolor{currentfill}%
\pgfsetlinewidth{0.501875pt}%
\definecolor{currentstroke}{rgb}{0.000000,0.000000,0.000000}%
\pgfsetstrokecolor{currentstroke}%
\pgfsetdash{}{0pt}%
\pgfsys@defobject{currentmarker}{\pgfqpoint{0.000000in}{0.000000in}}{\pgfqpoint{0.000000in}{0.041667in}}{%
\pgfpathmoveto{\pgfqpoint{0.000000in}{0.000000in}}%
\pgfpathlineto{\pgfqpoint{0.000000in}{0.041667in}}%
\pgfusepath{stroke,fill}%
}%
\begin{pgfscope}%
\pgfsys@transformshift{3.680470in}{0.352669in}%
\pgfsys@useobject{currentmarker}{}%
\end{pgfscope}%
\end{pgfscope}%
\begin{pgfscope}%
\definecolor{textcolor}{rgb}{0.000000,0.000000,0.000000}%
\pgfsetstrokecolor{textcolor}%
\pgfsetfillcolor{textcolor}%
\pgftext[x=3.680470in,y=0.304058in,,top]{\color{textcolor}{\sffamily\fontsize{8.000000}{9.600000}\selectfont\catcode`\^=\active\def^{\ifmmode\sp\else\^{}\fi}\catcode`\%=\active\def%{\%}1.0}}%
\end{pgfscope}%
\begin{pgfscope}%
\pgfsetbuttcap%
\pgfsetroundjoin%
\definecolor{currentfill}{rgb}{0.000000,0.000000,0.000000}%
\pgfsetfillcolor{currentfill}%
\pgfsetlinewidth{0.301125pt}%
\definecolor{currentstroke}{rgb}{0.000000,0.000000,0.000000}%
\pgfsetstrokecolor{currentstroke}%
\pgfsetdash{}{0pt}%
\pgfsys@defobject{currentmarker}{\pgfqpoint{0.000000in}{0.000000in}}{\pgfqpoint{0.000000in}{0.027778in}}{%
\pgfpathmoveto{\pgfqpoint{0.000000in}{0.000000in}}%
\pgfpathlineto{\pgfqpoint{0.000000in}{0.027778in}}%
\pgfusepath{stroke,fill}%
}%
\begin{pgfscope}%
\pgfsys@transformshift{2.397828in}{0.352669in}%
\pgfsys@useobject{currentmarker}{}%
\end{pgfscope}%
\end{pgfscope}%
\begin{pgfscope}%
\pgfsetbuttcap%
\pgfsetroundjoin%
\definecolor{currentfill}{rgb}{0.000000,0.000000,0.000000}%
\pgfsetfillcolor{currentfill}%
\pgfsetlinewidth{0.301125pt}%
\definecolor{currentstroke}{rgb}{0.000000,0.000000,0.000000}%
\pgfsetstrokecolor{currentstroke}%
\pgfsetdash{}{0pt}%
\pgfsys@defobject{currentmarker}{\pgfqpoint{0.000000in}{0.000000in}}{\pgfqpoint{0.000000in}{0.027778in}}{%
\pgfpathmoveto{\pgfqpoint{0.000000in}{0.000000in}}%
\pgfpathlineto{\pgfqpoint{0.000000in}{0.027778in}}%
\pgfusepath{stroke,fill}%
}%
\begin{pgfscope}%
\pgfsys@transformshift{2.465336in}{0.352669in}%
\pgfsys@useobject{currentmarker}{}%
\end{pgfscope}%
\end{pgfscope}%
\begin{pgfscope}%
\pgfsetbuttcap%
\pgfsetroundjoin%
\definecolor{currentfill}{rgb}{0.000000,0.000000,0.000000}%
\pgfsetfillcolor{currentfill}%
\pgfsetlinewidth{0.301125pt}%
\definecolor{currentstroke}{rgb}{0.000000,0.000000,0.000000}%
\pgfsetstrokecolor{currentstroke}%
\pgfsetdash{}{0pt}%
\pgfsys@defobject{currentmarker}{\pgfqpoint{0.000000in}{0.000000in}}{\pgfqpoint{0.000000in}{0.027778in}}{%
\pgfpathmoveto{\pgfqpoint{0.000000in}{0.000000in}}%
\pgfpathlineto{\pgfqpoint{0.000000in}{0.027778in}}%
\pgfusepath{stroke,fill}%
}%
\begin{pgfscope}%
\pgfsys@transformshift{2.532843in}{0.352669in}%
\pgfsys@useobject{currentmarker}{}%
\end{pgfscope}%
\end{pgfscope}%
\begin{pgfscope}%
\pgfsetbuttcap%
\pgfsetroundjoin%
\definecolor{currentfill}{rgb}{0.000000,0.000000,0.000000}%
\pgfsetfillcolor{currentfill}%
\pgfsetlinewidth{0.301125pt}%
\definecolor{currentstroke}{rgb}{0.000000,0.000000,0.000000}%
\pgfsetstrokecolor{currentstroke}%
\pgfsetdash{}{0pt}%
\pgfsys@defobject{currentmarker}{\pgfqpoint{0.000000in}{0.000000in}}{\pgfqpoint{0.000000in}{0.027778in}}{%
\pgfpathmoveto{\pgfqpoint{0.000000in}{0.000000in}}%
\pgfpathlineto{\pgfqpoint{0.000000in}{0.027778in}}%
\pgfusepath{stroke,fill}%
}%
\begin{pgfscope}%
\pgfsys@transformshift{2.600350in}{0.352669in}%
\pgfsys@useobject{currentmarker}{}%
\end{pgfscope}%
\end{pgfscope}%
\begin{pgfscope}%
\pgfsetbuttcap%
\pgfsetroundjoin%
\definecolor{currentfill}{rgb}{0.000000,0.000000,0.000000}%
\pgfsetfillcolor{currentfill}%
\pgfsetlinewidth{0.301125pt}%
\definecolor{currentstroke}{rgb}{0.000000,0.000000,0.000000}%
\pgfsetstrokecolor{currentstroke}%
\pgfsetdash{}{0pt}%
\pgfsys@defobject{currentmarker}{\pgfqpoint{0.000000in}{0.000000in}}{\pgfqpoint{0.000000in}{0.027778in}}{%
\pgfpathmoveto{\pgfqpoint{0.000000in}{0.000000in}}%
\pgfpathlineto{\pgfqpoint{0.000000in}{0.027778in}}%
\pgfusepath{stroke,fill}%
}%
\begin{pgfscope}%
\pgfsys@transformshift{2.735365in}{0.352669in}%
\pgfsys@useobject{currentmarker}{}%
\end{pgfscope}%
\end{pgfscope}%
\begin{pgfscope}%
\pgfsetbuttcap%
\pgfsetroundjoin%
\definecolor{currentfill}{rgb}{0.000000,0.000000,0.000000}%
\pgfsetfillcolor{currentfill}%
\pgfsetlinewidth{0.301125pt}%
\definecolor{currentstroke}{rgb}{0.000000,0.000000,0.000000}%
\pgfsetstrokecolor{currentstroke}%
\pgfsetdash{}{0pt}%
\pgfsys@defobject{currentmarker}{\pgfqpoint{0.000000in}{0.000000in}}{\pgfqpoint{0.000000in}{0.027778in}}{%
\pgfpathmoveto{\pgfqpoint{0.000000in}{0.000000in}}%
\pgfpathlineto{\pgfqpoint{0.000000in}{0.027778in}}%
\pgfusepath{stroke,fill}%
}%
\begin{pgfscope}%
\pgfsys@transformshift{2.802873in}{0.352669in}%
\pgfsys@useobject{currentmarker}{}%
\end{pgfscope}%
\end{pgfscope}%
\begin{pgfscope}%
\pgfsetbuttcap%
\pgfsetroundjoin%
\definecolor{currentfill}{rgb}{0.000000,0.000000,0.000000}%
\pgfsetfillcolor{currentfill}%
\pgfsetlinewidth{0.301125pt}%
\definecolor{currentstroke}{rgb}{0.000000,0.000000,0.000000}%
\pgfsetstrokecolor{currentstroke}%
\pgfsetdash{}{0pt}%
\pgfsys@defobject{currentmarker}{\pgfqpoint{0.000000in}{0.000000in}}{\pgfqpoint{0.000000in}{0.027778in}}{%
\pgfpathmoveto{\pgfqpoint{0.000000in}{0.000000in}}%
\pgfpathlineto{\pgfqpoint{0.000000in}{0.027778in}}%
\pgfusepath{stroke,fill}%
}%
\begin{pgfscope}%
\pgfsys@transformshift{2.870380in}{0.352669in}%
\pgfsys@useobject{currentmarker}{}%
\end{pgfscope}%
\end{pgfscope}%
\begin{pgfscope}%
\pgfsetbuttcap%
\pgfsetroundjoin%
\definecolor{currentfill}{rgb}{0.000000,0.000000,0.000000}%
\pgfsetfillcolor{currentfill}%
\pgfsetlinewidth{0.301125pt}%
\definecolor{currentstroke}{rgb}{0.000000,0.000000,0.000000}%
\pgfsetstrokecolor{currentstroke}%
\pgfsetdash{}{0pt}%
\pgfsys@defobject{currentmarker}{\pgfqpoint{0.000000in}{0.000000in}}{\pgfqpoint{0.000000in}{0.027778in}}{%
\pgfpathmoveto{\pgfqpoint{0.000000in}{0.000000in}}%
\pgfpathlineto{\pgfqpoint{0.000000in}{0.027778in}}%
\pgfusepath{stroke,fill}%
}%
\begin{pgfscope}%
\pgfsys@transformshift{2.937888in}{0.352669in}%
\pgfsys@useobject{currentmarker}{}%
\end{pgfscope}%
\end{pgfscope}%
\begin{pgfscope}%
\pgfsetbuttcap%
\pgfsetroundjoin%
\definecolor{currentfill}{rgb}{0.000000,0.000000,0.000000}%
\pgfsetfillcolor{currentfill}%
\pgfsetlinewidth{0.301125pt}%
\definecolor{currentstroke}{rgb}{0.000000,0.000000,0.000000}%
\pgfsetstrokecolor{currentstroke}%
\pgfsetdash{}{0pt}%
\pgfsys@defobject{currentmarker}{\pgfqpoint{0.000000in}{0.000000in}}{\pgfqpoint{0.000000in}{0.027778in}}{%
\pgfpathmoveto{\pgfqpoint{0.000000in}{0.000000in}}%
\pgfpathlineto{\pgfqpoint{0.000000in}{0.027778in}}%
\pgfusepath{stroke,fill}%
}%
\begin{pgfscope}%
\pgfsys@transformshift{3.072903in}{0.352669in}%
\pgfsys@useobject{currentmarker}{}%
\end{pgfscope}%
\end{pgfscope}%
\begin{pgfscope}%
\pgfsetbuttcap%
\pgfsetroundjoin%
\definecolor{currentfill}{rgb}{0.000000,0.000000,0.000000}%
\pgfsetfillcolor{currentfill}%
\pgfsetlinewidth{0.301125pt}%
\definecolor{currentstroke}{rgb}{0.000000,0.000000,0.000000}%
\pgfsetstrokecolor{currentstroke}%
\pgfsetdash{}{0pt}%
\pgfsys@defobject{currentmarker}{\pgfqpoint{0.000000in}{0.000000in}}{\pgfqpoint{0.000000in}{0.027778in}}{%
\pgfpathmoveto{\pgfqpoint{0.000000in}{0.000000in}}%
\pgfpathlineto{\pgfqpoint{0.000000in}{0.027778in}}%
\pgfusepath{stroke,fill}%
}%
\begin{pgfscope}%
\pgfsys@transformshift{3.140410in}{0.352669in}%
\pgfsys@useobject{currentmarker}{}%
\end{pgfscope}%
\end{pgfscope}%
\begin{pgfscope}%
\pgfsetbuttcap%
\pgfsetroundjoin%
\definecolor{currentfill}{rgb}{0.000000,0.000000,0.000000}%
\pgfsetfillcolor{currentfill}%
\pgfsetlinewidth{0.301125pt}%
\definecolor{currentstroke}{rgb}{0.000000,0.000000,0.000000}%
\pgfsetstrokecolor{currentstroke}%
\pgfsetdash{}{0pt}%
\pgfsys@defobject{currentmarker}{\pgfqpoint{0.000000in}{0.000000in}}{\pgfqpoint{0.000000in}{0.027778in}}{%
\pgfpathmoveto{\pgfqpoint{0.000000in}{0.000000in}}%
\pgfpathlineto{\pgfqpoint{0.000000in}{0.027778in}}%
\pgfusepath{stroke,fill}%
}%
\begin{pgfscope}%
\pgfsys@transformshift{3.207917in}{0.352669in}%
\pgfsys@useobject{currentmarker}{}%
\end{pgfscope}%
\end{pgfscope}%
\begin{pgfscope}%
\pgfsetbuttcap%
\pgfsetroundjoin%
\definecolor{currentfill}{rgb}{0.000000,0.000000,0.000000}%
\pgfsetfillcolor{currentfill}%
\pgfsetlinewidth{0.301125pt}%
\definecolor{currentstroke}{rgb}{0.000000,0.000000,0.000000}%
\pgfsetstrokecolor{currentstroke}%
\pgfsetdash{}{0pt}%
\pgfsys@defobject{currentmarker}{\pgfqpoint{0.000000in}{0.000000in}}{\pgfqpoint{0.000000in}{0.027778in}}{%
\pgfpathmoveto{\pgfqpoint{0.000000in}{0.000000in}}%
\pgfpathlineto{\pgfqpoint{0.000000in}{0.027778in}}%
\pgfusepath{stroke,fill}%
}%
\begin{pgfscope}%
\pgfsys@transformshift{3.275425in}{0.352669in}%
\pgfsys@useobject{currentmarker}{}%
\end{pgfscope}%
\end{pgfscope}%
\begin{pgfscope}%
\pgfsetbuttcap%
\pgfsetroundjoin%
\definecolor{currentfill}{rgb}{0.000000,0.000000,0.000000}%
\pgfsetfillcolor{currentfill}%
\pgfsetlinewidth{0.301125pt}%
\definecolor{currentstroke}{rgb}{0.000000,0.000000,0.000000}%
\pgfsetstrokecolor{currentstroke}%
\pgfsetdash{}{0pt}%
\pgfsys@defobject{currentmarker}{\pgfqpoint{0.000000in}{0.000000in}}{\pgfqpoint{0.000000in}{0.027778in}}{%
\pgfpathmoveto{\pgfqpoint{0.000000in}{0.000000in}}%
\pgfpathlineto{\pgfqpoint{0.000000in}{0.027778in}}%
\pgfusepath{stroke,fill}%
}%
\begin{pgfscope}%
\pgfsys@transformshift{3.410440in}{0.352669in}%
\pgfsys@useobject{currentmarker}{}%
\end{pgfscope}%
\end{pgfscope}%
\begin{pgfscope}%
\pgfsetbuttcap%
\pgfsetroundjoin%
\definecolor{currentfill}{rgb}{0.000000,0.000000,0.000000}%
\pgfsetfillcolor{currentfill}%
\pgfsetlinewidth{0.301125pt}%
\definecolor{currentstroke}{rgb}{0.000000,0.000000,0.000000}%
\pgfsetstrokecolor{currentstroke}%
\pgfsetdash{}{0pt}%
\pgfsys@defobject{currentmarker}{\pgfqpoint{0.000000in}{0.000000in}}{\pgfqpoint{0.000000in}{0.027778in}}{%
\pgfpathmoveto{\pgfqpoint{0.000000in}{0.000000in}}%
\pgfpathlineto{\pgfqpoint{0.000000in}{0.027778in}}%
\pgfusepath{stroke,fill}%
}%
\begin{pgfscope}%
\pgfsys@transformshift{3.477947in}{0.352669in}%
\pgfsys@useobject{currentmarker}{}%
\end{pgfscope}%
\end{pgfscope}%
\begin{pgfscope}%
\pgfsetbuttcap%
\pgfsetroundjoin%
\definecolor{currentfill}{rgb}{0.000000,0.000000,0.000000}%
\pgfsetfillcolor{currentfill}%
\pgfsetlinewidth{0.301125pt}%
\definecolor{currentstroke}{rgb}{0.000000,0.000000,0.000000}%
\pgfsetstrokecolor{currentstroke}%
\pgfsetdash{}{0pt}%
\pgfsys@defobject{currentmarker}{\pgfqpoint{0.000000in}{0.000000in}}{\pgfqpoint{0.000000in}{0.027778in}}{%
\pgfpathmoveto{\pgfqpoint{0.000000in}{0.000000in}}%
\pgfpathlineto{\pgfqpoint{0.000000in}{0.027778in}}%
\pgfusepath{stroke,fill}%
}%
\begin{pgfscope}%
\pgfsys@transformshift{3.545455in}{0.352669in}%
\pgfsys@useobject{currentmarker}{}%
\end{pgfscope}%
\end{pgfscope}%
\begin{pgfscope}%
\pgfsetbuttcap%
\pgfsetroundjoin%
\definecolor{currentfill}{rgb}{0.000000,0.000000,0.000000}%
\pgfsetfillcolor{currentfill}%
\pgfsetlinewidth{0.301125pt}%
\definecolor{currentstroke}{rgb}{0.000000,0.000000,0.000000}%
\pgfsetstrokecolor{currentstroke}%
\pgfsetdash{}{0pt}%
\pgfsys@defobject{currentmarker}{\pgfqpoint{0.000000in}{0.000000in}}{\pgfqpoint{0.000000in}{0.027778in}}{%
\pgfpathmoveto{\pgfqpoint{0.000000in}{0.000000in}}%
\pgfpathlineto{\pgfqpoint{0.000000in}{0.027778in}}%
\pgfusepath{stroke,fill}%
}%
\begin{pgfscope}%
\pgfsys@transformshift{3.612962in}{0.352669in}%
\pgfsys@useobject{currentmarker}{}%
\end{pgfscope}%
\end{pgfscope}%
\begin{pgfscope}%
\definecolor{textcolor}{rgb}{0.000000,0.000000,0.000000}%
\pgfsetstrokecolor{textcolor}%
\pgfsetfillcolor{textcolor}%
\pgftext[x=3.005395in,y=0.140972in,,top]{\color{textcolor}{\sffamily\fontsize{9.000000}{10.800000}\selectfont\catcode`\^=\active\def^{\ifmmode\sp\else\^{}\fi}\catcode`\%=\active\def%{\%}$x$}}%
\end{pgfscope}%
\begin{pgfscope}%
\pgfsetbuttcap%
\pgfsetroundjoin%
\definecolor{currentfill}{rgb}{0.000000,0.000000,0.000000}%
\pgfsetfillcolor{currentfill}%
\pgfsetlinewidth{0.501875pt}%
\definecolor{currentstroke}{rgb}{0.000000,0.000000,0.000000}%
\pgfsetstrokecolor{currentstroke}%
\pgfsetdash{}{0pt}%
\pgfsys@defobject{currentmarker}{\pgfqpoint{0.000000in}{0.000000in}}{\pgfqpoint{0.041667in}{0.000000in}}{%
\pgfpathmoveto{\pgfqpoint{0.000000in}{0.000000in}}%
\pgfpathlineto{\pgfqpoint{0.041667in}{0.000000in}}%
\pgfusepath{stroke,fill}%
}%
\begin{pgfscope}%
\pgfsys@transformshift{2.330321in}{0.352669in}%
\pgfsys@useobject{currentmarker}{}%
\end{pgfscope}%
\end{pgfscope}%
\begin{pgfscope}%
\definecolor{textcolor}{rgb}{0.000000,0.000000,0.000000}%
\pgfsetstrokecolor{textcolor}%
\pgfsetfillcolor{textcolor}%
\pgftext[x=2.019228in, y=0.310460in, left, base]{\color{textcolor}{\sffamily\fontsize{8.000000}{9.600000}\selectfont\catcode`\^=\active\def^{\ifmmode\sp\else\^{}\fi}\catcode`\%=\active\def%{\%}\ensuremath{-}0.2}}%
\end{pgfscope}%
\begin{pgfscope}%
\pgfsetbuttcap%
\pgfsetroundjoin%
\definecolor{currentfill}{rgb}{0.000000,0.000000,0.000000}%
\pgfsetfillcolor{currentfill}%
\pgfsetlinewidth{0.501875pt}%
\definecolor{currentstroke}{rgb}{0.000000,0.000000,0.000000}%
\pgfsetstrokecolor{currentstroke}%
\pgfsetdash{}{0pt}%
\pgfsys@defobject{currentmarker}{\pgfqpoint{0.000000in}{0.000000in}}{\pgfqpoint{0.041667in}{0.000000in}}{%
\pgfpathmoveto{\pgfqpoint{0.000000in}{0.000000in}}%
\pgfpathlineto{\pgfqpoint{0.041667in}{0.000000in}}%
\pgfusepath{stroke,fill}%
}%
\begin{pgfscope}%
\pgfsys@transformshift{2.330321in}{0.615039in}%
\pgfsys@useobject{currentmarker}{}%
\end{pgfscope}%
\end{pgfscope}%
\begin{pgfscope}%
\definecolor{textcolor}{rgb}{0.000000,0.000000,0.000000}%
\pgfsetstrokecolor{textcolor}%
\pgfsetfillcolor{textcolor}%
\pgftext[x=2.105006in, y=0.572830in, left, base]{\color{textcolor}{\sffamily\fontsize{8.000000}{9.600000}\selectfont\catcode`\^=\active\def^{\ifmmode\sp\else\^{}\fi}\catcode`\%=\active\def%{\%}0.0}}%
\end{pgfscope}%
\begin{pgfscope}%
\pgfsetbuttcap%
\pgfsetroundjoin%
\definecolor{currentfill}{rgb}{0.000000,0.000000,0.000000}%
\pgfsetfillcolor{currentfill}%
\pgfsetlinewidth{0.501875pt}%
\definecolor{currentstroke}{rgb}{0.000000,0.000000,0.000000}%
\pgfsetstrokecolor{currentstroke}%
\pgfsetdash{}{0pt}%
\pgfsys@defobject{currentmarker}{\pgfqpoint{0.000000in}{0.000000in}}{\pgfqpoint{0.041667in}{0.000000in}}{%
\pgfpathmoveto{\pgfqpoint{0.000000in}{0.000000in}}%
\pgfpathlineto{\pgfqpoint{0.041667in}{0.000000in}}%
\pgfusepath{stroke,fill}%
}%
\begin{pgfscope}%
\pgfsys@transformshift{2.330321in}{0.877409in}%
\pgfsys@useobject{currentmarker}{}%
\end{pgfscope}%
\end{pgfscope}%
\begin{pgfscope}%
\definecolor{textcolor}{rgb}{0.000000,0.000000,0.000000}%
\pgfsetstrokecolor{textcolor}%
\pgfsetfillcolor{textcolor}%
\pgftext[x=2.105006in, y=0.835200in, left, base]{\color{textcolor}{\sffamily\fontsize{8.000000}{9.600000}\selectfont\catcode`\^=\active\def^{\ifmmode\sp\else\^{}\fi}\catcode`\%=\active\def%{\%}0.2}}%
\end{pgfscope}%
\begin{pgfscope}%
\pgfsetbuttcap%
\pgfsetroundjoin%
\definecolor{currentfill}{rgb}{0.000000,0.000000,0.000000}%
\pgfsetfillcolor{currentfill}%
\pgfsetlinewidth{0.501875pt}%
\definecolor{currentstroke}{rgb}{0.000000,0.000000,0.000000}%
\pgfsetstrokecolor{currentstroke}%
\pgfsetdash{}{0pt}%
\pgfsys@defobject{currentmarker}{\pgfqpoint{0.000000in}{0.000000in}}{\pgfqpoint{0.041667in}{0.000000in}}{%
\pgfpathmoveto{\pgfqpoint{0.000000in}{0.000000in}}%
\pgfpathlineto{\pgfqpoint{0.041667in}{0.000000in}}%
\pgfusepath{stroke,fill}%
}%
\begin{pgfscope}%
\pgfsys@transformshift{2.330321in}{1.139780in}%
\pgfsys@useobject{currentmarker}{}%
\end{pgfscope}%
\end{pgfscope}%
\begin{pgfscope}%
\definecolor{textcolor}{rgb}{0.000000,0.000000,0.000000}%
\pgfsetstrokecolor{textcolor}%
\pgfsetfillcolor{textcolor}%
\pgftext[x=2.105006in, y=1.097571in, left, base]{\color{textcolor}{\sffamily\fontsize{8.000000}{9.600000}\selectfont\catcode`\^=\active\def^{\ifmmode\sp\else\^{}\fi}\catcode`\%=\active\def%{\%}0.4}}%
\end{pgfscope}%
\begin{pgfscope}%
\pgfsetbuttcap%
\pgfsetroundjoin%
\definecolor{currentfill}{rgb}{0.000000,0.000000,0.000000}%
\pgfsetfillcolor{currentfill}%
\pgfsetlinewidth{0.501875pt}%
\definecolor{currentstroke}{rgb}{0.000000,0.000000,0.000000}%
\pgfsetstrokecolor{currentstroke}%
\pgfsetdash{}{0pt}%
\pgfsys@defobject{currentmarker}{\pgfqpoint{0.000000in}{0.000000in}}{\pgfqpoint{0.041667in}{0.000000in}}{%
\pgfpathmoveto{\pgfqpoint{0.000000in}{0.000000in}}%
\pgfpathlineto{\pgfqpoint{0.041667in}{0.000000in}}%
\pgfusepath{stroke,fill}%
}%
\begin{pgfscope}%
\pgfsys@transformshift{2.330321in}{1.402150in}%
\pgfsys@useobject{currentmarker}{}%
\end{pgfscope}%
\end{pgfscope}%
\begin{pgfscope}%
\definecolor{textcolor}{rgb}{0.000000,0.000000,0.000000}%
\pgfsetstrokecolor{textcolor}%
\pgfsetfillcolor{textcolor}%
\pgftext[x=2.105006in, y=1.359941in, left, base]{\color{textcolor}{\sffamily\fontsize{8.000000}{9.600000}\selectfont\catcode`\^=\active\def^{\ifmmode\sp\else\^{}\fi}\catcode`\%=\active\def%{\%}0.6}}%
\end{pgfscope}%
\begin{pgfscope}%
\pgfsetbuttcap%
\pgfsetroundjoin%
\definecolor{currentfill}{rgb}{0.000000,0.000000,0.000000}%
\pgfsetfillcolor{currentfill}%
\pgfsetlinewidth{0.501875pt}%
\definecolor{currentstroke}{rgb}{0.000000,0.000000,0.000000}%
\pgfsetstrokecolor{currentstroke}%
\pgfsetdash{}{0pt}%
\pgfsys@defobject{currentmarker}{\pgfqpoint{0.000000in}{0.000000in}}{\pgfqpoint{0.041667in}{0.000000in}}{%
\pgfpathmoveto{\pgfqpoint{0.000000in}{0.000000in}}%
\pgfpathlineto{\pgfqpoint{0.041667in}{0.000000in}}%
\pgfusepath{stroke,fill}%
}%
\begin{pgfscope}%
\pgfsys@transformshift{2.330321in}{1.664521in}%
\pgfsys@useobject{currentmarker}{}%
\end{pgfscope}%
\end{pgfscope}%
\begin{pgfscope}%
\definecolor{textcolor}{rgb}{0.000000,0.000000,0.000000}%
\pgfsetstrokecolor{textcolor}%
\pgfsetfillcolor{textcolor}%
\pgftext[x=2.105006in, y=1.622311in, left, base]{\color{textcolor}{\sffamily\fontsize{8.000000}{9.600000}\selectfont\catcode`\^=\active\def^{\ifmmode\sp\else\^{}\fi}\catcode`\%=\active\def%{\%}0.8}}%
\end{pgfscope}%
\begin{pgfscope}%
\pgfsetbuttcap%
\pgfsetroundjoin%
\definecolor{currentfill}{rgb}{0.000000,0.000000,0.000000}%
\pgfsetfillcolor{currentfill}%
\pgfsetlinewidth{0.501875pt}%
\definecolor{currentstroke}{rgb}{0.000000,0.000000,0.000000}%
\pgfsetstrokecolor{currentstroke}%
\pgfsetdash{}{0pt}%
\pgfsys@defobject{currentmarker}{\pgfqpoint{0.000000in}{0.000000in}}{\pgfqpoint{0.041667in}{0.000000in}}{%
\pgfpathmoveto{\pgfqpoint{0.000000in}{0.000000in}}%
\pgfpathlineto{\pgfqpoint{0.041667in}{0.000000in}}%
\pgfusepath{stroke,fill}%
}%
\begin{pgfscope}%
\pgfsys@transformshift{2.330321in}{1.926891in}%
\pgfsys@useobject{currentmarker}{}%
\end{pgfscope}%
\end{pgfscope}%
\begin{pgfscope}%
\definecolor{textcolor}{rgb}{0.000000,0.000000,0.000000}%
\pgfsetstrokecolor{textcolor}%
\pgfsetfillcolor{textcolor}%
\pgftext[x=2.105006in, y=1.884682in, left, base]{\color{textcolor}{\sffamily\fontsize{8.000000}{9.600000}\selectfont\catcode`\^=\active\def^{\ifmmode\sp\else\^{}\fi}\catcode`\%=\active\def%{\%}1.0}}%
\end{pgfscope}%
\begin{pgfscope}%
\pgfsetbuttcap%
\pgfsetroundjoin%
\definecolor{currentfill}{rgb}{0.000000,0.000000,0.000000}%
\pgfsetfillcolor{currentfill}%
\pgfsetlinewidth{0.301125pt}%
\definecolor{currentstroke}{rgb}{0.000000,0.000000,0.000000}%
\pgfsetstrokecolor{currentstroke}%
\pgfsetdash{}{0pt}%
\pgfsys@defobject{currentmarker}{\pgfqpoint{0.000000in}{0.000000in}}{\pgfqpoint{0.027778in}{0.000000in}}{%
\pgfpathmoveto{\pgfqpoint{0.000000in}{0.000000in}}%
\pgfpathlineto{\pgfqpoint{0.027778in}{0.000000in}}%
\pgfusepath{stroke,fill}%
}%
\begin{pgfscope}%
\pgfsys@transformshift{2.330321in}{0.418261in}%
\pgfsys@useobject{currentmarker}{}%
\end{pgfscope}%
\end{pgfscope}%
\begin{pgfscope}%
\pgfsetbuttcap%
\pgfsetroundjoin%
\definecolor{currentfill}{rgb}{0.000000,0.000000,0.000000}%
\pgfsetfillcolor{currentfill}%
\pgfsetlinewidth{0.301125pt}%
\definecolor{currentstroke}{rgb}{0.000000,0.000000,0.000000}%
\pgfsetstrokecolor{currentstroke}%
\pgfsetdash{}{0pt}%
\pgfsys@defobject{currentmarker}{\pgfqpoint{0.000000in}{0.000000in}}{\pgfqpoint{0.027778in}{0.000000in}}{%
\pgfpathmoveto{\pgfqpoint{0.000000in}{0.000000in}}%
\pgfpathlineto{\pgfqpoint{0.027778in}{0.000000in}}%
\pgfusepath{stroke,fill}%
}%
\begin{pgfscope}%
\pgfsys@transformshift{2.330321in}{0.483854in}%
\pgfsys@useobject{currentmarker}{}%
\end{pgfscope}%
\end{pgfscope}%
\begin{pgfscope}%
\pgfsetbuttcap%
\pgfsetroundjoin%
\definecolor{currentfill}{rgb}{0.000000,0.000000,0.000000}%
\pgfsetfillcolor{currentfill}%
\pgfsetlinewidth{0.301125pt}%
\definecolor{currentstroke}{rgb}{0.000000,0.000000,0.000000}%
\pgfsetstrokecolor{currentstroke}%
\pgfsetdash{}{0pt}%
\pgfsys@defobject{currentmarker}{\pgfqpoint{0.000000in}{0.000000in}}{\pgfqpoint{0.027778in}{0.000000in}}{%
\pgfpathmoveto{\pgfqpoint{0.000000in}{0.000000in}}%
\pgfpathlineto{\pgfqpoint{0.027778in}{0.000000in}}%
\pgfusepath{stroke,fill}%
}%
\begin{pgfscope}%
\pgfsys@transformshift{2.330321in}{0.549447in}%
\pgfsys@useobject{currentmarker}{}%
\end{pgfscope}%
\end{pgfscope}%
\begin{pgfscope}%
\pgfsetbuttcap%
\pgfsetroundjoin%
\definecolor{currentfill}{rgb}{0.000000,0.000000,0.000000}%
\pgfsetfillcolor{currentfill}%
\pgfsetlinewidth{0.301125pt}%
\definecolor{currentstroke}{rgb}{0.000000,0.000000,0.000000}%
\pgfsetstrokecolor{currentstroke}%
\pgfsetdash{}{0pt}%
\pgfsys@defobject{currentmarker}{\pgfqpoint{0.000000in}{0.000000in}}{\pgfqpoint{0.027778in}{0.000000in}}{%
\pgfpathmoveto{\pgfqpoint{0.000000in}{0.000000in}}%
\pgfpathlineto{\pgfqpoint{0.027778in}{0.000000in}}%
\pgfusepath{stroke,fill}%
}%
\begin{pgfscope}%
\pgfsys@transformshift{2.330321in}{0.680632in}%
\pgfsys@useobject{currentmarker}{}%
\end{pgfscope}%
\end{pgfscope}%
\begin{pgfscope}%
\pgfsetbuttcap%
\pgfsetroundjoin%
\definecolor{currentfill}{rgb}{0.000000,0.000000,0.000000}%
\pgfsetfillcolor{currentfill}%
\pgfsetlinewidth{0.301125pt}%
\definecolor{currentstroke}{rgb}{0.000000,0.000000,0.000000}%
\pgfsetstrokecolor{currentstroke}%
\pgfsetdash{}{0pt}%
\pgfsys@defobject{currentmarker}{\pgfqpoint{0.000000in}{0.000000in}}{\pgfqpoint{0.027778in}{0.000000in}}{%
\pgfpathmoveto{\pgfqpoint{0.000000in}{0.000000in}}%
\pgfpathlineto{\pgfqpoint{0.027778in}{0.000000in}}%
\pgfusepath{stroke,fill}%
}%
\begin{pgfscope}%
\pgfsys@transformshift{2.330321in}{0.746224in}%
\pgfsys@useobject{currentmarker}{}%
\end{pgfscope}%
\end{pgfscope}%
\begin{pgfscope}%
\pgfsetbuttcap%
\pgfsetroundjoin%
\definecolor{currentfill}{rgb}{0.000000,0.000000,0.000000}%
\pgfsetfillcolor{currentfill}%
\pgfsetlinewidth{0.301125pt}%
\definecolor{currentstroke}{rgb}{0.000000,0.000000,0.000000}%
\pgfsetstrokecolor{currentstroke}%
\pgfsetdash{}{0pt}%
\pgfsys@defobject{currentmarker}{\pgfqpoint{0.000000in}{0.000000in}}{\pgfqpoint{0.027778in}{0.000000in}}{%
\pgfpathmoveto{\pgfqpoint{0.000000in}{0.000000in}}%
\pgfpathlineto{\pgfqpoint{0.027778in}{0.000000in}}%
\pgfusepath{stroke,fill}%
}%
\begin{pgfscope}%
\pgfsys@transformshift{2.330321in}{0.811817in}%
\pgfsys@useobject{currentmarker}{}%
\end{pgfscope}%
\end{pgfscope}%
\begin{pgfscope}%
\pgfsetbuttcap%
\pgfsetroundjoin%
\definecolor{currentfill}{rgb}{0.000000,0.000000,0.000000}%
\pgfsetfillcolor{currentfill}%
\pgfsetlinewidth{0.301125pt}%
\definecolor{currentstroke}{rgb}{0.000000,0.000000,0.000000}%
\pgfsetstrokecolor{currentstroke}%
\pgfsetdash{}{0pt}%
\pgfsys@defobject{currentmarker}{\pgfqpoint{0.000000in}{0.000000in}}{\pgfqpoint{0.027778in}{0.000000in}}{%
\pgfpathmoveto{\pgfqpoint{0.000000in}{0.000000in}}%
\pgfpathlineto{\pgfqpoint{0.027778in}{0.000000in}}%
\pgfusepath{stroke,fill}%
}%
\begin{pgfscope}%
\pgfsys@transformshift{2.330321in}{0.943002in}%
\pgfsys@useobject{currentmarker}{}%
\end{pgfscope}%
\end{pgfscope}%
\begin{pgfscope}%
\pgfsetbuttcap%
\pgfsetroundjoin%
\definecolor{currentfill}{rgb}{0.000000,0.000000,0.000000}%
\pgfsetfillcolor{currentfill}%
\pgfsetlinewidth{0.301125pt}%
\definecolor{currentstroke}{rgb}{0.000000,0.000000,0.000000}%
\pgfsetstrokecolor{currentstroke}%
\pgfsetdash{}{0pt}%
\pgfsys@defobject{currentmarker}{\pgfqpoint{0.000000in}{0.000000in}}{\pgfqpoint{0.027778in}{0.000000in}}{%
\pgfpathmoveto{\pgfqpoint{0.000000in}{0.000000in}}%
\pgfpathlineto{\pgfqpoint{0.027778in}{0.000000in}}%
\pgfusepath{stroke,fill}%
}%
\begin{pgfscope}%
\pgfsys@transformshift{2.330321in}{1.008595in}%
\pgfsys@useobject{currentmarker}{}%
\end{pgfscope}%
\end{pgfscope}%
\begin{pgfscope}%
\pgfsetbuttcap%
\pgfsetroundjoin%
\definecolor{currentfill}{rgb}{0.000000,0.000000,0.000000}%
\pgfsetfillcolor{currentfill}%
\pgfsetlinewidth{0.301125pt}%
\definecolor{currentstroke}{rgb}{0.000000,0.000000,0.000000}%
\pgfsetstrokecolor{currentstroke}%
\pgfsetdash{}{0pt}%
\pgfsys@defobject{currentmarker}{\pgfqpoint{0.000000in}{0.000000in}}{\pgfqpoint{0.027778in}{0.000000in}}{%
\pgfpathmoveto{\pgfqpoint{0.000000in}{0.000000in}}%
\pgfpathlineto{\pgfqpoint{0.027778in}{0.000000in}}%
\pgfusepath{stroke,fill}%
}%
\begin{pgfscope}%
\pgfsys@transformshift{2.330321in}{1.074187in}%
\pgfsys@useobject{currentmarker}{}%
\end{pgfscope}%
\end{pgfscope}%
\begin{pgfscope}%
\pgfsetbuttcap%
\pgfsetroundjoin%
\definecolor{currentfill}{rgb}{0.000000,0.000000,0.000000}%
\pgfsetfillcolor{currentfill}%
\pgfsetlinewidth{0.301125pt}%
\definecolor{currentstroke}{rgb}{0.000000,0.000000,0.000000}%
\pgfsetstrokecolor{currentstroke}%
\pgfsetdash{}{0pt}%
\pgfsys@defobject{currentmarker}{\pgfqpoint{0.000000in}{0.000000in}}{\pgfqpoint{0.027778in}{0.000000in}}{%
\pgfpathmoveto{\pgfqpoint{0.000000in}{0.000000in}}%
\pgfpathlineto{\pgfqpoint{0.027778in}{0.000000in}}%
\pgfusepath{stroke,fill}%
}%
\begin{pgfscope}%
\pgfsys@transformshift{2.330321in}{1.205372in}%
\pgfsys@useobject{currentmarker}{}%
\end{pgfscope}%
\end{pgfscope}%
\begin{pgfscope}%
\pgfsetbuttcap%
\pgfsetroundjoin%
\definecolor{currentfill}{rgb}{0.000000,0.000000,0.000000}%
\pgfsetfillcolor{currentfill}%
\pgfsetlinewidth{0.301125pt}%
\definecolor{currentstroke}{rgb}{0.000000,0.000000,0.000000}%
\pgfsetstrokecolor{currentstroke}%
\pgfsetdash{}{0pt}%
\pgfsys@defobject{currentmarker}{\pgfqpoint{0.000000in}{0.000000in}}{\pgfqpoint{0.027778in}{0.000000in}}{%
\pgfpathmoveto{\pgfqpoint{0.000000in}{0.000000in}}%
\pgfpathlineto{\pgfqpoint{0.027778in}{0.000000in}}%
\pgfusepath{stroke,fill}%
}%
\begin{pgfscope}%
\pgfsys@transformshift{2.330321in}{1.270965in}%
\pgfsys@useobject{currentmarker}{}%
\end{pgfscope}%
\end{pgfscope}%
\begin{pgfscope}%
\pgfsetbuttcap%
\pgfsetroundjoin%
\definecolor{currentfill}{rgb}{0.000000,0.000000,0.000000}%
\pgfsetfillcolor{currentfill}%
\pgfsetlinewidth{0.301125pt}%
\definecolor{currentstroke}{rgb}{0.000000,0.000000,0.000000}%
\pgfsetstrokecolor{currentstroke}%
\pgfsetdash{}{0pt}%
\pgfsys@defobject{currentmarker}{\pgfqpoint{0.000000in}{0.000000in}}{\pgfqpoint{0.027778in}{0.000000in}}{%
\pgfpathmoveto{\pgfqpoint{0.000000in}{0.000000in}}%
\pgfpathlineto{\pgfqpoint{0.027778in}{0.000000in}}%
\pgfusepath{stroke,fill}%
}%
\begin{pgfscope}%
\pgfsys@transformshift{2.330321in}{1.336558in}%
\pgfsys@useobject{currentmarker}{}%
\end{pgfscope}%
\end{pgfscope}%
\begin{pgfscope}%
\pgfsetbuttcap%
\pgfsetroundjoin%
\definecolor{currentfill}{rgb}{0.000000,0.000000,0.000000}%
\pgfsetfillcolor{currentfill}%
\pgfsetlinewidth{0.301125pt}%
\definecolor{currentstroke}{rgb}{0.000000,0.000000,0.000000}%
\pgfsetstrokecolor{currentstroke}%
\pgfsetdash{}{0pt}%
\pgfsys@defobject{currentmarker}{\pgfqpoint{0.000000in}{0.000000in}}{\pgfqpoint{0.027778in}{0.000000in}}{%
\pgfpathmoveto{\pgfqpoint{0.000000in}{0.000000in}}%
\pgfpathlineto{\pgfqpoint{0.027778in}{0.000000in}}%
\pgfusepath{stroke,fill}%
}%
\begin{pgfscope}%
\pgfsys@transformshift{2.330321in}{1.467743in}%
\pgfsys@useobject{currentmarker}{}%
\end{pgfscope}%
\end{pgfscope}%
\begin{pgfscope}%
\pgfsetbuttcap%
\pgfsetroundjoin%
\definecolor{currentfill}{rgb}{0.000000,0.000000,0.000000}%
\pgfsetfillcolor{currentfill}%
\pgfsetlinewidth{0.301125pt}%
\definecolor{currentstroke}{rgb}{0.000000,0.000000,0.000000}%
\pgfsetstrokecolor{currentstroke}%
\pgfsetdash{}{0pt}%
\pgfsys@defobject{currentmarker}{\pgfqpoint{0.000000in}{0.000000in}}{\pgfqpoint{0.027778in}{0.000000in}}{%
\pgfpathmoveto{\pgfqpoint{0.000000in}{0.000000in}}%
\pgfpathlineto{\pgfqpoint{0.027778in}{0.000000in}}%
\pgfusepath{stroke,fill}%
}%
\begin{pgfscope}%
\pgfsys@transformshift{2.330321in}{1.533335in}%
\pgfsys@useobject{currentmarker}{}%
\end{pgfscope}%
\end{pgfscope}%
\begin{pgfscope}%
\pgfsetbuttcap%
\pgfsetroundjoin%
\definecolor{currentfill}{rgb}{0.000000,0.000000,0.000000}%
\pgfsetfillcolor{currentfill}%
\pgfsetlinewidth{0.301125pt}%
\definecolor{currentstroke}{rgb}{0.000000,0.000000,0.000000}%
\pgfsetstrokecolor{currentstroke}%
\pgfsetdash{}{0pt}%
\pgfsys@defobject{currentmarker}{\pgfqpoint{0.000000in}{0.000000in}}{\pgfqpoint{0.027778in}{0.000000in}}{%
\pgfpathmoveto{\pgfqpoint{0.000000in}{0.000000in}}%
\pgfpathlineto{\pgfqpoint{0.027778in}{0.000000in}}%
\pgfusepath{stroke,fill}%
}%
\begin{pgfscope}%
\pgfsys@transformshift{2.330321in}{1.598928in}%
\pgfsys@useobject{currentmarker}{}%
\end{pgfscope}%
\end{pgfscope}%
\begin{pgfscope}%
\pgfsetbuttcap%
\pgfsetroundjoin%
\definecolor{currentfill}{rgb}{0.000000,0.000000,0.000000}%
\pgfsetfillcolor{currentfill}%
\pgfsetlinewidth{0.301125pt}%
\definecolor{currentstroke}{rgb}{0.000000,0.000000,0.000000}%
\pgfsetstrokecolor{currentstroke}%
\pgfsetdash{}{0pt}%
\pgfsys@defobject{currentmarker}{\pgfqpoint{0.000000in}{0.000000in}}{\pgfqpoint{0.027778in}{0.000000in}}{%
\pgfpathmoveto{\pgfqpoint{0.000000in}{0.000000in}}%
\pgfpathlineto{\pgfqpoint{0.027778in}{0.000000in}}%
\pgfusepath{stroke,fill}%
}%
\begin{pgfscope}%
\pgfsys@transformshift{2.330321in}{1.730113in}%
\pgfsys@useobject{currentmarker}{}%
\end{pgfscope}%
\end{pgfscope}%
\begin{pgfscope}%
\pgfsetbuttcap%
\pgfsetroundjoin%
\definecolor{currentfill}{rgb}{0.000000,0.000000,0.000000}%
\pgfsetfillcolor{currentfill}%
\pgfsetlinewidth{0.301125pt}%
\definecolor{currentstroke}{rgb}{0.000000,0.000000,0.000000}%
\pgfsetstrokecolor{currentstroke}%
\pgfsetdash{}{0pt}%
\pgfsys@defobject{currentmarker}{\pgfqpoint{0.000000in}{0.000000in}}{\pgfqpoint{0.027778in}{0.000000in}}{%
\pgfpathmoveto{\pgfqpoint{0.000000in}{0.000000in}}%
\pgfpathlineto{\pgfqpoint{0.027778in}{0.000000in}}%
\pgfusepath{stroke,fill}%
}%
\begin{pgfscope}%
\pgfsys@transformshift{2.330321in}{1.795706in}%
\pgfsys@useobject{currentmarker}{}%
\end{pgfscope}%
\end{pgfscope}%
\begin{pgfscope}%
\pgfsetbuttcap%
\pgfsetroundjoin%
\definecolor{currentfill}{rgb}{0.000000,0.000000,0.000000}%
\pgfsetfillcolor{currentfill}%
\pgfsetlinewidth{0.301125pt}%
\definecolor{currentstroke}{rgb}{0.000000,0.000000,0.000000}%
\pgfsetstrokecolor{currentstroke}%
\pgfsetdash{}{0pt}%
\pgfsys@defobject{currentmarker}{\pgfqpoint{0.000000in}{0.000000in}}{\pgfqpoint{0.027778in}{0.000000in}}{%
\pgfpathmoveto{\pgfqpoint{0.000000in}{0.000000in}}%
\pgfpathlineto{\pgfqpoint{0.027778in}{0.000000in}}%
\pgfusepath{stroke,fill}%
}%
\begin{pgfscope}%
\pgfsys@transformshift{2.330321in}{1.861298in}%
\pgfsys@useobject{currentmarker}{}%
\end{pgfscope}%
\end{pgfscope}%
\begin{pgfscope}%
\pgfsetbuttcap%
\pgfsetroundjoin%
\definecolor{currentfill}{rgb}{0.000000,0.000000,0.000000}%
\pgfsetfillcolor{currentfill}%
\pgfsetlinewidth{0.301125pt}%
\definecolor{currentstroke}{rgb}{0.000000,0.000000,0.000000}%
\pgfsetstrokecolor{currentstroke}%
\pgfsetdash{}{0pt}%
\pgfsys@defobject{currentmarker}{\pgfqpoint{0.000000in}{0.000000in}}{\pgfqpoint{0.027778in}{0.000000in}}{%
\pgfpathmoveto{\pgfqpoint{0.000000in}{0.000000in}}%
\pgfpathlineto{\pgfqpoint{0.027778in}{0.000000in}}%
\pgfusepath{stroke,fill}%
}%
\begin{pgfscope}%
\pgfsys@transformshift{2.330321in}{1.992484in}%
\pgfsys@useobject{currentmarker}{}%
\end{pgfscope}%
\end{pgfscope}%
\begin{pgfscope}%
\pgfsetbuttcap%
\pgfsetroundjoin%
\definecolor{currentfill}{rgb}{0.000000,0.000000,0.000000}%
\pgfsetfillcolor{currentfill}%
\pgfsetlinewidth{0.301125pt}%
\definecolor{currentstroke}{rgb}{0.000000,0.000000,0.000000}%
\pgfsetstrokecolor{currentstroke}%
\pgfsetdash{}{0pt}%
\pgfsys@defobject{currentmarker}{\pgfqpoint{0.000000in}{0.000000in}}{\pgfqpoint{0.027778in}{0.000000in}}{%
\pgfpathmoveto{\pgfqpoint{0.000000in}{0.000000in}}%
\pgfpathlineto{\pgfqpoint{0.027778in}{0.000000in}}%
\pgfusepath{stroke,fill}%
}%
\begin{pgfscope}%
\pgfsys@transformshift{2.330321in}{2.058076in}%
\pgfsys@useobject{currentmarker}{}%
\end{pgfscope}%
\end{pgfscope}%
\begin{pgfscope}%
\pgfsetbuttcap%
\pgfsetroundjoin%
\definecolor{currentfill}{rgb}{0.000000,0.000000,0.000000}%
\pgfsetfillcolor{currentfill}%
\pgfsetlinewidth{0.301125pt}%
\definecolor{currentstroke}{rgb}{0.000000,0.000000,0.000000}%
\pgfsetstrokecolor{currentstroke}%
\pgfsetdash{}{0pt}%
\pgfsys@defobject{currentmarker}{\pgfqpoint{0.000000in}{0.000000in}}{\pgfqpoint{0.027778in}{0.000000in}}{%
\pgfpathmoveto{\pgfqpoint{0.000000in}{0.000000in}}%
\pgfpathlineto{\pgfqpoint{0.027778in}{0.000000in}}%
\pgfusepath{stroke,fill}%
}%
\begin{pgfscope}%
\pgfsys@transformshift{2.330321in}{2.123669in}%
\pgfsys@useobject{currentmarker}{}%
\end{pgfscope}%
\end{pgfscope}%
\begin{pgfscope}%
\pgfsetrectcap%
\pgfsetmiterjoin%
\pgfsetlinewidth{0.602250pt}%
\definecolor{currentstroke}{rgb}{0.000000,0.000000,0.000000}%
\pgfsetstrokecolor{currentstroke}%
\pgfsetdash{}{0pt}%
\pgfpathmoveto{\pgfqpoint{2.330321in}{0.352669in}}%
\pgfpathlineto{\pgfqpoint{2.330321in}{2.123669in}}%
\pgfusepath{stroke}%
\end{pgfscope}%
\begin{pgfscope}%
\pgfsetrectcap%
\pgfsetmiterjoin%
\pgfsetlinewidth{0.602250pt}%
\definecolor{currentstroke}{rgb}{0.000000,0.000000,0.000000}%
\pgfsetstrokecolor{currentstroke}%
\pgfsetdash{}{0pt}%
\pgfpathmoveto{\pgfqpoint{2.330321in}{0.352669in}}%
\pgfpathlineto{\pgfqpoint{3.680470in}{0.352669in}}%
\pgfusepath{stroke}%
\end{pgfscope}%
\begin{pgfscope}%
\pgfpathrectangle{\pgfqpoint{2.330321in}{0.352669in}}{\pgfqpoint{1.350149in}{1.771000in}}%
\pgfusepath{clip}%
\pgfsetrectcap%
\pgfsetroundjoin%
\pgfsetlinewidth{0.903375pt}%
\definecolor{currentstroke}{rgb}{0.250980,0.400000,0.600000}%
\pgfsetstrokecolor{currentstroke}%
\pgfsetdash{}{0pt}%
\pgfpathmoveto{\pgfqpoint{2.330321in}{0.665495in}}%
\pgfpathlineto{\pgfqpoint{2.342484in}{0.662158in}}%
\pgfpathlineto{\pgfqpoint{2.355999in}{0.660078in}}%
\pgfpathlineto{\pgfqpoint{2.369514in}{0.659616in}}%
\pgfpathlineto{\pgfqpoint{2.383029in}{0.660672in}}%
\pgfpathlineto{\pgfqpoint{2.396544in}{0.663151in}}%
\pgfpathlineto{\pgfqpoint{2.411411in}{0.667409in}}%
\pgfpathlineto{\pgfqpoint{2.426277in}{0.673153in}}%
\pgfpathlineto{\pgfqpoint{2.442495in}{0.680971in}}%
\pgfpathlineto{\pgfqpoint{2.460065in}{0.691101in}}%
\pgfpathlineto{\pgfqpoint{2.478986in}{0.703732in}}%
\pgfpathlineto{\pgfqpoint{2.500610in}{0.720072in}}%
\pgfpathlineto{\pgfqpoint{2.523585in}{0.739314in}}%
\pgfpathlineto{\pgfqpoint{2.550615in}{0.763940in}}%
\pgfpathlineto{\pgfqpoint{2.583051in}{0.795565in}}%
\pgfpathlineto{\pgfqpoint{2.631705in}{0.845334in}}%
\pgfpathlineto{\pgfqpoint{2.703335in}{0.918423in}}%
\pgfpathlineto{\pgfqpoint{2.739825in}{0.953395in}}%
\pgfpathlineto{\pgfqpoint{2.769558in}{0.979938in}}%
\pgfpathlineto{\pgfqpoint{2.796588in}{1.002171in}}%
\pgfpathlineto{\pgfqpoint{2.822267in}{1.021363in}}%
\pgfpathlineto{\pgfqpoint{2.845242in}{1.036768in}}%
\pgfpathlineto{\pgfqpoint{2.868218in}{1.050371in}}%
\pgfpathlineto{\pgfqpoint{2.889842in}{1.061427in}}%
\pgfpathlineto{\pgfqpoint{2.910114in}{1.070181in}}%
\pgfpathlineto{\pgfqpoint{2.930387in}{1.077321in}}%
\pgfpathlineto{\pgfqpoint{2.949308in}{1.082488in}}%
\pgfpathlineto{\pgfqpoint{2.968229in}{1.086184in}}%
\pgfpathlineto{\pgfqpoint{2.987150in}{1.088387in}}%
\pgfpathlineto{\pgfqpoint{3.006071in}{1.089087in}}%
\pgfpathlineto{\pgfqpoint{3.024992in}{1.088280in}}%
\pgfpathlineto{\pgfqpoint{3.043913in}{1.085969in}}%
\pgfpathlineto{\pgfqpoint{3.062834in}{1.082168in}}%
\pgfpathlineto{\pgfqpoint{3.081755in}{1.076896in}}%
\pgfpathlineto{\pgfqpoint{3.102027in}{1.069647in}}%
\pgfpathlineto{\pgfqpoint{3.122300in}{1.060787in}}%
\pgfpathlineto{\pgfqpoint{3.143924in}{1.049623in}}%
\pgfpathlineto{\pgfqpoint{3.165548in}{1.036768in}}%
\pgfpathlineto{\pgfqpoint{3.188523in}{1.021363in}}%
\pgfpathlineto{\pgfqpoint{3.212850in}{1.003230in}}%
\pgfpathlineto{\pgfqpoint{3.238529in}{0.982247in}}%
\pgfpathlineto{\pgfqpoint{3.266910in}{0.957129in}}%
\pgfpathlineto{\pgfqpoint{3.299346in}{0.926385in}}%
\pgfpathlineto{\pgfqpoint{3.339891in}{0.885766in}}%
\pgfpathlineto{\pgfqpoint{3.457472in}{0.766501in}}%
\pgfpathlineto{\pgfqpoint{3.487205in}{0.739314in}}%
\pgfpathlineto{\pgfqpoint{3.511532in}{0.718997in}}%
\pgfpathlineto{\pgfqpoint{3.533156in}{0.702775in}}%
\pgfpathlineto{\pgfqpoint{3.553429in}{0.689438in}}%
\pgfpathlineto{\pgfqpoint{3.570998in}{0.679562in}}%
\pgfpathlineto{\pgfqpoint{3.587216in}{0.672003in}}%
\pgfpathlineto{\pgfqpoint{3.603434in}{0.666096in}}%
\pgfpathlineto{\pgfqpoint{3.618301in}{0.662264in}}%
\pgfpathlineto{\pgfqpoint{3.631816in}{0.660201in}}%
\pgfpathlineto{\pgfqpoint{3.645331in}{0.659590in}}%
\pgfpathlineto{\pgfqpoint{3.658846in}{0.660528in}}%
\pgfpathlineto{\pgfqpoint{3.671009in}{0.662776in}}%
\pgfpathlineto{\pgfqpoint{3.680470in}{0.665495in}}%
\pgfpathlineto{\pgfqpoint{3.680470in}{0.665495in}}%
\pgfusepath{stroke}%
\end{pgfscope}%
\begin{pgfscope}%
\pgfpathrectangle{\pgfqpoint{2.330321in}{0.352669in}}{\pgfqpoint{1.350149in}{1.771000in}}%
\pgfusepath{clip}%
\pgfsetrectcap%
\pgfsetroundjoin%
\pgfsetlinewidth{0.903375pt}%
\definecolor{currentstroke}{rgb}{0.768627,0.556863,0.211765}%
\pgfsetstrokecolor{currentstroke}%
\pgfsetdash{}{0pt}%
\pgfpathmoveto{\pgfqpoint{2.330321in}{0.665495in}}%
\pgfpathlineto{\pgfqpoint{2.334375in}{0.672852in}}%
\pgfpathlineto{\pgfqpoint{2.338430in}{0.677483in}}%
\pgfpathlineto{\pgfqpoint{2.342484in}{0.679857in}}%
\pgfpathlineto{\pgfqpoint{2.346539in}{0.680396in}}%
\pgfpathlineto{\pgfqpoint{2.350593in}{0.679477in}}%
\pgfpathlineto{\pgfqpoint{2.355999in}{0.676558in}}%
\pgfpathlineto{\pgfqpoint{2.364108in}{0.669914in}}%
\pgfpathlineto{\pgfqpoint{2.385732in}{0.650561in}}%
\pgfpathlineto{\pgfqpoint{2.393841in}{0.645787in}}%
\pgfpathlineto{\pgfqpoint{2.400599in}{0.643549in}}%
\pgfpathlineto{\pgfqpoint{2.407356in}{0.643054in}}%
\pgfpathlineto{\pgfqpoint{2.414114in}{0.644346in}}%
\pgfpathlineto{\pgfqpoint{2.420871in}{0.647391in}}%
\pgfpathlineto{\pgfqpoint{2.427629in}{0.652086in}}%
\pgfpathlineto{\pgfqpoint{2.435738in}{0.659682in}}%
\pgfpathlineto{\pgfqpoint{2.445198in}{0.670820in}}%
\pgfpathlineto{\pgfqpoint{2.457362in}{0.687780in}}%
\pgfpathlineto{\pgfqpoint{2.504664in}{0.757066in}}%
\pgfpathlineto{\pgfqpoint{2.515476in}{0.768751in}}%
\pgfpathlineto{\pgfqpoint{2.524937in}{0.776690in}}%
\pgfpathlineto{\pgfqpoint{2.533046in}{0.781633in}}%
\pgfpathlineto{\pgfqpoint{2.541155in}{0.784784in}}%
\pgfpathlineto{\pgfqpoint{2.549264in}{0.786130in}}%
\pgfpathlineto{\pgfqpoint{2.557373in}{0.785709in}}%
\pgfpathlineto{\pgfqpoint{2.565482in}{0.783607in}}%
\pgfpathlineto{\pgfqpoint{2.573591in}{0.779954in}}%
\pgfpathlineto{\pgfqpoint{2.583051in}{0.773962in}}%
\pgfpathlineto{\pgfqpoint{2.593863in}{0.765239in}}%
\pgfpathlineto{\pgfqpoint{2.608730in}{0.751027in}}%
\pgfpathlineto{\pgfqpoint{2.641166in}{0.719103in}}%
\pgfpathlineto{\pgfqpoint{2.651978in}{0.710776in}}%
\pgfpathlineto{\pgfqpoint{2.661438in}{0.705396in}}%
\pgfpathlineto{\pgfqpoint{2.669547in}{0.702560in}}%
\pgfpathlineto{\pgfqpoint{2.676305in}{0.701652in}}%
\pgfpathlineto{\pgfqpoint{2.683062in}{0.702220in}}%
\pgfpathlineto{\pgfqpoint{2.689820in}{0.704395in}}%
\pgfpathlineto{\pgfqpoint{2.696577in}{0.708295in}}%
\pgfpathlineto{\pgfqpoint{2.703335in}{0.714025in}}%
\pgfpathlineto{\pgfqpoint{2.710092in}{0.721676in}}%
\pgfpathlineto{\pgfqpoint{2.718201in}{0.733499in}}%
\pgfpathlineto{\pgfqpoint{2.726310in}{0.748296in}}%
\pgfpathlineto{\pgfqpoint{2.734419in}{0.766136in}}%
\pgfpathlineto{\pgfqpoint{2.743880in}{0.790835in}}%
\pgfpathlineto{\pgfqpoint{2.753340in}{0.819710in}}%
\pgfpathlineto{\pgfqpoint{2.764152in}{0.857727in}}%
\pgfpathlineto{\pgfqpoint{2.776316in}{0.906637in}}%
\pgfpathlineto{\pgfqpoint{2.789831in}{0.968066in}}%
\pgfpathlineto{\pgfqpoint{2.804697in}{1.043265in}}%
\pgfpathlineto{\pgfqpoint{2.822267in}{1.140529in}}%
\pgfpathlineto{\pgfqpoint{2.845242in}{1.276914in}}%
\pgfpathlineto{\pgfqpoint{2.893896in}{1.568400in}}%
\pgfpathlineto{\pgfqpoint{2.911466in}{1.663406in}}%
\pgfpathlineto{\pgfqpoint{2.926332in}{1.735505in}}%
\pgfpathlineto{\pgfqpoint{2.938496in}{1.787493in}}%
\pgfpathlineto{\pgfqpoint{2.949308in}{1.827645in}}%
\pgfpathlineto{\pgfqpoint{2.958768in}{1.857657in}}%
\pgfpathlineto{\pgfqpoint{2.968229in}{1.882565in}}%
\pgfpathlineto{\pgfqpoint{2.976338in}{1.899657in}}%
\pgfpathlineto{\pgfqpoint{2.983095in}{1.910797in}}%
\pgfpathlineto{\pgfqpoint{2.989853in}{1.919054in}}%
\pgfpathlineto{\pgfqpoint{2.995259in}{1.923553in}}%
\pgfpathlineto{\pgfqpoint{3.000665in}{1.926164in}}%
\pgfpathlineto{\pgfqpoint{3.006071in}{1.926876in}}%
\pgfpathlineto{\pgfqpoint{3.011477in}{1.925689in}}%
\pgfpathlineto{\pgfqpoint{3.016883in}{1.922605in}}%
\pgfpathlineto{\pgfqpoint{3.022289in}{1.917636in}}%
\pgfpathlineto{\pgfqpoint{3.027695in}{1.910797in}}%
\pgfpathlineto{\pgfqpoint{3.034452in}{1.899657in}}%
\pgfpathlineto{\pgfqpoint{3.041210in}{1.885691in}}%
\pgfpathlineto{\pgfqpoint{3.049319in}{1.865306in}}%
\pgfpathlineto{\pgfqpoint{3.058779in}{1.836724in}}%
\pgfpathlineto{\pgfqpoint{3.068240in}{1.803251in}}%
\pgfpathlineto{\pgfqpoint{3.079052in}{1.759452in}}%
\pgfpathlineto{\pgfqpoint{3.091215in}{1.703809in}}%
\pgfpathlineto{\pgfqpoint{3.106082in}{1.627967in}}%
\pgfpathlineto{\pgfqpoint{3.123651in}{1.529676in}}%
\pgfpathlineto{\pgfqpoint{3.147978in}{1.384124in}}%
\pgfpathlineto{\pgfqpoint{3.189875in}{1.132781in}}%
\pgfpathlineto{\pgfqpoint{3.208796in}{1.029055in}}%
\pgfpathlineto{\pgfqpoint{3.223662in}{0.955219in}}%
\pgfpathlineto{\pgfqpoint{3.237177in}{0.895226in}}%
\pgfpathlineto{\pgfqpoint{3.249341in}{0.847730in}}%
\pgfpathlineto{\pgfqpoint{3.260153in}{0.811037in}}%
\pgfpathlineto{\pgfqpoint{3.270965in}{0.779737in}}%
\pgfpathlineto{\pgfqpoint{3.280425in}{0.756833in}}%
\pgfpathlineto{\pgfqpoint{3.288534in}{0.740521in}}%
\pgfpathlineto{\pgfqpoint{3.296643in}{0.727221in}}%
\pgfpathlineto{\pgfqpoint{3.304752in}{0.716850in}}%
\pgfpathlineto{\pgfqpoint{3.311510in}{0.710362in}}%
\pgfpathlineto{\pgfqpoint{3.318267in}{0.705741in}}%
\pgfpathlineto{\pgfqpoint{3.325025in}{0.702890in}}%
\pgfpathlineto{\pgfqpoint{3.331782in}{0.701695in}}%
\pgfpathlineto{\pgfqpoint{3.338540in}{0.702029in}}%
\pgfpathlineto{\pgfqpoint{3.345297in}{0.703754in}}%
\pgfpathlineto{\pgfqpoint{3.353406in}{0.707452in}}%
\pgfpathlineto{\pgfqpoint{3.362867in}{0.713660in}}%
\pgfpathlineto{\pgfqpoint{3.373679in}{0.722688in}}%
\pgfpathlineto{\pgfqpoint{3.388545in}{0.737192in}}%
\pgfpathlineto{\pgfqpoint{3.416927in}{0.765239in}}%
\pgfpathlineto{\pgfqpoint{3.429091in}{0.774921in}}%
\pgfpathlineto{\pgfqpoint{3.438551in}{0.780664in}}%
\pgfpathlineto{\pgfqpoint{3.446660in}{0.784070in}}%
\pgfpathlineto{\pgfqpoint{3.454769in}{0.785899in}}%
\pgfpathlineto{\pgfqpoint{3.462878in}{0.786030in}}%
\pgfpathlineto{\pgfqpoint{3.469636in}{0.784784in}}%
\pgfpathlineto{\pgfqpoint{3.477745in}{0.781633in}}%
\pgfpathlineto{\pgfqpoint{3.485854in}{0.776690in}}%
\pgfpathlineto{\pgfqpoint{3.493963in}{0.770022in}}%
\pgfpathlineto{\pgfqpoint{3.503423in}{0.760227in}}%
\pgfpathlineto{\pgfqpoint{3.514235in}{0.746734in}}%
\pgfpathlineto{\pgfqpoint{3.527750in}{0.727296in}}%
\pgfpathlineto{\pgfqpoint{3.566944in}{0.669100in}}%
\pgfpathlineto{\pgfqpoint{3.576404in}{0.658280in}}%
\pgfpathlineto{\pgfqpoint{3.584513in}{0.651021in}}%
\pgfpathlineto{\pgfqpoint{3.591271in}{0.646646in}}%
\pgfpathlineto{\pgfqpoint{3.598028in}{0.643946in}}%
\pgfpathlineto{\pgfqpoint{3.604786in}{0.643010in}}%
\pgfpathlineto{\pgfqpoint{3.611543in}{0.643860in}}%
\pgfpathlineto{\pgfqpoint{3.618301in}{0.646433in}}%
\pgfpathlineto{\pgfqpoint{3.625058in}{0.650561in}}%
\pgfpathlineto{\pgfqpoint{3.634519in}{0.658379in}}%
\pgfpathlineto{\pgfqpoint{3.656143in}{0.677435in}}%
\pgfpathlineto{\pgfqpoint{3.661549in}{0.679924in}}%
\pgfpathlineto{\pgfqpoint{3.665603in}{0.680396in}}%
\pgfpathlineto{\pgfqpoint{3.669658in}{0.679290in}}%
\pgfpathlineto{\pgfqpoint{3.673712in}{0.676212in}}%
\pgfpathlineto{\pgfqpoint{3.677767in}{0.670727in}}%
\pgfpathlineto{\pgfqpoint{3.680470in}{0.665495in}}%
\pgfpathlineto{\pgfqpoint{3.680470in}{0.665495in}}%
\pgfusepath{stroke}%
\end{pgfscope}%
\begin{pgfscope}%
\pgfpathrectangle{\pgfqpoint{2.330321in}{0.352669in}}{\pgfqpoint{1.350149in}{1.771000in}}%
\pgfusepath{clip}%
\pgfsetrectcap%
\pgfsetroundjoin%
\pgfsetlinewidth{0.903375pt}%
\definecolor{currentstroke}{rgb}{0.619608,0.188235,0.172549}%
\pgfsetstrokecolor{currentstroke}%
\pgfsetdash{}{0pt}%
\pgfpathmoveto{\pgfqpoint{2.330321in}{0.665495in}}%
\pgfpathlineto{\pgfqpoint{2.337078in}{0.667094in}}%
\pgfpathlineto{\pgfqpoint{2.349242in}{0.667822in}}%
\pgfpathlineto{\pgfqpoint{2.364108in}{0.668995in}}%
\pgfpathlineto{\pgfqpoint{2.376272in}{0.671432in}}%
\pgfpathlineto{\pgfqpoint{2.393841in}{0.676707in}}%
\pgfpathlineto{\pgfqpoint{2.420871in}{0.684787in}}%
\pgfpathlineto{\pgfqpoint{2.437089in}{0.688006in}}%
\pgfpathlineto{\pgfqpoint{2.457362in}{0.690416in}}%
\pgfpathlineto{\pgfqpoint{2.487095in}{0.693918in}}%
\pgfpathlineto{\pgfqpoint{2.501961in}{0.697248in}}%
\pgfpathlineto{\pgfqpoint{2.515476in}{0.701774in}}%
\pgfpathlineto{\pgfqpoint{2.528991in}{0.707864in}}%
\pgfpathlineto{\pgfqpoint{2.543858in}{0.716263in}}%
\pgfpathlineto{\pgfqpoint{2.561427in}{0.727986in}}%
\pgfpathlineto{\pgfqpoint{2.629002in}{0.775085in}}%
\pgfpathlineto{\pgfqpoint{2.650626in}{0.787053in}}%
\pgfpathlineto{\pgfqpoint{2.687117in}{0.806863in}}%
\pgfpathlineto{\pgfqpoint{2.700632in}{0.816215in}}%
\pgfpathlineto{\pgfqpoint{2.711444in}{0.825405in}}%
\pgfpathlineto{\pgfqpoint{2.722256in}{0.836637in}}%
\pgfpathlineto{\pgfqpoint{2.731716in}{0.848511in}}%
\pgfpathlineto{\pgfqpoint{2.741177in}{0.862614in}}%
\pgfpathlineto{\pgfqpoint{2.750637in}{0.879220in}}%
\pgfpathlineto{\pgfqpoint{2.760098in}{0.898567in}}%
\pgfpathlineto{\pgfqpoint{2.770910in}{0.924279in}}%
\pgfpathlineto{\pgfqpoint{2.781722in}{0.954019in}}%
\pgfpathlineto{\pgfqpoint{2.792534in}{0.987881in}}%
\pgfpathlineto{\pgfqpoint{2.804697in}{1.030871in}}%
\pgfpathlineto{\pgfqpoint{2.818212in}{1.084450in}}%
\pgfpathlineto{\pgfqpoint{2.833079in}{1.149761in}}%
\pgfpathlineto{\pgfqpoint{2.850648in}{1.233944in}}%
\pgfpathlineto{\pgfqpoint{2.877678in}{1.372147in}}%
\pgfpathlineto{\pgfqpoint{2.908763in}{1.529522in}}%
\pgfpathlineto{\pgfqpoint{2.924981in}{1.604248in}}%
\pgfpathlineto{\pgfqpoint{2.938496in}{1.659787in}}%
\pgfpathlineto{\pgfqpoint{2.949308in}{1.698662in}}%
\pgfpathlineto{\pgfqpoint{2.958768in}{1.728003in}}%
\pgfpathlineto{\pgfqpoint{2.968229in}{1.752536in}}%
\pgfpathlineto{\pgfqpoint{2.976338in}{1.769464in}}%
\pgfpathlineto{\pgfqpoint{2.983095in}{1.780539in}}%
\pgfpathlineto{\pgfqpoint{2.989853in}{1.788768in}}%
\pgfpathlineto{\pgfqpoint{2.995259in}{1.793260in}}%
\pgfpathlineto{\pgfqpoint{3.000665in}{1.795868in}}%
\pgfpathlineto{\pgfqpoint{3.006071in}{1.796581in}}%
\pgfpathlineto{\pgfqpoint{3.011477in}{1.795394in}}%
\pgfpathlineto{\pgfqpoint{3.016883in}{1.792313in}}%
\pgfpathlineto{\pgfqpoint{3.022289in}{1.787354in}}%
\pgfpathlineto{\pgfqpoint{3.027695in}{1.780539in}}%
\pgfpathlineto{\pgfqpoint{3.034452in}{1.769464in}}%
\pgfpathlineto{\pgfqpoint{3.041210in}{1.755626in}}%
\pgfpathlineto{\pgfqpoint{3.049319in}{1.735520in}}%
\pgfpathlineto{\pgfqpoint{3.058779in}{1.707512in}}%
\pgfpathlineto{\pgfqpoint{3.069591in}{1.670004in}}%
\pgfpathlineto{\pgfqpoint{3.081755in}{1.621639in}}%
\pgfpathlineto{\pgfqpoint{3.095270in}{1.561562in}}%
\pgfpathlineto{\pgfqpoint{3.112839in}{1.476181in}}%
\pgfpathlineto{\pgfqpoint{3.147978in}{1.295357in}}%
\pgfpathlineto{\pgfqpoint{3.170954in}{1.181359in}}%
\pgfpathlineto{\pgfqpoint{3.188523in}{1.101640in}}%
\pgfpathlineto{\pgfqpoint{3.203390in}{1.041111in}}%
\pgfpathlineto{\pgfqpoint{3.216905in}{0.992404in}}%
\pgfpathlineto{\pgfqpoint{3.229068in}{0.954019in}}%
\pgfpathlineto{\pgfqpoint{3.239880in}{0.924279in}}%
\pgfpathlineto{\pgfqpoint{3.250692in}{0.898567in}}%
\pgfpathlineto{\pgfqpoint{3.261504in}{0.876685in}}%
\pgfpathlineto{\pgfqpoint{3.272316in}{0.858341in}}%
\pgfpathlineto{\pgfqpoint{3.283128in}{0.843166in}}%
\pgfpathlineto{\pgfqpoint{3.293940in}{0.830736in}}%
\pgfpathlineto{\pgfqpoint{3.304752in}{0.820585in}}%
\pgfpathlineto{\pgfqpoint{3.316916in}{0.811292in}}%
\pgfpathlineto{\pgfqpoint{3.331782in}{0.802043in}}%
\pgfpathlineto{\pgfqpoint{3.356109in}{0.789162in}}%
\pgfpathlineto{\pgfqpoint{3.381788in}{0.775085in}}%
\pgfpathlineto{\pgfqpoint{3.402060in}{0.762190in}}%
\pgfpathlineto{\pgfqpoint{3.426388in}{0.744794in}}%
\pgfpathlineto{\pgfqpoint{3.461527in}{0.719696in}}%
\pgfpathlineto{\pgfqpoint{3.477745in}{0.709988in}}%
\pgfpathlineto{\pgfqpoint{3.492611in}{0.702867in}}%
\pgfpathlineto{\pgfqpoint{3.506126in}{0.698031in}}%
\pgfpathlineto{\pgfqpoint{3.520993in}{0.694411in}}%
\pgfpathlineto{\pgfqpoint{3.538562in}{0.691869in}}%
\pgfpathlineto{\pgfqpoint{3.581810in}{0.686574in}}%
\pgfpathlineto{\pgfqpoint{3.598028in}{0.682646in}}%
\pgfpathlineto{\pgfqpoint{3.623707in}{0.674534in}}%
\pgfpathlineto{\pgfqpoint{3.639925in}{0.670182in}}%
\pgfpathlineto{\pgfqpoint{3.652088in}{0.668364in}}%
\pgfpathlineto{\pgfqpoint{3.680470in}{0.665495in}}%
\pgfpathlineto{\pgfqpoint{3.680470in}{0.665495in}}%
\pgfusepath{stroke}%
\end{pgfscope}%
\begin{pgfscope}%
\definecolor{textcolor}{rgb}{0.000000,0.000000,0.000000}%
\pgfsetstrokecolor{textcolor}%
\pgfsetfillcolor{textcolor}%
\pgftext[x=3.005395in,y=2.207002in,,base]{\color{textcolor}{\sffamily\fontsize{8.000000}{9.600000}\selectfont\catcode`\^=\active\def^{\ifmmode\sp\else\^{}\fi}\catcode`\%=\active\def%{\%}(b) Chebyshev nodes}}%
\end{pgfscope}%
\begin{pgfscope}%
\pgfpathrectangle{\pgfqpoint{2.330321in}{0.352669in}}{\pgfqpoint{1.350149in}{1.771000in}}%
\pgfusepath{clip}%
\pgfsetbuttcap%
\pgfsetroundjoin%
\definecolor{currentfill}{rgb}{0.250980,0.400000,0.600000}%
\pgfsetfillcolor{currentfill}%
\pgfsetlinewidth{1.003750pt}%
\definecolor{currentstroke}{rgb}{0.250980,0.400000,0.600000}%
\pgfsetstrokecolor{currentstroke}%
\pgfsetdash{}{0pt}%
\pgfsys@defobject{currentmarker}{\pgfqpoint{-0.013889in}{-0.013889in}}{\pgfqpoint{0.013889in}{0.013889in}}{%
\pgfpathmoveto{\pgfqpoint{0.000000in}{-0.013889in}}%
\pgfpathcurveto{\pgfqpoint{0.003683in}{-0.013889in}}{\pgfqpoint{0.007216in}{-0.012425in}}{\pgfqpoint{0.009821in}{-0.009821in}}%
\pgfpathcurveto{\pgfqpoint{0.012425in}{-0.007216in}}{\pgfqpoint{0.013889in}{-0.003683in}}{\pgfqpoint{0.013889in}{0.000000in}}%
\pgfpathcurveto{\pgfqpoint{0.013889in}{0.003683in}}{\pgfqpoint{0.012425in}{0.007216in}}{\pgfqpoint{0.009821in}{0.009821in}}%
\pgfpathcurveto{\pgfqpoint{0.007216in}{0.012425in}}{\pgfqpoint{0.003683in}{0.013889in}}{\pgfqpoint{0.000000in}{0.013889in}}%
\pgfpathcurveto{\pgfqpoint{-0.003683in}{0.013889in}}{\pgfqpoint{-0.007216in}{0.012425in}}{\pgfqpoint{-0.009821in}{0.009821in}}%
\pgfpathcurveto{\pgfqpoint{-0.012425in}{0.007216in}}{\pgfqpoint{-0.013889in}{0.003683in}}{\pgfqpoint{-0.013889in}{0.000000in}}%
\pgfpathcurveto{\pgfqpoint{-0.013889in}{-0.003683in}}{\pgfqpoint{-0.012425in}{-0.007216in}}{\pgfqpoint{-0.009821in}{-0.009821in}}%
\pgfpathcurveto{\pgfqpoint{-0.007216in}{-0.012425in}}{\pgfqpoint{-0.003683in}{-0.013889in}}{\pgfqpoint{0.000000in}{-0.013889in}}%
\pgfpathlineto{\pgfqpoint{0.000000in}{-0.013889in}}%
\pgfpathclose%
\pgfusepath{stroke,fill}%
}%
\begin{pgfscope}%
\pgfsys@transformshift{3.680470in}{0.665495in}%
\pgfsys@useobject{currentmarker}{}%
\end{pgfscope}%
\begin{pgfscope}%
\pgfsys@transformshift{3.551542in}{0.690595in}%
\pgfsys@useobject{currentmarker}{}%
\end{pgfscope}%
\begin{pgfscope}%
\pgfsys@transformshift{3.214005in}{1.002326in}%
\pgfsys@useobject{currentmarker}{}%
\end{pgfscope}%
\begin{pgfscope}%
\pgfsys@transformshift{2.796786in}{1.002326in}%
\pgfsys@useobject{currentmarker}{}%
\end{pgfscope}%
\begin{pgfscope}%
\pgfsys@transformshift{2.459248in}{0.690595in}%
\pgfsys@useobject{currentmarker}{}%
\end{pgfscope}%
\begin{pgfscope}%
\pgfsys@transformshift{2.330321in}{0.665495in}%
\pgfsys@useobject{currentmarker}{}%
\end{pgfscope}%
\end{pgfscope}%
\begin{pgfscope}%
\pgfpathrectangle{\pgfqpoint{2.330321in}{0.352669in}}{\pgfqpoint{1.350149in}{1.771000in}}%
\pgfusepath{clip}%
\pgfsetbuttcap%
\pgfsetroundjoin%
\definecolor{currentfill}{rgb}{0.768627,0.556863,0.211765}%
\pgfsetfillcolor{currentfill}%
\pgfsetlinewidth{1.003750pt}%
\definecolor{currentstroke}{rgb}{0.768627,0.556863,0.211765}%
\pgfsetstrokecolor{currentstroke}%
\pgfsetdash{}{0pt}%
\pgfsys@defobject{currentmarker}{\pgfqpoint{-0.013889in}{-0.013889in}}{\pgfqpoint{0.013889in}{0.013889in}}{%
\pgfpathmoveto{\pgfqpoint{0.000000in}{-0.013889in}}%
\pgfpathcurveto{\pgfqpoint{0.003683in}{-0.013889in}}{\pgfqpoint{0.007216in}{-0.012425in}}{\pgfqpoint{0.009821in}{-0.009821in}}%
\pgfpathcurveto{\pgfqpoint{0.012425in}{-0.007216in}}{\pgfqpoint{0.013889in}{-0.003683in}}{\pgfqpoint{0.013889in}{0.000000in}}%
\pgfpathcurveto{\pgfqpoint{0.013889in}{0.003683in}}{\pgfqpoint{0.012425in}{0.007216in}}{\pgfqpoint{0.009821in}{0.009821in}}%
\pgfpathcurveto{\pgfqpoint{0.007216in}{0.012425in}}{\pgfqpoint{0.003683in}{0.013889in}}{\pgfqpoint{0.000000in}{0.013889in}}%
\pgfpathcurveto{\pgfqpoint{-0.003683in}{0.013889in}}{\pgfqpoint{-0.007216in}{0.012425in}}{\pgfqpoint{-0.009821in}{0.009821in}}%
\pgfpathcurveto{\pgfqpoint{-0.012425in}{0.007216in}}{\pgfqpoint{-0.013889in}{0.003683in}}{\pgfqpoint{-0.013889in}{0.000000in}}%
\pgfpathcurveto{\pgfqpoint{-0.013889in}{-0.003683in}}{\pgfqpoint{-0.012425in}{-0.007216in}}{\pgfqpoint{-0.009821in}{-0.009821in}}%
\pgfpathcurveto{\pgfqpoint{-0.007216in}{-0.012425in}}{\pgfqpoint{-0.003683in}{-0.013889in}}{\pgfqpoint{0.000000in}{-0.013889in}}%
\pgfpathlineto{\pgfqpoint{0.000000in}{-0.013889in}}%
\pgfpathclose%
\pgfusepath{stroke,fill}%
}%
\begin{pgfscope}%
\pgfsys@transformshift{3.680470in}{0.665495in}%
\pgfsys@useobject{currentmarker}{}%
\end{pgfscope}%
\begin{pgfscope}%
\pgfsys@transformshift{3.647429in}{0.670596in}%
\pgfsys@useobject{currentmarker}{}%
\end{pgfscope}%
\begin{pgfscope}%
\pgfsys@transformshift{3.551542in}{0.690595in}%
\pgfsys@useobject{currentmarker}{}%
\end{pgfscope}%
\begin{pgfscope}%
\pgfsys@transformshift{3.402194in}{0.751162in}%
\pgfsys@useobject{currentmarker}{}%
\end{pgfscope}%
\begin{pgfscope}%
\pgfsys@transformshift{3.214005in}{1.002326in}%
\pgfsys@useobject{currentmarker}{}%
\end{pgfscope}%
\begin{pgfscope}%
\pgfsys@transformshift{3.005395in}{1.926891in}%
\pgfsys@useobject{currentmarker}{}%
\end{pgfscope}%
\begin{pgfscope}%
\pgfsys@transformshift{2.796786in}{1.002326in}%
\pgfsys@useobject{currentmarker}{}%
\end{pgfscope}%
\begin{pgfscope}%
\pgfsys@transformshift{2.608596in}{0.751162in}%
\pgfsys@useobject{currentmarker}{}%
\end{pgfscope}%
\begin{pgfscope}%
\pgfsys@transformshift{2.459248in}{0.690595in}%
\pgfsys@useobject{currentmarker}{}%
\end{pgfscope}%
\begin{pgfscope}%
\pgfsys@transformshift{2.363361in}{0.670596in}%
\pgfsys@useobject{currentmarker}{}%
\end{pgfscope}%
\begin{pgfscope}%
\pgfsys@transformshift{2.330321in}{0.665495in}%
\pgfsys@useobject{currentmarker}{}%
\end{pgfscope}%
\end{pgfscope}%
\begin{pgfscope}%
\pgfpathrectangle{\pgfqpoint{2.330321in}{0.352669in}}{\pgfqpoint{1.350149in}{1.771000in}}%
\pgfusepath{clip}%
\pgfsetbuttcap%
\pgfsetroundjoin%
\definecolor{currentfill}{rgb}{0.619608,0.188235,0.172549}%
\pgfsetfillcolor{currentfill}%
\pgfsetlinewidth{1.003750pt}%
\definecolor{currentstroke}{rgb}{0.619608,0.188235,0.172549}%
\pgfsetstrokecolor{currentstroke}%
\pgfsetdash{}{0pt}%
\pgfsys@defobject{currentmarker}{\pgfqpoint{-0.013889in}{-0.013889in}}{\pgfqpoint{0.013889in}{0.013889in}}{%
\pgfpathmoveto{\pgfqpoint{0.000000in}{-0.013889in}}%
\pgfpathcurveto{\pgfqpoint{0.003683in}{-0.013889in}}{\pgfqpoint{0.007216in}{-0.012425in}}{\pgfqpoint{0.009821in}{-0.009821in}}%
\pgfpathcurveto{\pgfqpoint{0.012425in}{-0.007216in}}{\pgfqpoint{0.013889in}{-0.003683in}}{\pgfqpoint{0.013889in}{0.000000in}}%
\pgfpathcurveto{\pgfqpoint{0.013889in}{0.003683in}}{\pgfqpoint{0.012425in}{0.007216in}}{\pgfqpoint{0.009821in}{0.009821in}}%
\pgfpathcurveto{\pgfqpoint{0.007216in}{0.012425in}}{\pgfqpoint{0.003683in}{0.013889in}}{\pgfqpoint{0.000000in}{0.013889in}}%
\pgfpathcurveto{\pgfqpoint{-0.003683in}{0.013889in}}{\pgfqpoint{-0.007216in}{0.012425in}}{\pgfqpoint{-0.009821in}{0.009821in}}%
\pgfpathcurveto{\pgfqpoint{-0.012425in}{0.007216in}}{\pgfqpoint{-0.013889in}{0.003683in}}{\pgfqpoint{-0.013889in}{0.000000in}}%
\pgfpathcurveto{\pgfqpoint{-0.013889in}{-0.003683in}}{\pgfqpoint{-0.012425in}{-0.007216in}}{\pgfqpoint{-0.009821in}{-0.009821in}}%
\pgfpathcurveto{\pgfqpoint{-0.007216in}{-0.012425in}}{\pgfqpoint{-0.003683in}{-0.013889in}}{\pgfqpoint{0.000000in}{-0.013889in}}%
\pgfpathlineto{\pgfqpoint{0.000000in}{-0.013889in}}%
\pgfpathclose%
\pgfusepath{stroke,fill}%
}%
\begin{pgfscope}%
\pgfsys@transformshift{3.680470in}{0.665495in}%
\pgfsys@useobject{currentmarker}{}%
\end{pgfscope}%
\begin{pgfscope}%
\pgfsys@transformshift{3.665718in}{0.667683in}%
\pgfsys@useobject{currentmarker}{}%
\end{pgfscope}%
\begin{pgfscope}%
\pgfsys@transformshift{3.622106in}{0.675039in}%
\pgfsys@useobject{currentmarker}{}%
\end{pgfscope}%
\begin{pgfscope}%
\pgfsys@transformshift{3.551542in}{0.690595in}%
\pgfsys@useobject{currentmarker}{}%
\end{pgfscope}%
\begin{pgfscope}%
\pgfsys@transformshift{3.457108in}{0.722626in}%
\pgfsys@useobject{currentmarker}{}%
\end{pgfscope}%
\begin{pgfscope}%
\pgfsys@transformshift{3.342932in}{0.795984in}%
\pgfsys@useobject{currentmarker}{}%
\end{pgfscope}%
\begin{pgfscope}%
\pgfsys@transformshift{3.214005in}{1.002326in}%
\pgfsys@useobject{currentmarker}{}%
\end{pgfscope}%
\begin{pgfscope}%
\pgfsys@transformshift{3.075960in}{1.645434in}%
\pgfsys@useobject{currentmarker}{}%
\end{pgfscope}%
\begin{pgfscope}%
\pgfsys@transformshift{2.934831in}{1.645434in}%
\pgfsys@useobject{currentmarker}{}%
\end{pgfscope}%
\begin{pgfscope}%
\pgfsys@transformshift{2.796786in}{1.002326in}%
\pgfsys@useobject{currentmarker}{}%
\end{pgfscope}%
\begin{pgfscope}%
\pgfsys@transformshift{2.667858in}{0.795984in}%
\pgfsys@useobject{currentmarker}{}%
\end{pgfscope}%
\begin{pgfscope}%
\pgfsys@transformshift{2.553682in}{0.722626in}%
\pgfsys@useobject{currentmarker}{}%
\end{pgfscope}%
\begin{pgfscope}%
\pgfsys@transformshift{2.459248in}{0.690595in}%
\pgfsys@useobject{currentmarker}{}%
\end{pgfscope}%
\begin{pgfscope}%
\pgfsys@transformshift{2.388684in}{0.675039in}%
\pgfsys@useobject{currentmarker}{}%
\end{pgfscope}%
\begin{pgfscope}%
\pgfsys@transformshift{2.345073in}{0.667683in}%
\pgfsys@useobject{currentmarker}{}%
\end{pgfscope}%
\begin{pgfscope}%
\pgfsys@transformshift{2.330321in}{0.665495in}%
\pgfsys@useobject{currentmarker}{}%
\end{pgfscope}%
\end{pgfscope}%
\begin{pgfscope}%
\pgfpathrectangle{\pgfqpoint{2.330321in}{0.352669in}}{\pgfqpoint{1.350149in}{1.771000in}}%
\pgfusepath{clip}%
\pgfsetrectcap%
\pgfsetroundjoin%
\pgfsetlinewidth{1.405250pt}%
\definecolor{currentstroke}{rgb}{0.000000,0.000000,0.000000}%
\pgfsetstrokecolor{currentstroke}%
\pgfsetdash{}{0pt}%
\pgfpathmoveto{\pgfqpoint{2.330321in}{0.665495in}}%
\pgfpathlineto{\pgfqpoint{2.377623in}{0.673036in}}%
\pgfpathlineto{\pgfqpoint{2.419520in}{0.681194in}}%
\pgfpathlineto{\pgfqpoint{2.456010in}{0.689757in}}%
\pgfpathlineto{\pgfqpoint{2.488446in}{0.698810in}}%
\pgfpathlineto{\pgfqpoint{2.518179in}{0.708595in}}%
\pgfpathlineto{\pgfqpoint{2.545209in}{0.719012in}}%
\pgfpathlineto{\pgfqpoint{2.569536in}{0.729898in}}%
\pgfpathlineto{\pgfqpoint{2.592512in}{0.741767in}}%
\pgfpathlineto{\pgfqpoint{2.614136in}{0.754630in}}%
\pgfpathlineto{\pgfqpoint{2.634408in}{0.768470in}}%
\pgfpathlineto{\pgfqpoint{2.653329in}{0.783234in}}%
\pgfpathlineto{\pgfqpoint{2.670899in}{0.798826in}}%
\pgfpathlineto{\pgfqpoint{2.687117in}{0.815104in}}%
\pgfpathlineto{\pgfqpoint{2.703335in}{0.833491in}}%
\pgfpathlineto{\pgfqpoint{2.718201in}{0.852493in}}%
\pgfpathlineto{\pgfqpoint{2.733068in}{0.873871in}}%
\pgfpathlineto{\pgfqpoint{2.747934in}{0.897992in}}%
\pgfpathlineto{\pgfqpoint{2.761449in}{0.922657in}}%
\pgfpathlineto{\pgfqpoint{2.774964in}{0.950307in}}%
\pgfpathlineto{\pgfqpoint{2.788479in}{0.981357in}}%
\pgfpathlineto{\pgfqpoint{2.801994in}{1.016271in}}%
\pgfpathlineto{\pgfqpoint{2.815509in}{1.055557in}}%
\pgfpathlineto{\pgfqpoint{2.829024in}{1.099755in}}%
\pgfpathlineto{\pgfqpoint{2.842539in}{1.149413in}}%
\pgfpathlineto{\pgfqpoint{2.856054in}{1.205041in}}%
\pgfpathlineto{\pgfqpoint{2.869569in}{1.267037in}}%
\pgfpathlineto{\pgfqpoint{2.884436in}{1.342783in}}%
\pgfpathlineto{\pgfqpoint{2.900654in}{1.434011in}}%
\pgfpathlineto{\pgfqpoint{2.920926in}{1.557863in}}%
\pgfpathlineto{\pgfqpoint{2.953362in}{1.757248in}}%
\pgfpathlineto{\pgfqpoint{2.965526in}{1.821673in}}%
\pgfpathlineto{\pgfqpoint{2.974986in}{1.863558in}}%
\pgfpathlineto{\pgfqpoint{2.981744in}{1.887834in}}%
\pgfpathlineto{\pgfqpoint{2.988501in}{1.906669in}}%
\pgfpathlineto{\pgfqpoint{2.993907in}{1.917462in}}%
\pgfpathlineto{\pgfqpoint{2.997962in}{1.922927in}}%
\pgfpathlineto{\pgfqpoint{3.002016in}{1.926070in}}%
\pgfpathlineto{\pgfqpoint{3.006071in}{1.926858in}}%
\pgfpathlineto{\pgfqpoint{3.008774in}{1.926070in}}%
\pgfpathlineto{\pgfqpoint{3.012828in}{1.922927in}}%
\pgfpathlineto{\pgfqpoint{3.016883in}{1.917462in}}%
\pgfpathlineto{\pgfqpoint{3.020937in}{1.909734in}}%
\pgfpathlineto{\pgfqpoint{3.026343in}{1.896053in}}%
\pgfpathlineto{\pgfqpoint{3.033101in}{1.873882in}}%
\pgfpathlineto{\pgfqpoint{3.041210in}{1.840650in}}%
\pgfpathlineto{\pgfqpoint{3.050670in}{1.794285in}}%
\pgfpathlineto{\pgfqpoint{3.062834in}{1.725849in}}%
\pgfpathlineto{\pgfqpoint{3.081755in}{1.608968in}}%
\pgfpathlineto{\pgfqpoint{3.110136in}{1.434011in}}%
\pgfpathlineto{\pgfqpoint{3.126354in}{1.342783in}}%
\pgfpathlineto{\pgfqpoint{3.141221in}{1.267037in}}%
\pgfpathlineto{\pgfqpoint{3.156087in}{1.199197in}}%
\pgfpathlineto{\pgfqpoint{3.169602in}{1.144186in}}%
\pgfpathlineto{\pgfqpoint{3.183117in}{1.095098in}}%
\pgfpathlineto{\pgfqpoint{3.196632in}{1.051416in}}%
\pgfpathlineto{\pgfqpoint{3.210147in}{1.012591in}}%
\pgfpathlineto{\pgfqpoint{3.223662in}{0.978086in}}%
\pgfpathlineto{\pgfqpoint{3.237177in}{0.947396in}}%
\pgfpathlineto{\pgfqpoint{3.250692in}{0.920063in}}%
\pgfpathlineto{\pgfqpoint{3.265559in}{0.893384in}}%
\pgfpathlineto{\pgfqpoint{3.280425in}{0.869791in}}%
\pgfpathlineto{\pgfqpoint{3.295292in}{0.848871in}}%
\pgfpathlineto{\pgfqpoint{3.311510in}{0.828678in}}%
\pgfpathlineto{\pgfqpoint{3.327728in}{0.810849in}}%
\pgfpathlineto{\pgfqpoint{3.345297in}{0.793817in}}%
\pgfpathlineto{\pgfqpoint{3.362867in}{0.778815in}}%
\pgfpathlineto{\pgfqpoint{3.381788in}{0.764593in}}%
\pgfpathlineto{\pgfqpoint{3.402060in}{0.751244in}}%
\pgfpathlineto{\pgfqpoint{3.423685in}{0.738820in}}%
\pgfpathlineto{\pgfqpoint{3.446660in}{0.727340in}}%
\pgfpathlineto{\pgfqpoint{3.470987in}{0.716798in}}%
\pgfpathlineto{\pgfqpoint{3.498017in}{0.706696in}}%
\pgfpathlineto{\pgfqpoint{3.527750in}{0.697193in}}%
\pgfpathlineto{\pgfqpoint{3.560186in}{0.688389in}}%
\pgfpathlineto{\pgfqpoint{3.596677in}{0.680050in}}%
\pgfpathlineto{\pgfqpoint{3.637222in}{0.672327in}}%
\pgfpathlineto{\pgfqpoint{3.680470in}{0.665495in}}%
\pgfpathlineto{\pgfqpoint{3.680470in}{0.665495in}}%
\pgfusepath{stroke}%
\end{pgfscope}%
\begin{pgfscope}%
\pgfsetrectcap%
\pgfsetroundjoin%
\pgfsetlinewidth{1.405250pt}%
\definecolor{currentstroke}{rgb}{0.000000,0.000000,0.000000}%
\pgfsetstrokecolor{currentstroke}%
\pgfsetdash{}{0pt}%
\pgfpathmoveto{\pgfqpoint{2.396987in}{2.016419in}}%
\pgfpathlineto{\pgfqpoint{2.480321in}{2.016419in}}%
\pgfpathlineto{\pgfqpoint{2.563654in}{2.016419in}}%
\pgfusepath{stroke}%
\end{pgfscope}%
\begin{pgfscope}%
\definecolor{textcolor}{rgb}{0.000000,0.000000,0.000000}%
\pgfsetstrokecolor{textcolor}%
\pgfsetfillcolor{textcolor}%
\pgftext[x=2.630321in,y=1.987252in,left,base]{\color{textcolor}{\sffamily\fontsize{6.000000}{7.200000}\selectfont\catcode`\^=\active\def^{\ifmmode\sp\else\^{}\fi}\catcode`\%=\active\def%{\%}$f(x)$}}%
\end{pgfscope}%
\begin{pgfscope}%
\pgfsetrectcap%
\pgfsetroundjoin%
\pgfsetlinewidth{0.903375pt}%
\definecolor{currentstroke}{rgb}{0.250980,0.400000,0.600000}%
\pgfsetstrokecolor{currentstroke}%
\pgfsetdash{}{0pt}%
\pgfpathmoveto{\pgfqpoint{2.396987in}{1.892022in}}%
\pgfpathlineto{\pgfqpoint{2.480321in}{1.892022in}}%
\pgfpathlineto{\pgfqpoint{2.563654in}{1.892022in}}%
\pgfusepath{stroke}%
\end{pgfscope}%
\begin{pgfscope}%
\definecolor{textcolor}{rgb}{0.000000,0.000000,0.000000}%
\pgfsetstrokecolor{textcolor}%
\pgfsetfillcolor{textcolor}%
\pgftext[x=2.630321in,y=1.862855in,left,base]{\color{textcolor}{\sffamily\fontsize{6.000000}{7.200000}\selectfont\catcode`\^=\active\def^{\ifmmode\sp\else\^{}\fi}\catcode`\%=\active\def%{\%}$N = 5$}}%
\end{pgfscope}%
\begin{pgfscope}%
\pgfsetrectcap%
\pgfsetroundjoin%
\pgfsetlinewidth{0.903375pt}%
\definecolor{currentstroke}{rgb}{0.768627,0.556863,0.211765}%
\pgfsetstrokecolor{currentstroke}%
\pgfsetdash{}{0pt}%
\pgfpathmoveto{\pgfqpoint{3.037533in}{2.016419in}}%
\pgfpathlineto{\pgfqpoint{3.120866in}{2.016419in}}%
\pgfpathlineto{\pgfqpoint{3.204199in}{2.016419in}}%
\pgfusepath{stroke}%
\end{pgfscope}%
\begin{pgfscope}%
\definecolor{textcolor}{rgb}{0.000000,0.000000,0.000000}%
\pgfsetstrokecolor{textcolor}%
\pgfsetfillcolor{textcolor}%
\pgftext[x=3.270866in,y=1.987252in,left,base]{\color{textcolor}{\sffamily\fontsize{6.000000}{7.200000}\selectfont\catcode`\^=\active\def^{\ifmmode\sp\else\^{}\fi}\catcode`\%=\active\def%{\%}$N = 10$}}%
\end{pgfscope}%
\begin{pgfscope}%
\pgfsetrectcap%
\pgfsetroundjoin%
\pgfsetlinewidth{0.903375pt}%
\definecolor{currentstroke}{rgb}{0.619608,0.188235,0.172549}%
\pgfsetstrokecolor{currentstroke}%
\pgfsetdash{}{0pt}%
\pgfpathmoveto{\pgfqpoint{3.037533in}{1.894104in}}%
\pgfpathlineto{\pgfqpoint{3.120866in}{1.894104in}}%
\pgfpathlineto{\pgfqpoint{3.204199in}{1.894104in}}%
\pgfusepath{stroke}%
\end{pgfscope}%
\begin{pgfscope}%
\definecolor{textcolor}{rgb}{0.000000,0.000000,0.000000}%
\pgfsetstrokecolor{textcolor}%
\pgfsetfillcolor{textcolor}%
\pgftext[x=3.270866in,y=1.864938in,left,base]{\color{textcolor}{\sffamily\fontsize{6.000000}{7.200000}\selectfont\catcode`\^=\active\def^{\ifmmode\sp\else\^{}\fi}\catcode`\%=\active\def%{\%}$N = 15$}}%
\end{pgfscope}%
\end{pgfpicture}%
\makeatother%
\endgroup%

\caption[Runge phenomenon on equidistant vs.\ Chebyshev nodes]{%
  Polynomial interpolation of $f(x) = 1/(1 + 25x^2)$ on $[-1,1]$.
  (a) Equidistant nodes: the interpolant converges near the centre
  but develops violent oscillations near the endpoints that grow
  with increasing polynomial degree --- the Runge phenomenon.
  (b) Chebyshev Gauss--Lobatto nodes (\cref{sec:spec-chebyshev}):
  the interpolant converges uniformly across the entire interval,
  with the error decreasing monotonically as $N$ increases.  The
  clustered node placement near the endpoints suppresses the node
  polynomial $\omega_{N+1}$, eliminating the instability.}
\label{fig:runge-phenomenon}
\end{figure}

\paragraph{The cure: clustered nodes.}

The Runge phenomenon is not a failure of polynomial interpolation
itself but of the node placement.  The fix is to cluster nodes toward
the endpoints of the interval, where the equidistant node polynomial
is largest.  The idea is to make $\abs{\omega_{N+1}(x)}$ as uniform as
possible across $[a, b]$, so that no region of the interval suffers
disproportionately large error.

The optimal node distribution on $[-1, 1]$ --- the one that minimises
$\max_{x} \abs{\omega_{N+1}(x)}$ --- turns out to be the zeros of the
$(N+1)$-th Chebyshev polynomial of the first kind:
%
\begin{equation}
  x_j = \cos\!\left(\frac{2j + 1}{2(N+1)}\,\pi\right) ,
  \qquad j = 0, 1, \ldots, N \,.
  \label{eq:chebyshev-zeros}
\end{equation}
%
These nodes cluster quadratically near $\pm 1$, with spacing
$\order{1/N^2}$ at the endpoints compared to $\order{1/N}$ in the
interior.  For this distribution, the node polynomial satisfies
$\max_{x \in [-1,1]} \abs{\omega_{N+1}(x)} = 2^{-N}$, the smallest
supremum norm on $[-1,1]$ achievable by any monic polynomial of
degree $N+1$ --- a classical result of Chebyshev \cite{Trefethen2000}.  The Lebesgue constant
grows only logarithmically, $\Lambda_N \sim (2/\pi)\ln N$, compared
to the exponential growth for equidistant nodes.

A closely related distribution --- the \emph{Chebyshev--Gauss--Lobatto}
(CGL) nodes --- includes the endpoints $\pm 1$:
%
\begin{equation}
  x_j = \cos\!\left(\frac{j}{N}\,\pi\right) ,
  \qquad j = 0, 1, \ldots, N \,.
  \label{eq:cgl-nodes}
\end{equation}
%
The CGL nodes are the extrema of the $N$-th Chebyshev polynomial
(which include the endpoints $\pm 1$).  They share the logarithmic Lebesgue
growth and the quadratic endpoint clustering of the Chebyshev zeros,
with the practical advantage that including $\pm 1$ simplifies
boundary conditions and inter-element coupling in a spectral-element
or DG context.
\xrefnote{The next section (\cref{sec:spec-chebyshev}) develops the
Chebyshev polynomial theory that explains why these particular nodes
are optimal.  \Cref{sec:spec-lgl} derives the Legendre--Gauss--Lobatto
(LGL) nodes used by the codebase's DG solver, which share the
essential clustering property.}

\paragraph{From interpolation to spectral methods.}

The Runge phenomenon teaches that polynomial approximation can
achieve arbitrarily high accuracy for smooth functions, but only if
the nodes are chosen correctly.  Spectral methods build on this
insight: they represent the numerical solution as a polynomial on
clustered nodes and use the polynomial structure to compute
derivatives, integrals, and inter-element fluxes.  The rest of this
chapter develops the machinery: Chebyshev and Legendre polynomial
bases that make clustered nodes natural, differentiation and
quadrature operators built directly from them, and the Vandermonde
matrices that convert between nodal and modal representations in the
DG solver.
